\documentclass[12pt, a4paper]{article}

\def\islight{true} 

\usepackage[utf8]{inputenc}
\usepackage[italian]{babel}
\usepackage[T1]{fontenc}
\usepackage{lmodern}
\usepackage{amsmath, amssymb}
\usepackage{geometry}
\usepackage{enumitem}
\usepackage{booktabs}
\usepackage{eurosym}
\usepackage{graphicx}
\usepackage{tcolorbox}
\usepackage{tocloft}
\usepackage{xcolor}
\usepackage{hyperref}
\usepackage{microtype}
\usepackage{datetime} 

\newcommand{\myfootertext}{}
\newcommand{\myicon}{}

\ifx\islight\undefined
    \definecolor{lightergray}{gray}{0.85}
    \definecolor{tocsec}{gray}{0.85}
    \definecolor{tocsubsec}{gray}{0.75}
    \renewcommand{\myicon}{logo_unitn_scuro.png}
    \renewcommand{\myfootertext}{Appunti redatti per gli studenti del DISI. \\ Eventuali errori possono essere segnalati su GitHub.}
    \pagecolor{black}
    \color{white}
\else
    \definecolor{lightergray}{gray}{0.35}
    \definecolor{tocsec}{gray}{0.35}
    \definecolor{tocsubsec}{gray}{0.45}
    \renewcommand{\myicon}{logo_unitn_chiaro.png}
    \renewcommand{\myfootertext}{Versione Light Mode. Disponibile anche la \href{https://github.com/mirconegri/University/tree/main/2nd_semester/statistica/schemi/secondo_parziale/schemi_probabilità_secondo_parziale_darkV.pdf}{versione Dark Mode}.}    \pagecolor{white}
    \color{black}
\fi

\hypersetup{
    colorlinks=true,
    linkcolor=lightergray,
    urlcolor=cyan,
    pdftitle={Appunti di Probabilità - Secondo Parziale},
    pdfauthor={Mirco Negri}
}

\renewcommand{\cftsecfont}{\color{tocsec}\bfseries}
\renewcommand{\cftsecpagefont}{\color{tocsec}\bfseries}
\renewcommand{\cftsubsecfont}{\color{tocsubsec}}
\renewcommand{\cftsubsecpagefont}{\color{tocsubsec}}
\renewcommand{\cftsubsecleader}{\color{tocsubsec}\cftdotfill{\cftdotsep}}

\geometry{margin=2.5cm}

\begin{document}

\begin{titlepage}
    \centering
    
    \includegraphics[width=0.25\textwidth]{\myicon}
    \vspace{1cm}
    
    {\scshape\LARGE Università degli Studi di Trento \par}
    \vspace{0.4cm}
    {\large Dipartimento di Ingegneria e Scienza dell'Informazione \par}
    
    \vspace{2cm}
    
    {\color{cyan}\rule{\textwidth}{1pt}}\vspace{0.5cm}
    {\Huge\bfseries Appunti di Probabilità e Statistica \par}
    \vspace{0.3cm}
    {\Large\itshape Secondo parziale \par}
    \vspace{0.4cm}{\color{cyan}\rule{\textwidth}{1pt}}
    
    \vspace{2.5cm}
    
    {\Large A cura di: \par}
    \vspace{0.3cm}
    {\Large \textbf{\href{https://github.com/mirconegri}{Mirco Negri}} \par}
    
    \vfill
    
    {\small \textit{\myfootertext} \par}
    \vspace{0.8cm}
    
    {\large \monthname\ \the\year \par}
\end{titlepage}

\newpage
\begingroup
\hypersetup{hidelinks}
\tableofcontents
\endgroup
\newpage

\section*{Probabilità e Statistica - Esercizi 5}
\addcontentsline{toc}{section}{Probabilità e Statistica - Esercizi 5}
\markboth{Esercizi 5}{Esercizi 5}

% ==========================================
%               ESERCIZIO 1
% ==========================================
\begin{tcolorbox}[colback=lightergray!10, colframe=tocsec, title=\textbf{Esercizio 1}]
Abbiamo una variabile aleatoria $X$ di distribuzione binomiale di parametri $n=17$ e $p=0.203$.
\end{tcolorbox}

\subsection*{Quesiti e soluzioni}

\subsubsection*{Quesito 1}
Qual è il valore atteso di $X$?

\textbf{Soluzione:}
Per definizione $\mathbb{E}(X)=\sum_{x=0}^{n} x p_X(x) = \dots = np$.
Essendo $X$ una variabile aleatoria con distribuzione binomiale, il valore atteso si espande e si calcola in modo diretto come il prodotto tra il numero di prove $n$ e la probabilità di successo $p$:
$$ \mathbb{E}(X) = 17 \cdot 0.203 = 3.451 $$
La risposta è: 3.451

\subsubsection*{Quesito 2}
Sia ora $Y$ una variabile aleatoria di Poisson con parametro $\lambda=np$ (ove $n,p$ sono quelli del quesito precedente). Qual è il valore atteso di $23Y+14.791$?

\textbf{Soluzione:}
Sappiamo che $\mathbb{E}(aY+b)=a\mathbb{E}(Y)+b$ per cui ci basta calcolare il valore atteso della v.a. di Poisson $Y$. Incolonniamo i passaggi della sommatoria:
$$ \begin{aligned} \mathbb{E}(Y) &= \sum_{y \in R_Y} y p(y) = \sum_{y \geq 0} y \frac{\lambda^y e^{-\lambda}}{y!} \\ &= \sum_{y=1}^\infty y \frac{\lambda^y e^{-\lambda}}{y!} = \sum_{y=1}^\infty y \frac{\lambda^y e^{-\lambda}}{y(y-1)!} \\ &= \lambda e^{-\lambda} \sum_{y=1}^\infty \frac{\lambda^{y-1}}{(y-1)!} = \lambda e^{-\lambda} \sum_{k=0}^\infty \frac{\lambda^k}{k!} \end{aligned} $$
Nell’ultimo passaggio abbiamo semplicemente sostituito $k=y-1$ così la serie parte nuovamente da 0. Ricordando che la definizione di $e^x=\exp(x)=\sum_{n=0}^\infty \frac{x^n}{n!}$, otteniamo che $\mathbb{E}(Y)=\lambda e^{-\lambda} e^\lambda = \lambda$.
Infine, applicando la proprietà di linearità dimostrata all'inizio: $\mathbb{E}(aY+b)=a\lambda+b$. 
Sostituendo i valori numerici:
$$ \mathbb{E}(23Y+14.791) = 23(3.451) + 14.791 = 79.373 + 14.791 = 94.164 $$
La risposta è: 94.164

\newpage

\subsubsection*{Quesito 3}
Sia la $Y$ la stessa del quesito precedente, qual è la probabilità $P(Y=0)$?

\textbf{Soluzione:}
Qui basta sostituire il valore desiderato nella funzione di densità discreta della v.a. Poissoniana $Y$, $P(Y=0)=p(0)=\frac{\lambda^0 e^{-\lambda}}{0!}$. Svolgendo i calcoli con $\lambda = 3.451$:
$$ P(Y=0) = \frac{3.451^0 \cdot e^{-3.451}}{1} = e^{-3.451} \approx 0.0317139 $$
La risposta è: 0.0317139


% ==========================================
%               ESERCIZIO 2
% ==========================================
\vspace{0.5cm}
\begin{tcolorbox}[colback=lightergray!10, colframe=tocsec, title=\textbf{Esercizio 2}]
Si consideri un dado non regolarmente equilibrato, in particolare la probabilità di ogni faccia del dado è
\begin{verbatim}
                     1    2    3    4    5    6 
                    0.07 0.20 0.25 0.06 0.06 0.36
\end{verbatim}

Se esce 6 o 5 per noi è un successo. Tutti gli altri numeri costituiscono insuccesso. Come avete visto a lezione, questo problema rientra nella classe di problemi in cui un esperimento elementare con due possibili risultati, detti successo e insuccesso.
L’evento elementare è un lancio del dado.
Consideriamo poi $N=18$ lanci del dado. Come avete visto si può modellare il problema con una variabile casuale discreta distribuita come una binomiale, $X \sim Bi(N,p)$
\end{tcolorbox}

\subsubsection*{Quesito 1}
1. Qual è la probabilità di successo $p$ dell’evento elementare?

\textbf{Soluzione:}
$p=P(\{6\} \cup \{5\}) = P(\{6\}) + P(\{5\}) = 0.36 + 0.06 = 0.42$.

Potete anche controllare che la probabilità di insuccesso è uguale al complementare dei due eventi che costituiscono successo, i.e. $1-p=$
La risposta è: 0.42

\subsubsection*{Quesito 2}
2. Qual è la probabilità di ottenere esattamente un successo (in $N=18$ lanci)?

\textbf{Soluzione:}
Utilizziamo la funzione di massa di probabilità della distribuzione Binomiale per un numero esatto di successi pari a $k=1$. La formula generale è $P(X=k) = \binom{N}{k} p^k (1-p)^{N-k}$. Sostituendo i parametri in nostro possesso:
$$ P(X=1) = \binom{18}{1} (0.42)^1 (1-0.42)^{18-1} = 18 \cdot 0.42 \cdot (0.58)^{17} \approx 7.191139 \times 10^{-4} $$
La risposta è: $7.191139 \times 10^{-4}$

\newpage

\subsubsection*{Quesito 3}
3. Qual è la probabilità di ottenere almeno 8 successi (in $N=18$ lanci)?

\textbf{Soluzione:}
Si osserva che almeno 8 successi, vuol dire 8 o più. Il numero massimo di successi ottenibili in 18 lanci è 18 da cui
$$ P(8 \leq X \leq 18) = \sum_{k=8}^{18} \binom{18}{k} p^k (1-p)^{18-k} $$
Sviluppando la somma per $p=0.42$ e accumulando le singole probabilità discrete:
$$ \sum_{k=8}^{18} \binom{18}{k} (0.42)^k (0.58)^{18-k} \approx 0.5061698 $$
La risposta è: 0.5061698


% ==========================================
%               ESERCIZIO 3
% ==========================================
\vspace{0.5cm}
\begin{tcolorbox}[colback=lightergray!10, colframe=tocsec, title=\textbf{Esercizio 3}]
Un ricercatore dopo mesi di progettazione ha fatto produrre il suo chip per la guida autonoma. Adesso lo deve testare. Dopo aver fatto moltissime simulazioni è giunto alla conclusione che il chip funziona correttamente nel 94.1\% dei casi.
\end{tcolorbox}

\subsubsection*{Quesito 1}
1. Qual è la probabilità che il chip sbagli la misurazione?

\textbf{Soluzione:}
Vediamo che la singola prova è rappresentata da una variabile $Y$ tale che:
$Y(\omega) = 0$ se il chip funziona correttamente
$Y(\omega) = 1$ se il chip non funziona correttamente.
Abbiamo che il chip non funziona correttamente nel 5.9\% dei casi. Da cui $p=0.059$.
In questo modo $Y \sim Bernoulli(p)$ con $p=0.059$.

La risposta è: 0.059

\subsubsection*{Quesito 2}
2. Indichiamo con $X \sim G(p)$ la v.a. geometrica che modella l’esperimento.
Per vedere se gli errori nelle misurazioni sono dovuti all’impiego continuo del chip, il ricercatore decide di tenere traccia di quando il chip fallisce.
Qual è la probabilità che le prime $m=27$ misurazioni siano esatte (i.e. la $m+1=28$ sia sbagliata)?

\textbf{Soluzione:}
E’ richiesta la probabilità che per osservare il primo successo – i.e. il chip non funziona correttamente – siano necessarie $m+1$ prove, i.e. 
$$ P(X=m+1) = p(1-p)^m $$
Trattandosi della funzione di densità di una variabile Geometrica, si descrive una sequenza di $m$ fallimenti seguiti da $1$ successo all'ultimo tentativo. Inserendo i parametri ricavati in precedenza:
$$ P(X=28) = 0.059 \cdot (1 - 0.059)^{27} = 0.059 \cdot (0.941)^{27} \approx 0.0114228 $$
La risposta è: 0.0114228

\subsubsection*{Quesito 3}
3. Le prime $m=27$ simulazioni si sono susseguite senza sbagli, qual è la probabilità che un errore si presenti per la prima volta dopo ulteriori $l=5$ simulazioni?

\textbf{Soluzione:}
Domanda sulla proprietà di assenza di memoria della variabile geometrica. Si veda l’Esempio 13 pag.88 note B.
$$ P(X=m+l \mid X>m) = P(X=l) = p(1-p)^{l-1} $$
L'informazione passata (il fatto di aver eseguito 27 prove senza errori) non modifica la probabilità degli eventi futuri. Valutiamo quindi la probabilità al quinto tentativo aggiuntivo:
$$ P(X=5) = 0.059 \cdot (0.941)^{5-1} = 0.059 \cdot (0.941)^4 \approx 0.0462605 $$
la risposta è: 0.0462605


% ==========================================
%               ESERCIZIO 4
% ==========================================
\vspace{0.5cm}
\begin{tcolorbox}[colback=lightergray!10, colframe=tocsec, title=\textbf{Esercizio 4}]
Jane Doe si appresta a fare un esame di italiano per potersi iscrivere all’Università di Trento. L’esame consiste in una successione di domande, tutte della stessa difficoltà. Appena Jane totalizza 16 risposte esatte, l’esame è superato. (Ovviamente il voto dipenderà dal tempo impiegato, ma questo non c’interessa.) Jane ha fatto molte simulazioni d’esame e sa che il suo tasso di risposte esatte è del 58.2\%.
\end{tcolorbox}

\subsubsection*{Quesito 1}
1. Si può descrivere questo esperimento con una variabile casuale distribuita come una Binomiale Negativa?
Rispondere TRUE o FALSE.

\textbf{Soluzione:}
\begin{itemize}
    \item prove ripetute: successione di domande;
    \item ogni prova ha solo due esiti possibili (che in genere indichiamo con successo e fallimento): ad ogni domanda Jane può dare una risposta giusta o sbagliata;
    \item la probabilità di successo $p$ è la stessa per ogni prova: ci viene detto che tutte le domande hanno la stessa difficoltà e la probabilità che Jane ha di dare la risposta giusta è 0.567 (dato del problema);
    \item le prove sono indipendenti: le domande lo sono;
    \item l’esperimento continua finché non si osservano $r$ successi, con $r$ fissato: ``Appena Jane totalizza 16 risposte esatte, l’esame è superato.''
\end{itemize}
La risposta è: TRUE

\subsubsection*{Quesito 2}
2. Quali sono i parametri $(r,p)$ della distribuzione?
La risposta deve avere la forma c(, ) (nell’ordine ordine richiesto). Se si è risposto FALSE alla domanda precedente inserire c(NA, NA).

\textbf{Soluzione:}
La risposta è: c(16, 0.567)

\subsubsection*{Quesito 3}
3. Indipendentemente dalle risposte date, si consideri la variabile aleatoria $X \sim BiNe(r,p)$ con $r=16$ e $p=0.567$, che rappresenta il numero di domande a cui Jane deve rispondere affinché l’esame termini.
Qual è la probabilità che Jane superi l’esame dopo 16 domande? Usare la distribuzione di $X$.

\textbf{Soluzione:}
Ricordiamo che l’esame è superato appena Jane totalizza 16 risposte esatte. Per cui la domanda può trovare risposta in molti modi (controllare che sono equivalenti);

Usando la funzione di probabilità di $X$:
$$ Pr(X=x) = \begin{cases} \binom{x-1}{r-1} p^r (1-p)^{x-r} & x=r, r+1, r+2, \dots \\ 0 & \text{altrove} \end{cases} $$
\begin{itemize}
    \item equivalentemente possiamo chiederci quale sia la probabilità che Jane non faccia alcun errore in 16 domande e possiamo utilizzare anche la distribuzione binomiale
    \item in R, controllate la documentazione della funzione dnbinom, poiché il risultato giusto si ottiene così: \texttt{dnbinom(x = (r - r), size = r, prob = p)}
\end{itemize}
Sostituendo direttamente $x=16$ per calcolare la probabilità di assenza di fallimenti:
$$ Pr(X=16) = \binom{15}{15} (0.567)^{16} (1-0.567)^0 = 1 \cdot (0.567)^{16} \cdot 1 \approx 1.1411138 \times 10^{-4} $$
La risposta è: $1.1411138 \times 10^{-4}$

\subsubsection*{Quesito 4}
4. Qual è la probabilità che Jane superi l’esame dopo 28 domande? Usare la distribuzione di $X$.

\textbf{Soluzione:}
Applichiamo la formula generale della Binomiale Negativa calcolando la probabilità che siano necessari esattamente $x=28$ tentativi totali per raccogliere l'ultimo dei 16 successi previsti:
$$ Pr(X=28) = \binom{28-1}{16-1} (0.567)^{16} (1-0.567)^{28-16} = \binom{27}{15} (0.567)^{16} (0.433)^{12} \approx 0.0861648 $$
La risposta è: 0.0861648

\subsubsection*{Quesito 5}
5. Calcolare $P(X \leq 20.28)$?

\textbf{Soluzione:}
Trattandosi di una variabile aleatoria discreta definita sui numeri naturali (conta il numero di domande), la richiesta $P(X \leq 20.28)$ è matematicamente equivalente a calcolare la funzione di ripartizione $P(X \leq 20)$. Raccogliamo quindi la probabilità che Jane superi l'esame con 16, 17, 18, 19 o 20 domande totali:
$$ P(X \leq 20) = \sum_{x=16}^{20} \binom{x-1}{15} (0.567)^{16} (0.433)^{x-16} \approx 0.0269213 $$
La risposta è: 0.0269213


% ==========================================
%               ESERCIZIO 5
% ==========================================
\vspace{0.5cm}
\begin{tcolorbox}[colback=lightergray!10, colframe=tocsec, title=\textbf{Esercizio 5}]
In media, in un autolavaggio arrivano 5 macchine in un’ora. Sia $X$ la v.c. che descrive il numero di auto che arrivano tra le 7 e le 8 di mattina. $X \sim Pois(\lambda=5)$.
\end{tcolorbox}

\subsubsection*{Quesito 1}
1. Qual è la probabilità che non arrivino macchine nell’ora indicata?

\textbf{Soluzione:}
Per una variabile distribuita secondo il modello di Poisson, la probabilità di osservare un numero esatto di eventi $x$ è data dalla massa discreta $P(X=x) = \frac{\lambda^x e^{-\lambda}}{x!}$. Fissando $x=0$:
$$ P(X=0) = \frac{5^0 e^{-5}}{0!} = e^{-5} \approx 0.0067379 $$
La risposta è: 0.0067379

\subsubsection*{Quesito 2}
2. Qual è la probabilità che arrivino esattamente 3 macchine nell’ora indicata?

\textbf{Soluzione:}
Sfruttando la medesima formula e ponendo le occorrenze a $x=3$:
$$ P(X=3) = \frac{5^3 e^{-5}}{3!} = \frac{125 e^{-5}}{6} \approx 0.1403739 $$
La risposta è: 0.1403739

\subsubsection*{Quesito 3}
3. Qual è la probabilità che arrivino almeno 7 e non più di 10 macchine?

\textbf{Soluzione:}
Dobbiamo calcolare la probabilità cumulata per l'intervallo ristretto compreso tra 7 e 10 (estremi inclusi), ovvero la sommatoria delle singole probabilità puntuali:
$$ P(7 \leq X \leq 10) = \sum_{x=7}^{10} \frac{5^x e^{-5}}{x!} \approx 0.2241213 $$
La risposta è: 0.2241213


\newpage

% ==========================================
%               ESERCIZIO 6
% ==========================================
\vspace{0.5cm}
\begin{tcolorbox}[colback=lightergray!10, colframe=tocsec, title=\textbf{Esercizio 6}]
In una famiglia la probabilità che nasca un maschio è 0.22. Possiamo rappresentare il numero di figli maschi nella famiglia come una v.a. binomiale $X \sim Bin(n,p)$ con $n=6$ e $P(\text{figlio maschio})=P(X=1)=p=0.22$.
Sapendo che il numero di figli nella famiglia è 6, determinare la probabilità che:
\end{tcolorbox}


\subsubsection*{Quesito 1}
1. almeno uno sia maschio

\textbf{Soluzione:}
Utilizziamo la regola dell'evento complementare per semplificare le operazioni algebriche. L'evento "almeno uno" ($X \geq 1$) è il complemento dell'evento "nessuno" ($X = 0$).
$$ P(X \geq 1) = 1 - P(X=0) $$
Sostituendo nella funzione di massa di probabilità della Binomiale:
$$ 1 - \binom{6}{0} (0.22)^0 (0.78)^6 = 1 - (0.78)^6 \approx 1 - 0.2251996 = 0.7748004 $$
La risposta è: 0.7748004

\subsubsection*{Quesito 2}
2. ci siamo (esattamente) 6 femmine

\textbf{Soluzione:}
Sotto le assunzioni del nostro modello, l'evento equivale logicamente a chiedere che ci siano esattamente 0 maschi, i.e. $P(X=0)$.
$$ P(X=0) = \binom{6}{0} (0.22)^0 (0.78)^6 = (0.78)^6 \approx 0.2251996 $$
La risposta è: 0.2251996

\subsubsection*{Quesito 3}
3. non ci siano più di 5 maschi

\textbf{Soluzione:}
In termini analitici, ciò richiede la probabilità $P(X \leq 5)$.
Anche in questo caso ci appoggiamo alle proprietà dell'evento complementare: il contrario di "non più di 5 maschi" su un totale di 6 figli è averne "esattamente 6 maschi".
$$ P(X \leq 5) = 1 - P(X=6) $$
$$ 1 - \binom{6}{6} (0.22)^6 (0.78)^0 = 1 - (0.22)^6 \approx 1 - 0.0001134 = 0.9998866 $$
La risposta è: 0.9998866

\newpage

% ==========================================
%               ESERCIZIO 7
% ==========================================
\vspace{0.5cm}
\begin{tcolorbox}[colback=lightergray!10, colframe=tocsec, title=\textbf{Esercizio 7}]
Si consideri la funzione $f$ così definita: $f(x)=cx^4$ per $0 \leq x \leq 1$ e $f(x)=0$ altrove, dove $c>0$ è una costante.
\end{tcolorbox}

\subsubsection*{Quesito 1}
1. Funzione di densità. Si determini $c>0$ in modo che $f(x)$ sia la funzione di densità di una variabile casuale $X$.

\textbf{Soluzione:}
Per essere una funzione di densità di probabilità deve integrare a 1 sul suo dominio, i.e.
$$ 1 = \int_0^1 f(x) dx = c \int_0^1 x^4 dx = c \left[ \frac{x^5}{5} \right]_0^1 = \frac{c}{5} $$
da cui si ricava immediatamente $c=5$.

La risposta è: 5

\vspace{0.3cm}
\textit{D’ora in poi, sia $X$ una v.c. con densità $f(x)$, in cui si usa la costante $c$ trovata sopra.}

\subsubsection*{Quesito 2}
2. Determinare il valor medio (anche detto aspettazione o valore atteso) di $X$.

\textbf{Soluzione:}
Applichiamo la definizione di valore atteso per variabili aleatorie continue:
$$ \mathbb{E}(X) = \int_0^1 x f(x) dx = (5) \int_0^1 x^5 dx = 5 \left[ \frac{x^6}{6} \right]_0^1 = \frac{5}{6} $$
La risposta è: 0.8333333

\subsubsection*{Quesito 3}
3. Determinare la varianza di $X$.

\textbf{Soluzione:}
La varianza è definita come $Var(X) = \int_0^1 \left(x - \frac{5}{6}\right)^2 f(x) dx$. 
Indichiamo con $k=4$ l’esponente di $x$, per cui $f(x) = (k+1)x^k$ e $\mathbb{E}(X) = \frac{k+1}{k+2}$. 
Sviluppando l'integrale del binomio al quadrato:
$$ Var(X) = \int_0^1 \left(x - \frac{k+1}{k+2}\right)^2 f(x) dx = \int_0^1 \left( x^2 - 2\frac{k+1}{k+2}x + \left(\frac{k+1}{k+2}\right)^2 \right) f(x) dx $$
Vediamo i vari pezzi separatamente sfruttando la linearità dell'integrale:
$$ \int_0^1 x^2 f(x) dx = (k+1) \int_0^1 x^{k+2} dx = \frac{k+1}{k+3} = \frac{5}{7} $$
$$ -2\frac{k+1}{k+2} \int_0^1 x f(x) dx = -2\frac{k+1}{k+2} \mathbb{E}(X) = -2\mathbb{E}(X)^2 $$
$$ \int_0^1 \left(\frac{k+1}{k+2}\right)^2 f(x) dx = \left(\frac{k+1}{k+2}\right)^2 \int_0^1 f(x) dx = \mathbb{E}(X)^2 $$
Sommando i termini abbiamo ritrovato la familiare uguaglianza:
$$ Var(X) = \mathbb{E}(X^2) - \mathbb{E}(X)^2 = \frac{5}{7} - \left(\frac{5}{6}\right)^2 = \frac{5}{7} - \frac{25}{36} = \frac{180 - 175}{252} = \frac{5}{252} \approx 0.0198413 $$
La risposta è: 0.0198413

\subsubsection*{Quesito 4}
4. Determinare la probabilità $P(X \leq 0.24)$.

\textbf{Soluzione:}
Ciò equivale a valutare la funzione di ripartizione $F(x)$ nel punto $x=0.24$:
$$ P(X \leq 0.24) = F(0.24) = \int_0^{0.24} f(x) dx = \left[ x^5 \right]_0^{0.24} = (0.24)^5 $$
La risposta è: 0.0039813


% ==========================================
%               ESERCIZIO 8
% ==========================================
\vspace{0.5cm}
\begin{tcolorbox}[colback=lightergray!10, colframe=tocsec, title=\textbf{Esercizio 8}]
Quanto a lungo vive un certo organismo biologico, può essere rappresentato da una variabile casuale (continua) $X$ con funzione di densità di probabilità 
$$ f(x) = 2\lambda x \exp\{-\lambda x^2\} \quad \text{per } \lambda, x > 0 $$
con $\lambda=1.29$.
\end{tcolorbox}

\vspace{0.3cm}
\textbf{Soluzione preliminare:} Determiniamo la funzione di ripartizione di $X$, 
$$ F(x) = \int_0^x f(t) dt = \int_0^x 2\lambda t \exp\{-\lambda t^2\} dt = 1 - \exp\{-\lambda x^2\} $$
per rispondere alle domande 1 e 2 usare il fatto che $P(a < X \leq b) = F(b) - F(a)$.

\subsubsection*{Quesito 1}
1. Calcolare $P(X \leq 0.4)$.

\textbf{Soluzione:}
Sostituiamo il valore direttamente nella funzione di ripartizione $F(x)$ precedentemente trovata:
$$ P(X \leq 0.4) = F(0.4) = 1 - \exp\{-1.29 \cdot 0.4^2\} = 1 - \exp\{-0.2064\} \approx 0.1864924 $$
La risposta è: 0.1864924

\subsubsection*{Quesito 2}
2. Calcolare $P(0.24 \leq X \leq 1.99)$.

\textbf{Soluzione:}
Sfruttiamo l'operazione di differenza sulla funzione di ripartizione:
$$ P(0.24 \leq X \leq 1.99) = F(1.99) - F(0.24) \approx 0.9223445 $$
La risposta è: 0.9223445

\newpage

\subsubsection*{Quesito 3}
3. Momenti. Ora si consideri una v.c. $Y$ con funzione di densità $f(y) = \alpha \exp\{-\alpha y\}$ per $y \geq 0$ con $\alpha=1.05$. Si osservi che $Y \sim E(\alpha)$ e $\mathbb{E}(Y^n) = \frac{n!}{\alpha^n}$.

\textbf{Soluzione (momento di ordine 1):}
Calcolare il momento (non centrato) di ordine 1 di $Y$:
$$ \mathbb{E}(Y) = \frac{1!}{\alpha^1} = \frac{1}{1.05} \approx 0.952381 $$
La risposta è: 0.952381

\textbf{Soluzione (momento di ordine 2):}
Calcolare il momento (non centrato) di ordine 2 di $Y$:
$$ \mathbb{E}(Y^2) = \frac{2!}{\alpha^2} = \frac{2}{1.05^2} = \frac{2}{1.1025} \approx 1.814059 $$
La risposta è: 1.814059


% ==========================================
%               ESERCIZIO 9
% ==========================================
\vspace{0.5cm}
\begin{tcolorbox}[colback=lightergray!10, colframe=tocsec, title=\textbf{Esercizio 9}]
Si consideri una variabile casuale $X$ con funzione di densità 
$$ f(x) = \begin{cases} k(8.9 - 2x) & \text{se } 0 < x < 4 \\ 0 & \text{altrimenti} \end{cases} $$
\end{tcolorbox}

\subsubsection*{Quesito 1}
1. Determinare la costante $k$ affinchè $f(x)$ sia funzione di densità.

\textbf{Soluzione:}
Affinchè $f(x)$ sia una funzione di densità, l’integrale della $f$ sul suo supporto – l’intervallo $(0,4)$ – deve essere uguale a 1, $\int_0^4 f(x) dx = 1$.
$$ \begin{aligned} 1 &= \int_0^4 k(8.9 - 2x) dx \\ &= k \int_0^4 (8.9 - 2x) dx \\ &= k \left( \int_0^4 8.9 dx - \int_0^4 2x dx \right) \\ &= k \left( \left[ 8.9x \right]_0^4 - \left[ x^2 \right]_0^4 \right) \end{aligned} $$
Sostituendo i limiti di integrazione otteniamo $k(8.9 \cdot 4 - 4^2) = k(35.6 - 16) = 19.6k$.
Ponendo $19.6k = 1$ otteniamo $k = \frac{1}{19.6} \approx 0.0510204$.

La risposta è: 0.0510204

\subsubsection*{Quesito 2}
2. Calcolare il valore atteso di $X$.

\textbf{Soluzione:}
$$ \mathbb{E}(X) = \int_0^4 x \cdot k(8.9 - 2x) dx = k \int_0^4 (8.9x - 2x^2) dx = k \left[ 8.9 \frac{x^2}{2} - 2 \frac{x^3}{3} \right]_0^4 $$
$$ \mathbb{E}(X) = k \left( 8.9 \cdot 8 - 2 \cdot \frac{64}{3} \right) = k (71.2 - 42.666) \approx 1.4557823 $$
La risposta è: 1.4557823

\subsubsection*{Quesito 3}
3. Calcolare il valore atteso di $14.8 + (6.1)X$.

\textbf{Soluzione:}
Usiamo la linearità dell'operatore valore atteso $\mathbb{E}(aX+b) = a\mathbb{E}(X)+b$:
$$ \mathbb{E}(14.8 + 6.1X) = 14.8 + 6.1 \cdot \mathbb{E}(X) = 14.8 + 6.1(1.4557823) \approx 23.6802721 $$
La risposta è: 23.6802721

\subsubsection*{Quesito 4}
4. Dato il momento non centrato di ordine 2, calcolare la varianza di $X$:
$$ \mathbb{E}(X^2) = 3.1564626 $$

\textbf{Soluzione:}
Usiamo la formula della varianza basata sui momenti:
$$ Var(X) = \mathbb{E}(X^2) - \mathbb{E}(X)^2 = 3.1564626 - (1.4557823)^2 \approx 1.0371604 $$
La risposta è: 1.0371604


% ==========================================
%               ESERCIZIO 10
% ==========================================
\vspace{0.5cm}
\begin{tcolorbox}[colback=lightergray!10, colframe=tocsec, title=\textbf{Esercizio 10}]
L’altezza media di una donna in Italia è distribuita normalmente con media $\mu=162.5$ cm (fonti) e scarto quadratico medio, anche detta deviazione standard, $\sigma=4.98$ (valore inventato). Indichiamo con $X \sim N(\mu, \sigma^2)$ la corrispondente variabile casuale.
\end{tcolorbox}

\vspace{0.3cm}
\textbf{Soluzione preliminare:} Servendosi dell’operazione di standardizzazione rispondere ai seguenti quesiti. 
$X \sim N(\mu=162.5, \sigma^2=24.8004)$ per cui $Z := \frac{X-\mu}{\sigma} \sim N(0,1)$. In questo modo si ha che $P(X < x) = P\left(Z < \frac{x-\mu}{\sigma}\right)$.

\subsubsection*{Quesito 1}
1. Si vuole determinare la probabilità che una donna abbia un’altezza inferiore a 169.223. Qual è il valore standardizzato corrispondente a 169.223?

\textbf{Soluzione:}
$$ z = \frac{169.223 - 162.5}{4.98} = \frac{6.723}{4.98} = 1.35 $$
La risposta è: 1.35

\subsubsection*{Quesito 2}
2. Calcolare $P(X \leq 169.223)$ attraverso standardizzazione.

\textbf{Soluzione:}
Utilizziamo il quantile calcolato al punto precedente:
$$ P(X \leq 169.223) = P(Z \leq 1.35) = \Phi(1.35) \approx 0.9115 $$
La risposta è: 0.9115

\subsubsection*{Quesito 3}
3. Calcolare $P(X \geq 172.0118)$ attraverso standardizzazione.

\textbf{Soluzione:}
Troviamo innanzitutto il valore z-score:
$$ z = \frac{172.0118 - 162.5}{4.98} = 1.91 $$
Calcoliamo quindi la probabilità della coda superiore:
$$ P(X \geq 172.0118) = P(Z \geq 1.91) = 1 - \Phi(1.91) \approx 1 - 0.9719 = 0.0281 $$
La risposta è: 0.0281


% ==========================================
%               ESERCIZIO 11
% ==========================================
\vspace{0.5cm}
\begin{tcolorbox}[colback=lightergray!10, colframe=tocsec, title=\textbf{Esercizio 11}]
Delle caramelle artigianali hanno un peso distribuito come una normale con media $\mu=5.85$ g e scarto quadratico medio $\sigma=0.57$ g. Al controllo qualità le caramelle con un peso superiore a 6.6508908 g o inferiore a 5.0857696 g vengono scartate.
Sia $X \sim N(\mu=5.85, \sigma^2=0.3249)$ la v.a. che rappresenta il peso di una caramella artigianale.
\end{tcolorbox}

\subsubsection*{Quesito 1}
1. Qual è la probabilità che una caramella sia sopra il peso soglia 6.6508908 g?

\textbf{Soluzione:}
E’ richiesta la probabilità $P(X > 6.6508908) = 1 - P(X \leq 6.6508908)$.
Standardizziamo il valore della soglia superiore:
$$ z = \frac{6.6508908 - 5.85}{0.57} = 1.40507 $$
Pertanto, $P(Z > 1.405) = 1 - \Phi(1.405) \approx 0.08$.

La risposta è: 0.08

\subsubsection*{Quesito 2}
2. In media, ogni quante caramelle se ne presenta una da scartare?

\textbf{Soluzione:}
La probabilità totale che una caramella sia scartata è la somma delle probabilità delle due code di reiezione: 
$$ p = P(X > 6.6508908) + P(X < 5.0857696) = 0.17 $$
Si noti che $0.17 = \frac{17}{100}$, ossia ogni 100 caramelle 17 sono da scartare, da cui (ad esempio facendo una proporzione $17:100 = 1:x$) il valore desiderato si ottiene calcolando l'inverso della probabilità:
$$ x = \frac{1}{p} = \frac{1}{0.17} \approx 5.8823529 $$
La risposta è: 5.8823529

\subsubsection*{Quesito 3}
3. Qual è il valore standardizzato corrispondente a 5.0857696?

\textbf{Soluzione:}
Usiamo la formula della standardizzazione $Z := \frac{X-\mu}{\sigma}$ per quanto riguarda il valore della soglia inferiore:
$$ z = \frac{5.0857696 - 5.85}{0.57} = \frac{-0.7642304}{0.57} \approx -1.340755 $$
La risposta è: -1.340755

\section*{Probabilità e Statistica - Esercizi 6}
\addcontentsline{toc}{section}{Probabilità e Statistica - Esercizi 6}
\markboth{Esercizi 6}{Esercizi 6}

% ==========================================
%               ESERCIZIO 1
% ==========================================
\begin{tcolorbox}[colback=lightergray!10, colframe=tocsec, title=\textbf{Esercizio 1}]
Sia $X \sim N(\mu,\sigma^2)$ con $\mu=6.7$, $\sigma^2=3.24$.
\end{tcolorbox}

\subsection*{Quesiti e soluzioni}

\subsubsection*{Quesito 1}
1. Quanto vale il secondo momento non centrato della variabile $X$?

\textbf{Soluzione:}
Sfruttiamo la definizione della varianza in termini di momenti. Sappiamo che $Var(X) = \mathbb{E}(X^2) - \mathbb{E}^2(X)$, da cui possiamo isolare il momento non centrato di ordine 2:
$$ \mathbb{E}(X^2) = Var(X) + \mathbb{E}^2(X) = \sigma^2 + \mu^2 $$
Sostituendo i valori forniti dal problema:
$$ \mathbb{E}(X^2) = 3.24 + (6.7)^2 = 3.24 + 44.89 = 48.13 $$
La risposta è: 48.13

\subsubsection*{Quesito 2}
2. Trasformazione di una v.c.
Si consideri ora la trasformazione $Y=g(X)=\left(\dfrac{X-\mu}{\sigma}\right)^2$
Qual è la media di $Y$?

\textbf{Soluzione:}
Si osservi che $Z = \dfrac{X-\mu}{\sigma} \sim N(0,1)$ è una normale standardizzata e $Y=Z^2$, per cui $\mathbb{E}(Y) = \mathbb{E}(Z^2)$, ossia il momento non centrato di ordine 2 di una v.a. normale standard. 
Usando la stessa osservazione del punto precedente: 
$$ \mathbb{E}(Z^2) = Var(Z) + \mathbb{E}^2(Z) = Var(Z) = 1 $$
La risposta è: 1

\newpage

\subsubsection*{Quesito 3}
3. Qual è la varianza di $Y$?

\textbf{Soluzione:}
Incolonniamo i passaggi sfruttando le proprietà dei momenti della normale standard:
$$ \begin{aligned} Var(Y) &= \mathbb{E}(Y^2) - \mathbb{E}(Y)^2 \\ &= \mathbb{E}(Z^4) - 1^2 \\ &= \int_{-\infty}^{+\infty} z^4 f_Z(z) dz - 1 = \dots = 3 - 1 = 2 \end{aligned} $$
ove $f_Z(z)$ è la densità gaussiana standard e $\mathbb{E}(Z^4) = 3$ è un risultato notevole noto.

Si osservi che con questo esercizio abbiamo verificato (empiricamente) la seguente proprietà delle v.a. normali standard:
Siano $X_1,\dots,X_n$ i.i.d. da $N(0,1)$, allora $Q := \sum_{i=1}^n X_i^2 \sim \chi^2_n$, ossia $Q$ è distribuita come una chi-quadro con $n$ gradi di libertà.
Inoltre $\mathbb{E}(Q) = n$ e $Var(Q) = 2n$. Nel nostro caso $n=1$.

La risposta è: 2

\subsubsection*{Quesito 4}
4. Scrivere la funzione di densità della variabile $Z=\dfrac{X-\mu}{\sigma}$, dove $X$ è la variabile del punto 1.

\textbf{Soluzione:}
Trattandosi di una variabile aleatoria normale standard $Z \sim N(0,1)$, la sua funzione di densità è:
$$ f_Z(z) = \frac{1}{\sqrt{2\pi}} \exp\left(-\frac{z^2}{2}\right) $$
Codice R corrispondente: \texttt{1 / sqrt(2 * pi) * exp( - (z \textasciicircum{} 2) / 2)}


% ==========================================
%               ESERCIZIO 2
% ==========================================
\vspace{0.5cm}
\begin{tcolorbox}[colback=lightergray!10, colframe=tocsec, title=\textbf{Esercizio 2}]
Sapendo che il 17\% delle persone che acquistano un biglietto online per un film poi non si presenta al cinema, i gestori del sito di booking di un cinema vendono 14 biglietti per il film proiettato nella sala 1, che ha 12 posti e 18 biglietti per quello nella sala 2, da 15 posti.
Conviene definire ``successo'' l’evento ``La persona che ha acquistato il biglietto si presenta al cinema''. La probabilità di successo è $p=0.83$. Chiamiamo inoltre $X_1$ la v.a. binomiale con distribuzione $Bin(N_1=14,p)$ e $X_2 \sim Bin(N_2=18,p)$.
\end{tcolorbox}

\subsubsection*{Quesito 1}
1. Calcolare la probabilità che una persona che ha comprato il biglietto online (per il primo film) non trovi posto nella sala 1, da 12 posti.

\textbf{Soluzione:}
La domanda equivale a ``qual è la probabilità che vi siano più successi (persone che si presentano) di posti nella sala''? In termini matematici viene chiesto di determinare $P(X_1 > 12)$ che si ottiene come la somma delle probabilità per i valori superiori alla capienza:
$$ P(X_1 > 12) = \sum_{k=13}^{14} \binom{N_1}{k} p^k (1-p)^{N_1-k} = \binom{14}{13}(0.83)^{13}(0.17)^1 + \binom{14}{14}(0.83)^{14}(0.17)^0 \approx 0.284787 $$
La risposta è: 0.284787

\subsubsection*{Quesito 2}
2. Calcolare la probabilità che una persona che ha comprato il biglietto online (per il secondo film) non trovi posto nella sala 2, da 15 posti.

\textbf{Soluzione:}
Idem $P(X_2 > 15)$, cambiando i parametri della distribuzione:
$$ P(X_2 > 15) = \sum_{k=16}^{18} \binom{18}{k} (0.83)^k (0.17)^{18-k} \approx 0.388091 $$
La risposta è: 0.388091

\subsubsection*{Quesito 3}
3. Calcolare il numero di biglietti che si potrebbero vendere per la sala con 12 posti accettando che la probabilità di scontentare una persona (i.e. lasciarla fuori dalla sala) non sia superiore a 0.38.

\textbf{Soluzione:}
Sia $\alpha=0.38$. Come abbiamo già visto in esercizi precedenti, si tratta del ``problema inverso'': fissata la probabilità determinare il valore assunto dalla v.a.
$$ \alpha = P(X_1 > n[1]) = \sum_{k=13}^{N_1} \binom{N_1}{k} p^k (1-p)^{N_1-k} $$
Provando ad aumentare iterativamente il parametro $N_1$, notiamo che con 14 biglietti venduti la probabilità di overbooking è $\sim 0.28$ (accettabile), mentre con 15 biglietti supererebbe la soglia imposta di 0.38.
La risposta è: 14

\subsubsection*{Quesito 4}
4. In media quante persone si presentano per il film della sala 2?

\textbf{Soluzione:}
Si richiede il valore atteso di una v.a. binomiale, che è definito dal prodotto i.e. $N_2 \cdot p$:
$$ \mathbb{E}(X_2) = 18 \cdot 0.83 = 14.94 $$
La risposta è: 14.94

\newpage

% ==========================================
%               ESERCIZIO 3
% ==========================================
\vspace{0.5cm}
\begin{tcolorbox}[colback=lightergray!10, colframe=tocsec, title=\textbf{Esercizio 3}]
Consideriamo il seguente esperimento casuale: lanciamo un dado (equilibrato) a 9 facce finché non esce 9; quando esce 9 ci fermiamo. Sia $N$ la variabile aleatoria che rappresenta il numero di lanci e sia $X$ la v.a. che conta quante volte è uscita la faccia con il numero 4.
\end{tcolorbox}

\subsection*{Quesiti e soluzioni}

\subsubsection*{Quesito 1}
Qual è il valore atteso di $X$ condizionato a $N=11$, $\mathbb{E}[X|N=11]$?

\textbf{Soluzione:}
Sapendo che ci sono stati 11 lanci, questo vuol dire che l’ultimo è stato un 9, mentre nei precedenti 10 lanci abbiamo ottenuto valori diversi da 9. Tra questi, solo la faccia con il numero 4 costituisce un successo, e ha probabilità $p = \dfrac{1}{8} = 0.125$.
La v.a. $X|N=11$ è quindi una binomiale di parametri $n=10$ e $p=0.125$. Il suo valore atteso è quindi $np = 10 \cdot \dfrac{1}{8}$.
La risposta è: 1.25

\subsubsection*{Quesito 2}
Quanto vale $\mathbb{E}[X^2|N=11]$?

\textbf{Soluzione:}
Sfruttando nuovamente il fatto che $X|N=11$ ha una distribuzione binomiale di parametri $n=10$ e $p=0.125$, e sapendo che la varianza di una variabile casuale binomiale è $np(1-p)$, possiamo scrivere:
$$ np(1-p) = Var(X|N=11) = \mathbb{E}(X^2|N=11) - (\mathbb{E}(X|N=11))^2 $$
da cui
$$ \mathbb{E}(X^2|N=11) = np(1-p) + (np)^2 $$
Sostituendo i valori: $10(0.125)(0.875) + (1.25)^2 = 1.09375 + 1.5625 = 2.65625$.
La risposta è: 2.65625

\subsubsection*{Quesito 3}
Scrivere la funzione di probabilità $P(z(N)=x)$ dove $z(n)=\mathbb{E}(X|N=n)$.

\textbf{Soluzione:}
Sappiamo, dal quesito 1, che $z(n) = \mathbb{E}(X|N=n) = (n-1)\dfrac{1}{8}$ per cui possiamo riscrivere la v.a. $z(N) = \mathbb{E}(X|N)$ come trasformazione (affine) della variabile $N$:
$$ z(N) = (N-1)\frac{1}{8} $$
Ora $N$ conta il numero di prove per avere un successo: è quindi una v.a. geometrica di parametro $p^*=\dfrac{1}{9}$ (poiché il dado è equilibrato) e sappiamo che 
$$ P(N=k) = p^*(1-p^*)^{k-1} \quad \text{per } k \in R_N=\{1,2,\dots\} $$
mentre è nulla altrimenti. Il supporto di $z(N)$ è quindi $R_{z(N)} = \{ \frac{k-1}{8} : k \in R_N \}$ e 
$$ P(\{z(N)=x\}) = P\left(\left\{(N-1)\frac{1}{8}=x\right\}\right) = P\left(N = 1 + x\left(\frac{1}{8}\right)^{-1}\right) $$
per $x \in R_{z(N)}$ o, alternativamente per $1+x\left(\frac{1}{8}\right)^{-1} \in R_N$, e nulla altrimenti.


% ==========================================
%               ESERCIZIO 4
% ==========================================
\vspace{0.5cm}
\begin{tcolorbox}[colback=lightergray!10, colframe=tocsec, title=\textbf{Esercizio 4}]
Sia data la seguente variabile aleatoria bivariata discreta 
\begin{center}
\begin{tabular}{c|ccc}
\toprule
$X \setminus Y$ & -0.44 & -0.07 & 1.7 \\
\midrule
-0.57 & 3k & 8k & 1k \\
0.12 & 4k & 5k & 6k \\
2.3 & 2k & 7k & 9k \\
\bottomrule
\end{tabular}
\end{center}
dove sulle righe abbiamo la variabile $X$ e sulle colonne la variabile $Y$.
\end{tcolorbox}

\subsection*{Quesiti e soluzioni}

\subsubsection*{Quesito 1}
Calcolare la costante $k$ che deve essere utilizzata per rendere la tabella una funzione di probabilità congiunta.

\textbf{Soluzione:}
Abbiamo che la somma per righe e colonne nella tabella sopra devono dare 1. La somma totale di tutte le celle è $(3+8+1) + (4+5+6) + (2+7+9) = 12 + 15 + 18 = 45$, i.e. $45k = 1$, da cui:
$$ k = \frac{1}{45} \approx 0.0222222 $$
La risposta è: 0.0222222

\subsubsection*{Quesito 2}
Calcolare la funzione di probabilità marginale di $X$. Inserire il solo valore di probabilità per $X=-0.57$.

\textbf{Soluzione:}
Data la probabilità congiunta $p_{X,Y}(x,y)=P(X=x,Y=y)$ è facile ottenere le marginali.
Indichiamo con $R_Y$ il range o immagine di $Y$, i.e. l’insieme di valori assunti da $Y$, 
$$ p_X(x) = \sum_{y \in R_Y} p_{X,Y}(x,y) $$
Sommando la prima riga della matrice corrispondente a $X=-0.57$:
$$ P(X=-0.57) = 3k + 8k + 1k = 12k = 12 \left(\frac{1}{45}\right) = \frac{12}{45} \approx 0.2666667 $$
La risposta è: 0.2666667

\newpage

\subsubsection*{Quesito 3}
Calcolare la funzione di probabilità marginale di $Y$. Inserire il solo valore di probabilità per $Y=-0.07$.

\textbf{Soluzione:}
Applicando la somma sulle colonne della matrice:
$$ p_Y(y) = \sum_{x \in R_X} p_{X,Y}(x,y) $$
Per la colonna corrispondente a $Y=-0.07$:
$$ P(Y=-0.07) = 8k + 5k + 7k = 20k = 20 \left(\frac{1}{45}\right) = \frac{20}{45} \approx 0.4444444 $$
La risposta è: 0.4444444

\subsubsection*{Quesito 4}
Qual è il valore atteso di $X$?

\textbf{Soluzione:}
Calcoliamo tutte le marginali di $X$: $p_X(-0.57)=\frac{12}{45}$, $p_X(0.12)=\frac{15}{45}$, $p_X(2.3)=\frac{18}{45}$.
Applicando la definizione di valore atteso per V.A. discrete:
$$ \mathbb{E}(X) = \sum_{x \in R_X} x p_X(x) = (-0.57)\frac{12}{45} + (0.12)\frac{15}{45} + (2.3)\frac{18}{45} = \frac{-6.84 + 1.8 + 41.4}{45} = \frac{36.36}{45} = 0.808 $$
La risposta è: 0.808

\subsubsection*{Quesito 5}
Qual è la varianza di $Y$?

\textbf{Soluzione:}
Calcoliamo tutte le marginali di $Y$: $p_Y(-0.44)=\frac{9}{45}$, $p_Y(-0.07)=\frac{20}{45}$, $p_Y(1.7)=\frac{16}{45}$.
Troviamo la media di $Y$:
$$ \mathbb{E}(Y) = (-0.44)\frac{9}{45} + (-0.07)\frac{20}{45} + (1.7)\frac{16}{45} = \frac{21.84}{45} \approx 0.4853333 $$
Calcoliamo il secondo momento non centrato:
$$ \mathbb{E}(Y^2) = (-0.44)^2 \frac{9}{45} + (-0.07)^2 \frac{20}{45} + (1.7)^2 \frac{16}{45} = \frac{48.0804}{45} \approx 1.0684533 $$
Quindi la varianza è $Var(Y) = \mathbb{E}(Y^2) - \mathbb{E}(Y)^2 = 1.0684533 - (0.4853333)^2 \approx 0.8329049$.
La risposta è: 0.8329049

\newpage

% ==========================================
%               ESERCIZIO 5
% ==========================================
\vspace{0.5cm}
\begin{tcolorbox}[colback=lightergray!10, colframe=tocsec, title=\textbf{Esercizio 5}]
Consideriamo una moneta non equilibrata, per cui la probabilità di avere testa è $p=0.422$. Lanciamo un dado a 6 facce e poi la moneta, tante volte quante il valore uscito nel lancio del dado. Sia $D$ la variabile aleatoria che ci dà l’esito del dado e $N$ la variabile aleatoria che conta il numero di teste uscite.
\end{tcolorbox}

\subsection*{Quesiti e soluzioni}

\subsubsection*{Quesito 1}
Qual è il valore atteso di $N$?

\textbf{Soluzione:}
Per calcolare il valore atteso, o momento primo, ci occorre innanzitutto conoscere le probabilità discrete marginali di $N$, $p_N(n)=P(N=n)$, per $n=0,\dots,6$, infatti la definizione di valore atteso per una variabile aleatoria discreta a valori in $0,\dots,6$ è:
$$ \mathbb{E}[N] = \sum_{n=0}^6 n \cdot P(N=n). $$
Ricordiamo che la densità discreta marginale di $N$, $p_N$, si ottiene come:
$$ p_N(n) = P(N=n) = \sum_{d \in R_D} p_{D,N}(d,n), $$
dove $p_{D,N}(d,n)=P(D=d,N=n)$ è la densità discreta congiunta di $D$ e $N$, e $R_D=\{1,\dots,6\}$ sono i possibili esiti del lancio del dado.
Quindi per ricavarci le probabilità marginali di $N$ dobbiamo prima calcolare la probabilità discreta congiunta di $D$ e $N$:
$$ p_{D,N}(d,n) = P(D=d)P(N=n|D=d) = \frac{1}{6} \cdot \binom{d}{n} p^n (1-p)^{d-n} $$
dove $p=0.422$. Osserviamo infatti che, per quanto riguarda la probabilità condizionata $P(N=n|D=d)$, se lanciando il dado otteniamo la faccia con il numero $d$, lanceremo la moneta $d$ volte e dunque conteremo il numero di successi (teste) sui $d$ tentativi (lanci), cioè abbiamo una distribuzione binomiale di parametri $d$ e $p$.
Sfruttando la legge delle aspettative iterate (Teorema della Speranza Matematica Totale), $\mathbb{E}(N) = \mathbb{E}[\mathbb{E}(N|D)] = \mathbb{E}[D \cdot p] = p\mathbb{E}[D] = 0.422 \cdot 3.5 = 1.477$.
La risposta è: 1.477

\subsubsection*{Quesito 2}
Qual è la media (valore atteso) di $N$ sapendo che $D=5$?

\textbf{Soluzione:}
Qui ci interessa la media di $N$ condizionata all’evento $D=5$. Come abbiamo visto, conoscendo il valore di $D=5$ sappiamo che il numero di successi $N$ si distribuisce come una binomiale di parametri $d=5$ e $p=0.422$, di cui il valore atteso è noto, cioè:
$$ \mathbb{E}[N|D=d] = d \cdot p. $$
In alternativa possiamo usare la definizione di valore atteso per una variabile aleatoria discreta, con l’accortezza di usare $p_{N|D=d}$ come densità discreta.
Sostituendo i valori: $\mathbb{E}[N|D=5] = 5 \cdot 0.422 = 2.11$.
La risposta è: 2.11

\subsubsection*{Quesito 3}
Sapendo che alla fine dell’esperimento abbiamo avuto 5 teste, qual è la media del risultato del lancio del dado?

\textbf{Soluzione:}
Dobbiamo calcolare:
$$ \mathbb{E}[D|N=5] = \sum_{d=1}^6 d \cdot P(D=d|N=5), $$
dove possiamo usare la densità discreta congiunta e la marginale calcolate per i punti precedenti, dato che (per il Teorema di Bayes):
$$ P(D=d|N=n) = \frac{P(D=d,N=n)}{P(N=n)} $$
Il calcolo produce una massa di probabilità a posteriori concentrata sui valori 5 e 6 del dado. La media pesata risultante sarà:
La risposta corretta è: 5.7761862


% ==========================================
%               ESERCIZIO 6
% ==========================================
\vspace{0.5cm}
\begin{tcolorbox}[colback=lightergray!10, colframe=tocsec, title=\textbf{Esercizio 6}]
Sia $X$ la variabile aleatoria che descrive il lancio di un dado equo con le seguenti 20 facce: 1, 1, 1, 1, 2, 2, 3, 3, 4, 4, 5, 6, 6, 6, 6, 7, 8, 9, 9, 10. Una volta lanciato il dado, tiriamo una moneta (probabilità di testa $p=0.41$) tante volte quante indicato dal dado e contiamo quante teste sono uscite. Sia $Y$ la variabile aleatoria che rappresenta il numero di teste.
\end{tcolorbox}

\subsection*{Quesiti e soluzioni}

\subsubsection*{Quesito 1}
Qual è la probabilità che $Y=4$?

\textbf{Soluzione:}
È sufficiente applicare la legge delle probabilità totali: 
$$ P(Y=4) = \sum_{i=4}^{10} P(X=i)P(Y=4|X=i) $$
e usare il fatto che $P(X=i) = \dfrac{n_i}{20}$ e $P(Y=4|X=i) = \binom{i}{4} p^4 (1-p)^{i-4}$, dove $n_i$ è il numero di facce del dado su cui compare il numero $i$.
La risposta corretta è: 0.0966134

\subsubsection*{Quesito 2}
Qual è la probabilità che $X=2Y$?

\textbf{Soluzione:}
Anche in questo caso è sufficiente la legge delle probabilità totali. Notiamo anzitutto che $X$ deve essere necessariamente pari, quindi: 
$$ P(X=2Y) = \sum_{x=1}^{10} P(x=2Y|X=x)P(X=x) = \sum_{i=1}^5 P(2i=2Y|X=2i)P(X=2i) = \sum_{i=1}^5 P(Y=i|X=2i)\frac{n_{2i}}{20}. $$
La risposta corretta è: 0.1625297

\subsubsection*{Quesito 3}
Sapendo che $Y=4$, qual è la probabilità che $X=9$?

\textbf{Soluzione:}
Basta applicare il teorema di Bayes:
$$ P(X=9|Y=4) = \frac{P(Y=4|X=9)P(X=9)}{P(Y=4)}. $$
Il denominatore è stato calcolato nel Quesito 1, $P(X=9) = \dfrac{n_9}{20}$ e la restante probabilità è la funzione di massa della binomiale $\binom{9}{4}(0.41)^4(0.59)^5$.
La risposta corretta è: 0.2634686

% ==========================================
%               ESERCIZIO 7
% ==========================================
\vspace{0.5cm}
\begin{tcolorbox}[colback=lightergray!10, colframe=tocsec, title=\textbf{Esercizio 7}]
Sia $X$ una variabile aleatoria Poissoniana di parametro $\lambda_X$. Sappiamo che $P(X=9)=P(X=10)$.
Sia inoltre $Y$ un’altra variabile aleatoria Poissoniana, indipendente da $X$, di media (valore atteso) 8.
Sia infine $Z=X+Y$ la somma delle due variabili aleatorie precedenti, che è nota essere anch’essa una variabile di Poisson con media $\lambda_X+\lambda_Y$.
\end{tcolorbox}

\subsection*{Quesiti e soluzioni}

\subsubsection*{Quesito 1}
Qual è il valore di $\lambda_X$?

\textbf{Soluzione:}
Scriviamo la funzione di densità discreta per generici $k$ e $k+1$:
$$ P(X=k) = \frac{\lambda_X^k e^{-\lambda_X}}{k!} \quad \text{e} \quad P(X=k+1) = \frac{\lambda_X^{k+1} e^{-\lambda_X}}{(k+1)!} $$
Ora, la condizione $P(X=k)=P(X=k+1)$ con $k=9$ equivale a:
$$ \frac{\lambda_X^k e^{-\lambda_X}}{k!} = \frac{\lambda_X^{k+1} e^{-\lambda_X}}{(k+1)!} $$
Semplificando $e^{-\lambda_X}$ e $\lambda_X^k$:
$$ \frac{1}{k!} = \frac{\lambda_X}{(k+1)k!} \implies 1 = \frac{\lambda_X}{k+1} \implies \lambda_X = k+1 $$
Sostituendo $k=9$ otteniamo $\lambda_X = 10$.
La risposta è: 10

\newpage

\subsubsection*{Quesito 2}
Qual è la probabilità $P(X=13|Z=24)$?

\textbf{Soluzione:}
Suggerimento: scrivere la funzione di densità discreta condizionata in modo astratto, con parametri generici (cioè $P(X=k|Z=n)$ e, nelle densità, usare i parametri generici $\lambda_X$ e $\lambda_X+\lambda_Y$).
La media di $Y \sim Pois(\lambda_Y)$ è $\lambda_Y$, quindi $Y \sim Pois(\lambda_Y=8)$.
Per la proprietà di riproducibilità delle Poissoniane, sappiamo anche che $Z \sim Pois(18)$.
Con parametri generici $k$ e $n$, possiamo scrivere:
$$ \begin{aligned} P(X=k|Z=n) &= \frac{P(X=k,Z=n)}{P(Z=n)} \\ &= \frac{P(X=k,X+Y=n)}{P(Z=n)} \\ &= \frac{P(X=k,Y=n-k)}{P(Z=n)} \\ &= \frac{P(X=k)P(Y=n-k)}{P(Z=n)} \end{aligned} $$
in cui abbiamo usato la definizione di probabilità condizionata (nella prima uguaglianza), la definizione di $Z=X+Y$ nella seconda, l’osservazione che l’intersezione di eventi $(X=k, X+Y=n)$ si può avere solo se $Y=n-k$ e infine l’indipendenza di $X$ e $Y$.
Possiamo sviluppare ulteriormente questa identità, seguendo il suggerimento:
$$ \begin{aligned} P(X=k|Z=n) &= \frac{\frac{\lambda_X^k e^{-\lambda_X}}{k!} \cdot \frac{\lambda_Y^{n-k} e^{-\lambda_Y}}{(n-k)!}}{\frac{(\lambda_X+\lambda_Y)^n e^{-(\lambda_X+\lambda_Y)}}{n!}} \\ &= \frac{n!}{k!(n-k)!} \frac{\lambda_X^k \lambda_Y^{n-k}}{(\lambda_X+\lambda_Y)^n} \\ &= \binom{n}{k} \left(\frac{\lambda_X}{\lambda_X+\lambda_Y}\right)^k \left(\frac{\lambda_Y}{\lambda_X+\lambda_Y}\right)^{n-k} \end{aligned} $$
Inserendo i valori assegnati per i parametri ($n=24, k=13, \lambda_X=10, \lambda_Y=8$):
$$ P(X=13|Z=24) = \binom{24}{13} \left(\frac{10}{18}\right)^{13} \left(\frac{8}{18}\right)^{11} \approx 0.160221 $$
La risposta è: 0.160221

\subsubsection*{Quesito 3}
Qual è il valore atteso condizionato $\mathbb{E}(X|Z=24)$?

\textbf{Soluzione:}
Qui torna utile l’espressione derivata sopra: osserviamo che la funzione di probabilità condizionata è una legge binomiale con parametri $n$ e $p = \frac{\lambda_X}{\lambda_X+\lambda_Y}$ per cui sappiamo immediatamente che il valore atteso di $X|Z=n$ è $np = n\frac{\lambda_X}{\lambda_X+\lambda_Y}$.
Sostituendo i valori dei parametri:
$$ \mathbb{E}(X|Z=24) = 24 \cdot \frac{10}{18} = \frac{240}{18} \approx 13.3333333 $$
La risposta è: 13.3333333


% ==========================================
%               ESERCIZIO 8
% ==========================================
\vspace{0.5cm}
\begin{tcolorbox}[colback=lightergray!10, colframe=tocsec, title=\textbf{Esercizio 8}]
Un gioco al Luna Park consiste in un distributore contenente palline numerate da 1 a 5. Una volta inserita una moneta la probabilità di ogni pallina è data dalla seguente tabella:
\begin{center}
\begin{tabular}{c|ccccc}
\toprule
$x$ & 1.00 & 2.00 & 3.00 & 4.00 & 5.00 \\
\midrule
$p$ & 0.12 & 0.29 & 0.34 & 0.02 & 0.23 \\
\bottomrule
\end{tabular}
\end{center}
Indichiamo con $X$ la variabile aleatoria discreta che descrive il risultato ``otteniamo la pallina con numero\dots'', $X(\omega) \in \{1,2,3,4,5\}$.
\end{tcolorbox}

\subsection*{Soluzione}

\subsubsection*{Quesito 1}
1. Calcolare il valore atteso di $X$.

\textbf{Soluzione:}
$$ \mu_1 = \mathbb{E}(X) = \sum_{i=1}^5 i \cdot P(X=i) $$
Svolgendo i calcoli: $1(0.12) + 2(0.29) + 3(0.34) + 4(0.02) + 5(0.23) = 0.12 + 0.58 + 1.02 + 0.08 + 1.15 = 2.95$.
La risposta è: 2.95

\subsubsection*{Quesito 2}
2. Calcolare il terzo momento centrato di $X$.

\textbf{Soluzione:}
Il terzo momento centrato è definito come:
$$ \bar{\mu}_3 = \mathbb{E}((X-\mu_1)^3) = \sum_{i=1}^5 (i-\mu_1)^3 P(X=i) $$
Sostituendo $\mu_1 = 2.95$ e calcolando la sommatoria per ogni $i$, si ottiene:
La risposta è: 0.86625

\subsubsection*{Quesito 3}
3. Calcolare il terzo momento non centrato di $Y=5.25X$

\textbf{Soluzione:}
$$ \mu_3 = \mathbb{E}(Y^3) = \mathbb{E}((bX)^3) = b^3 \mathbb{E}(X^3) $$
Calcoliamo preliminarmente $\mathbb{E}(X^3) = \sum i^3 P(X=i) = 1(0.12) + 8(0.29) + 27(0.34) + 64(0.02) + 125(0.23) = 41.65$.
Quindi: $\mathbb{E}(Y^3) = (5.25)^3 \cdot 41.65 \approx 6026.8851562$.
(Alcuni di voi potrebbero avere una soluzione sbagliata, verranno ricalcolate a fine anno.)
La risposta è: 6026.8851562

\newpage

\subsubsection*{Quesito 4}
4. Probabilità condizionata
Supponiamo ora che siamo due amici, entrambi vorremmo tentare la fortuna con il distributore ma, siccome abbiamo una sola moneta, tiriamo a sorte a chi tocca. La moneta è equilibrata. Io punto su testa. Qual è la probabilità che io ottenga la pallina con il numero 3?

\textbf{Soluzione:}
Se esce croce, la probabilità che io ottenga la pallina con il numero 3 è, ovviamente nulla. Per cui è richiesta la probabilità tramite legge delle probabilità totali:
$$ \begin{aligned} P(\text{io ottengo pallina } 3) &= P(\text{io ottengo pallina } 3 \mid \text{testa})P(\text{testa}) \\ &\quad + P(\text{io ottengo pallina } 3 \mid \text{croce})P(\text{croce}) \end{aligned} $$
Il secondo termine si annulla, quindi:
$$ P(\text{io ottengo pallina } 3) = P(\text{io ottengo pallina } 3 \mid \text{testa})P(\text{testa}) $$
o, alternativamente:
$$ P(\text{io ottengo pallina } 3) = P(\{\text{io ottengo pallina } 3\} \cap \{\text{testa}\}) = \frac{1}{2} P(X=3) = 0.5 \cdot 0.34 = 0.17. $$
La risposta è: 0.17


% ==========================================
%               ESERCIZIO 9
% ==========================================
\vspace{0.5cm}
\begin{tcolorbox}[colback=lightergray!10, colframe=tocsec, title=\textbf{Esercizio 9}]
Sia $X \sim N(\mu=35.63, \sigma^2=9.55)$.
\end{tcolorbox}

\subsubsection*{Quesito 1}
1. Calcolare $Pr(30.96 < X \leq 39.31)$

\textbf{Soluzione:}
Standardizzando i valori estremi con $Z = \frac{X - \mu}{\sigma} = \frac{X - 35.63}{\sqrt{9.55}}$:
$$ P\left(\frac{30.96 - 35.63}{\sqrt{9.55}} < Z \leq \frac{39.31 - 35.63}{\sqrt{9.55}}\right) = \Phi(1.1908) - \Phi(-1.5111) \approx 0.8177661 $$
La risposta è: 0.8177661

\subsubsection*{Quesito 2}
2. Calcolare $Pr(X < 33.98)$

\textbf{Soluzione:}
Standardizzando:
$$ P\left(Z < \frac{33.98 - 35.63}{\sqrt{9.55}}\right) = \Phi(-0.5339) \approx 0.2966959 $$
La risposta è: 0.2966959

\subsubsection*{Quesito 3}
3. Calcolare $Pr(|X| > 39.65)$

\textbf{Soluzione:}
Scomponendo il modulo assoluto: $P(X > 39.65) + P(X < -39.65)$. Essendo la probabilità di un valore così negativo essenzialmente zero per questa gaussiana, si calcola:
$$ 1 - P(X \leq 39.65) = 1 - \Phi\left(\frac{39.65 - 35.63}{\sqrt{9.55}}\right) \approx 0.0966564 $$
La risposta è: 0.0966564


% ==========================================
%               ESERCIZIO 10
% ==========================================
\vspace{0.5cm}
\begin{tcolorbox}[colback=lightergray!10, colframe=tocsec, title=\textbf{Esercizio 10}]
Sappiamo che certi eventi accadono seguendo una distribuzione di Poisson con una media di 5 eventi alla settimana. Sia $X \sim Pois(\lambda)$ con $\lambda=5$.
\end{tcolorbox}

\subsubsection*{Quesito 1}
1. Calcolare la probabilità che più di 7 eventi si verifichino in una settimana.

\textbf{Soluzione:}
Ossia $P(X > 7) = 1 - P(X \leq 7)$. Accumulando le probabilità discrete fino a 7:
$$ 1 - \sum_{k=0}^7 \frac{5^k e^{-5}}{k!} \approx 0.1333717 $$
La risposta è: 0.1333717

\subsubsection*{Quesito 2}
2. Calcolare la probabilità che al più 9 eventi si verifichino in una settimana.

\textbf{Soluzione:}
Ossia $P(X \leq 9)$. Sostituendo e cumulando le densità:
$$ \sum_{k=0}^9 \frac{5^k e^{-5}}{k!} \approx 0.9681719 $$
La risposta è: 0.9681719

\subsubsection*{Quesito 3}
3. Calcolare la probabilità che trascorrano almeno 7 settimane tra due eventi successivi.

\textbf{Soluzione:}
Suggerimento: l’unità di misura di tempo è una settimana; che variabile casuale descrive il tempo trascorso tra eventi successivi?
Abbiamo che le distribuzioni di Poisson, che descrive il numero di eventi un intervallo di tempo, ed Esponenziale sono così legate:
se $X$ è la v.a. che conta il numero di eventi che accadono in una settimana, allora $Y=k \cdot X$ conta il numero di eventi che accadono in $k$ settimane e 
$$ P(Y=y) = \frac{(\lambda k)^y}{y!} e^{-\lambda k} \quad \text{per } \lambda>0; k \geq 0; y=0,1,2,\dots $$
$P(Y=0) = e^{-\lambda k}$ è la probabilità di non osservare alcun evento in un intervallo di $k$ settimane. Pensiamo all’evento come ad un guasto di un elettrodomestico, allora la stessa probabilità può anche essere interpretata come l’elettrodomestico funziona bene almeno per un tempo $k$ (dove l’unità di misura nel nostro caso è la settimana).
Chiamiamo $T$ il tempo di buon funzionamento dell’elettrodomestico, $e^{-\lambda k} = P(T \geq k) = 1 - P(T \leq k)$ da cui $P(T \leq k) = 1 - e^{-\lambda k}$ in cui riconosciamo la funzione di ripartizione di una variabile casuale esponenziale, con parametro $\lambda$.
La probabilità richiesta è $P(T \geq 7)$:
$$ P(T \geq 7) = e^{-5 \cdot 7} = e^{-35} \approx 6.6613381 \times 10^{-16} $$
La risposta è: $6.6613381^{-16}$


% ==========================================
%               ESERCIZIO 11
% ==========================================
\vspace{0.5cm}
\begin{tcolorbox}[colback=lightergray!10, colframe=tocsec, title=\textbf{Esercizio 11}]
Sia data la seguente variabile aleatoria bivariata discreta 
\begin{center}
\begin{tabular}{c|ccc}
\toprule
$X \setminus Y$ & -0.61 & -0.37 & 5.33 \\
\midrule
-1.24 & 9k & 2k & 6k \\
0.89 & 7k & 3k & 5k \\
2.1 & 8k & 1k & 4k \\
\bottomrule
\end{tabular}
\end{center}
dove sulle righe abbiamo la variabile $X$ e sulle colonne la variabile $Y$.
\end{tcolorbox}

\subsubsection*{Quesito 1}
1. Calcolare la costante $k$ che deve essere utilizzata per rendere la tabella una funzione di probabilità congiunta.

\textbf{Soluzione:}
Abbiamo che la somma per righe e colonne nella tabella sopra devono dare 1. La somma dei coefficienti è $9+2+6+7+3+5+8+1+4 = 45$, i.e. $45k = 1$, da cui:
La risposta è: 0.0222222

\subsubsection*{Quesito 2}
2. Calcolare la funzione di probabilità marginale di $X$. Inserire il solo valore di probabilità per $X=2.1$.

\textbf{Soluzione:}
Data la probabilità congiunta $p_{X,Y}(x,y)=P(X=x,Y=y)$ è facile ottenere le marginali.
Indichiamo con $R_Y$ il range o immagine di $Y$, i.e. l’insieme di valori assunti da $Y$, 
$$ p_X(x) = \sum_{y \in R_Y} p_{X,Y}(x,y) $$
Sommando l'ultima riga: $8k + 1k + 4k = 13k = \frac{13}{45} \approx 0.2888889$.
La risposta è: 0.2888889

\subsubsection*{Quesito 3}
3. Calcolare la funzione di probabilità marginale di $Y$. Inserire il solo valore di probabilità per $Y=-0.61$.

\textbf{Soluzione:}
$$ p_Y(y) = \sum_{x \in R_X} p_{X,Y}(x,y) $$
Sommando la prima colonna: $9k + 7k + 8k = 24k = \frac{24}{45} \approx 0.5333333$.
La risposta è: 0.5333333

\vspace{0.3cm}
\textit{Date le funzioni di probabilità marginali, calcolare media e varianza è ormai semplice.}

\subsubsection*{Quesito 4}
4. Calcolare il valore atteso di $X$.

\textbf{Soluzione:}
Sviluppando $\mathbb{E}(X) = \sum x p_X(x) = (-1.24)\frac{17}{45} + (0.89)\frac{15}{45} + (2.1)\frac{13}{45} \approx 0.4348889$.
La risposta è: 0.4348889

\subsubsection*{Quesito 5}
5. Calcolare la varianza di $Y$.

\textbf{Soluzione:}
Calcolando i momenti di $Y$ tramite la marginale e applicando $Var(Y) = \mathbb{E}(Y^2) - \mathbb{E}(Y)^2$.
La risposta è: 7.720736


% ==========================================
%               ESERCIZIO 12
% ==========================================
\vspace{0.5cm}
\begin{tcolorbox}[colback=lightergray!10, colframe=tocsec, title=\textbf{Esercizio 12}]
Sia data la seguente variabile aleatoria bivariata discreta 
\begin{center}
\begin{tabular}{c|cccc}
\toprule
$X \setminus Y$ & -1.19 & 0.67 & 1.38 & 2.2 \\
\midrule
-0.06 & 3k & 6k & 4k & 7k \\
1.19 & 10k & 1k & 8k & 12k \\
2.95 & 2k & 9k & 5k & 11k \\
\bottomrule
\end{tabular}
\end{center}
dove sulle righe abbiamo la variabile $X$ e sulle colonne la variabile $Y$.
\end{tcolorbox}

\subsubsection*{Quesito 1}
1. Calcolare la costante $k$ che deve essere utilizzata per rendere la tabella una funzione di probabilità congiunta.

\textbf{Soluzione:}
Abbiamo che la somma per righe e colonne nella tabella sopra devono dare 1. La somma dei coefficienti è $20 + 31 + 27 = 78$, i.e. $78k = 1$, i.e.
$$ k = \frac{1}{78} \approx 0.0128205 $$
La risposta è: 0.0128205

\subsubsection*{Quesito 2}
2. Calcolare la distribuzione di probabilità marginale di $X$. Determinare $P(X \leq 1.19)$.

\textbf{Soluzione:}
Indichiamo con $R_Y$ il range o immagine di $Y$, i.e. l’insieme di valori assunti da $Y$, la funzione di probabilità di $X$ è data da $p_X(x) = \sum_{y \in R_Y} p_{X,Y}(x,y)$ da cui:
$$ P(X \leq x) = \sum_{k:x_k \leq x} p_X(x_k) $$
Sommando le probabilità marginali per le prime due righe: $\frac{20}{78} + \frac{31}{78} = \frac{51}{78} \approx 0.6538462$.
La risposta è: 0.6538462

\subsubsection*{Quesito 3}
3. Calcolare la distribuzione di probabilità marginale di $Y$. Determinare $P(Y > 1.38)$.

\textbf{Soluzione:}
Simile alla domanda precedente. Sommando le probabilità della colonna in cui $Y=2.2$: $\frac{7k+12k+11k}{78k} = \frac{30}{78} \approx 0.3846154$.
La risposta è: 0.3846154

\subsubsection*{Quesito 4}
4. Calcolare il valore atteso condizionato $\mathbb{E}[X|Y=0.67]$.

\textbf{Soluzione:}
Ora abbiamo bisogno della funzione di probabilità condizionata:
$$ p_{X|Y}(x|y) = \frac{p_{X,Y}(x,y)}{p_Y(y)} \quad \text{per } y \in R_Y : p_Y(y)>0 $$
Ricaviamo prima $p_Y(0.67) = \frac{6+1+9}{78} = \frac{16}{78}$. 
Poi ricalcoliamo le probabilità scalate per $X$: $P(X=-0.06|Y=0.67) = \frac{6}{16}$, $P(X=1.19|Y=0.67) = \frac{1}{16}$, $P(X=2.95|Y=0.67) = \frac{9}{16}$.
Moltiplicando e sommando si ricava il valore atteso condizionato.
La risposta è: 1.71125

\subsubsection*{Quesito 5}
5. Determinare la varianza condizionata $Var[Y|X=1.19]$.

\textbf{Soluzione:}
Simile a sopra. Calcoliamo la distribuzione condizionata di $Y$ data la riga $X=1.19$, ne calcoliamo il momento di ordine 1 e di ordine 2 e determiniamo la varianza.
La risposta è: 2.1214506


% ==========================================
%               ESERCIZIO 13
% ==========================================
\vspace{0.5cm}
\begin{tcolorbox}[colback=lightergray!10, colframe=tocsec, title=\textbf{Esercizio 13}]
Consideriamo un esperimento aleatorio che consiste nel lancio di due dadi non equilibrati. In particolare la funzione di probabilità di massa $p_{X,Y}(x,y)$ – dove $X(\omega) \in \{1,\dots,6\}$ rappresenta il primo dado e $Y(\omega) \in \{1,\dots,6\}$ il secondo dado – è riportata in tabella
\begin{center}
\begin{tabular}{c|cccccc}
\toprule
$X \setminus Y$ & 1 & 2 & 3 & 4 & 5 & 6 \\
\midrule
1 & 0.00120 & 0.00750 & 0.00660 & 0.00900 & 0.00180 & 0.00390 \\
2 & 0.00640 & 0.04000 & 0.03520 & 0.04800 & 0.00960 & 0.02080 \\
3 & 0.00800 & 0.05000 & 0.04400 & 0.06000 & 0.01200 & 0.02600 \\
4 & 0.00680 & 0.04250 & 0.03740 & 0.05100 & 0.01020 & 0.02210 \\
5 & 0.00920 & 0.05750 & 0.05060 & 0.06900 & 0.01380 & 0.02990 \\
6 & 0.00840 & 0.05250 & 0.04620 & 0.06300 & 0.01260 & 0.02730 \\
\bottomrule
\end{tabular}
\end{center}
Come al solito, sulle righe abbiamo la variabile $X$ e sulle colonne la variabile $Y$.
\end{tcolorbox}

\newpage

\subsubsection*{Quesito 1}
1. Determinare la probabilità che sul primo dado esca un numero pari.

\textbf{Soluzione:}
Abbiamo bisogno della funzione di probabilità marginale di $X$, che otteniamo dalla somma per righe:
\begin{center}
\begin{tabular}{rr}
\hline
 $X$ & $p_X(x)$ \\ 
\hline
  1 & 0.03 \\ 
  2 & 0.16 \\ 
  3 & 0.20 \\ 
  4 & 0.17 \\ 
  5 & 0.23 \\ 
  6 & 0.21 \\ 
\hline
\end{tabular}
\end{center}
e la probabilità richiesta si ottiene come $P(X=2)+P(X=4)+P(X=6) = 0.16 + 0.17 + 0.21$.
La risposta è: 0.54

\subsubsection*{Quesito 2}
2. Determinare la probabilità che sul secondo dado non escano 2 o 6.

\textbf{Soluzione:}
Dalla somma sulle colonne otteniamo la marginale di $Y$, dopodiché la probabilità richiesta è data da $1 - P(Y=2) - P(Y=6)$.
$p_Y(y)$: 0.04, 0.25, 0.22, 0.30, 0.06, 0.13.
Pertanto, $1 - 0.25 - 0.13 = 1 - 0.38 = 0.62$.
La risposta è: 0.62

\subsubsection*{Quesito 3}
3. Dopo aver determinato la funzione di probabilità condizionata $p_{X|Y}(x|y=2)$. Determinare $P(X \geq 4|Y=2)$.

\textbf{Soluzione:}
Si osservi che $X$ e $Y$ sono indipendenti (in ogni cella si verifica che $P(X=x, Y=y) = P(X=x)P(Y=y)$), per cui $P(X=x|Y) = P(X=x)$, per cui la probabilità richiesta altri non è che $P(X \geq 4)$.
$$ P(X \geq 4) = 0.17 + 0.23 + 0.21 = 0.61 $$
La risposta è: 0.61

\subsubsection*{Quesito 4}
4. Calcolare il valore atteso condizionato $\mathbb{E}[X|Y=6]$.

\textbf{Soluzione:}
Per la stessa osservazione precedente, il valore atteso condizionato è uguale a $\mathbb{E}(X)$.
$$ \mathbb{E}(X) = 1(0.03) + 2(0.16) + 3(0.20) + 4(0.17) + 5(0.23) + 6(0.21) = 4.04 $$
La risposta è: 4.04

\subsubsection*{Quesito 5}
5. Determinare la varianza condizionata $Var[Y|X=4]$.

\textbf{Soluzione:}
Allo stesso modo, vediamo che la varianza richiesta a causa dell'indipendenza corrisponde alla varianza non condizionata $Var(Y)$.
La risposta è: 1.8896

\newpage

\section*{Probabilità e Statistica - Esercizi 7}
\addcontentsline{toc}{section}{Probabilità e Statistica - Esercizi 7}
\markboth{Esercizi 7}{Esercizi 7}

% ==========================================
%               ESERCIZIO 1
% ==========================================
\begin{tcolorbox}[colback=lightergray!10, colframe=tocsec, title=\textbf{Esercizio 1}]
Per essere una densità di probabilità, una funzione deve avere integrale su $\mathbb{R}$ uguale a 1. In questo esercizio calcoleremo alcune costanti di rinormalizzazione.
\end{tcolorbox}

\subsubsection*{Quesito 1}
Per quale valore di $c$ la funzione definita come $f(x)=c \cdot e^{-8x}$ per $x \geq 0$ e costantemente nulla altrimenti è una densità di probabilità?

\textbf{Soluzione:}
Dobbiamo prendere come costante l'inverso dell'integrale sul supporto:
$$ c = \left(\int_0^{+\infty} e^{-8x} dx\right)^{-1} = 8 $$
dove l’ultima uguaglianza segue dall’integrale elementare della funzione esponenziale $\left[ -\frac{1}{8} e^{-8x} \right]_0^{+\infty} = \frac{1}{8}$.
La risposta è: 8

\subsubsection*{Quesito 2}
Per quale valore di $c$ la funzione definita come $f(x)=c \cdot \dfrac{8}{1+x^2}$ per $x \in \mathbb{R}$ è una densità di probabilità?

\textbf{Soluzione:}
Anche qui abbiamo da fare una semplice integrazione sull'intero asse reale:
$$ c = \left(\int_{-\infty}^{+\infty} \frac{8}{1+x^2} dx\right)^{-1} = \left(8 \int_{-\infty}^{+\infty} \frac{1}{1+x^2} dx\right)^{-1} = (8\pi)^{-1} = \frac{1}{8\pi} $$
dove l’ultima uguaglianza segue dalle proprietà dell’arcotangente (sapendo che $\left[ \arctan(x) \right]_{-\infty}^{+\infty} = \frac{\pi}{2} - \left(-\frac{\pi}{2}\right) = \pi$).
La risposta è: 0.0397887

\subsubsection*{Quesito 3}
Consideriamo ora la funzione $f(x)$ definita per $x \in [0,1]$ da $f(x)=c \cdot x^3(1-x)^b$ e costantemente nulla altrimenti. Sia $b$ un numero intero e non negativo. Calcolare il valore di $c$ in funzione di $b$ in modo che $f(x)$ sia una densità.

\textbf{Soluzione:}
Qui dobbiamo fare un’integrazione per parti ripetuta, sfruttando il fatto che entrambi gli esponenti sono interi non negativi. Chiamiamo
$$ I(a,b) = \int_0^1 x^a(1-x)^b dx $$
Da una prima integrazione per parti abbiamo:
$$ I(a,b) = \frac{a}{b+1} I(a-1, b+1) $$
e possiamo iterare, ottenendo:
$$ I(a,b) = \frac{a(a-1)}{(b+1)(b+2)} I(a-2, b+2) = \frac{a(a-1)(a-2)}{(b+1)(b+2)(b+3)} I(a-3, b+3) $$
Ora sfruttiamo il fatto che $a=3$, quindi abbiamo che:
$$ I(3,b) = \frac{3!}{(b+1)(b+2)(b+3)} I(0,b+3) = \frac{3! \cdot b!}{(b+3)!} \int_0^1 (1-x)^{b+3} dx = \frac{3! \cdot b!}{(b+4)!} $$
In generale questo si potrebbe estendere (mediante la funzione $\Gamma$) anche a numeri complessi di parte reale maggiore di $-1$, tuttavia in questo esercizio era indicato specificamente che $b$ era intero non negativo, quindi possiamo usare il fattoriale.
Per quanto riguarda $c$, allora, essendo l'inverso dell'integrale abbiamo:
$$ c = \frac{(b+4)!}{b! 3!} = \frac{(b+1)(b+2)(b+3)(b+4)}{6} $$
che può essere implementato come funzione (con un controllo sulle altre condizioni richieste) come:
\begin{verbatim}
function(b) {
  ifelse((b%%1==0 & b >= 0), factorial(b+4)/(factorial(3)*factorial(b)), NA)
}
\end{verbatim}


% ==========================================
%               ESERCIZIO 2
% ==========================================
\vspace{0.5cm}
\begin{tcolorbox}[colback=lightergray!10, colframe=tocsec, title=\textbf{Esercizio 2}]
Di una variabile aleatoria $X$ sappiamo che ha la seguente funzione di densità: $f(x)=c \cdot x^{-9}$ per $x>9$ e identicamente nulla altrimenti.
\end{tcolorbox}

\subsection*{Quesiti e soluzioni}
\textit{Siccome ne conosciamo la funzione di densità, sappiamo che la variabile aleatoria $X$ è assolutamente continua.}

\subsubsection*{Quesito 1}
Quanto vale $c$?

\textbf{Soluzione:}
La costante di rinormalizzazione $c$ deve essere tale che $\int_9^{+\infty} c \cdot x^{-9} dx = 1$, ossia:
$$ c \left[ \frac{x^{-8}}{-8} \right]_9^{+\infty} = 1 \implies c \left( \frac{9^{-8}}{8} \right) = 1 $$
Da cui si ricava isolando la costante: $c = 8 \cdot 9^8 = (9-1) \cdot 9^{9-1}$ (si tratta di integrare un monomio).
La risposta corretta è: $3.4437377 \times 10^8$

\newpage

\subsubsection*{Quesito 2}
Qual è il valore atteso di $X$?

\textbf{Soluzione:}
Per definizione, il valore atteso di una variabile aleatoria con densità $f$ a supporto in $(9, +\infty)$ è:
$$ \mathbb{E}[X] = \int_9^{+\infty} x \cdot f(x) dx $$
In questo caso:
$$ \begin{aligned} \mathbb{E}[X] &= \int_9^{+\infty} x \cdot (9-1) \cdot 9^{9-1} \cdot x^{-9} dx \\ &= \int_9^{+\infty} (9-1) \cdot 9^{9-1} \cdot x^{-8} dx \\ &= (9-1) \cdot 9^{9-1} \left[ \frac{x^{-7}}{-7} \right]_9^{+\infty} \\ &= (9-1) \cdot 9^{9-1} \cdot \frac{9^{-7}}{7} = \frac{9(9-1)}{9-2} = \frac{72}{7} \approx 10.2857143 \end{aligned} $$
La risposta corretta è: 10.2857143

\subsubsection*{Quesito 3}
Quanto valgono le seguenti probabilità?
\begin{itemize}
    \item $P(X=26.1)$
    \item $P(X \in [26.1, 27))$
    \item $P(X \in (26.1, 27])$
    \item $P(X \in (-0.92, 26.1))$
\end{itemize}
\textit{Inserire un vettore di risposte del tipo c(valore1, valore2, valore3, valore4).}

\textbf{Soluzione:}
Poiché la $X$ ha densità, la probabilità su un singolo punto è nulla: $P(X=x)=0$ per ogni $x \in \mathbb{R}$.
Per lo stesso motivo (continuità della funzione di ripartizione, $F$, di $X$), gli estremi dell'intervallo non modificano il calcolo dell'area:
$$ P(X \in [26.1, 27)) = P(X \in (26.1, 27]) $$
e questa probabilità si trova come da definizione:
$$ P(X \in (26.1, 27]) = \int_{26.1}^{27} f(x) dx $$
Alternativamente possiamo prima trovare la funzione di ripartizione calcolando l'integrale indefinito:
$$ F(x) = \int_9^x f(t) dt = 1 - \frac{9}{9-1} \cdot x^{-(9-1)} = 1 - \left(\frac{9}{x}\right)^8 \quad \text{per } x>9 $$
e nulla altrove, e poi calcolare $P(X \in (26.1, 27]) = F(27) - F(26.1)$.
Lo stesso vale per l’ultima probabilità richiesta, osservando che $F(-0.92)=0$ e quindi $P(X \in (-0.92, 26.1)) = F(26.1) - F(-0.92) = F(26.1)$.


% ==========================================
%               ESERCIZIO 3
% ==========================================
\vspace{0.5cm}
\begin{tcolorbox}[colback=lightergray!10, colframe=tocsec, title=\textbf{Esercizio 3}]
Abbiamo in un piccolo contenitore, un moscerino della frutta, la cui durata di vita (in settimane) può essere descritta da una variabile aleatoria $X$ con densità di probabilità 
$$ f_X(x)=c \lambda x e^{-\lambda x^2} \quad \text{per } \lambda, x > 0 \text{ con } \lambda=1.42 $$
Il contenitore è monitorato attraverso un sensore, la cui durata (sempre in settimane) può essere descritta da una variabile aleatoria $Y$ di densità 
$$ f_Y(y)=\alpha e^{-\alpha y} \quad \text{per } y \geq 0 \text{ con } \alpha=0.64 $$
\end{tcolorbox}

\subsubsection*{Quesito 1}
Determinare $c$ affinché $f_X$ sia una densità di probabilità.

\textbf{Soluzione:}
Osserviamo che la derivata dell'esponente fornisce proprio i termini in $x$: $\frac{d}{dx}(-\lambda x^2) = -2\lambda x$.
Costruiamo la funzione di ripartizione sfruttando questa derivata:
$$ F_X(t) = \int_0^t c \lambda x e^{-\lambda x^2} dx = -\frac{c}{2} \int_0^t -2\lambda x e^{-\lambda x^2} dx = -\frac{c}{2} \left[ e^{-\lambda x^2} \right]_0^t = \frac{c}{2} - \frac{c}{2} e^{-\lambda t^2} $$
Per $t \to +\infty$ la Funzione di Ripartizione deve essere uguale a 1, quindi affinché $F_X(+\infty) = \frac{c}{2} = 1$, deve essere $c=2$. 
Abbiamo anche già la funzione di ripartizione di $X$, $F_X(x) = 1 - e^{-\lambda x^2}$, che ci verrà comoda a breve.
La risposta è: 2

\subsubsection*{Quesito 2}
Calcolare $P(X \in [0.17, 1.51])$.

\textbf{Soluzione:}
Avendo la funzione di ripartizione, ci basta calcolarla in $0.17$ e sottrarre questo valore a quello di $F_X$ calcolata in $1.51$. (NB: stiamo sfruttando che la funzione di ripartizione è continua.)
$$ P(0.17 \leq X \leq 1.51) = F_X(1.51) - F_X(0.17) \approx 0.9205402 $$
La risposta è: 0.9205402

\subsubsection*{Quesito 3}
Calcolare il momento primo (o momento di ordine 1) di $Y$.

\textbf{Soluzione:}
Il momento primo, noto anche come valore atteso o speranza, si calcola per le variabili aleatorie continue nel modo seguente:
$$ \mathbb{E}(Y) = \int_{-\infty}^{+\infty} y \cdot f_Y(y) dy = \int_0^{+\infty} y \cdot \alpha e^{-\alpha y} dy = \frac{1}{\alpha} $$
usando l’integrazione per parti. Sostituendo $\alpha=0.64$: $\frac{1}{0.64} = 1.5625$.
La risposta è: 1.5625

\newpage

\subsubsection*{Quesito 4}
Calcolare il momento secondo di $Y$.

\textbf{Soluzione:}
Il momento secondo (che non è necessariamente la varianza, ad esempio, non lo è in questo caso) si ottiene in modo analogo al momento primo calcolando $\mathbb{E}(Y^2)$:
$$ \mathbb{E}(Y^2) = \int_{-\infty}^{+\infty} y^2 \cdot f_Y(y) dy = \int_0^{+\infty} y^2 \cdot \alpha e^{-\alpha y} dy = \frac{2}{\alpha^2} $$
usando due volte l’integrazione per parti. Sostituendo $\alpha=0.64$: $\frac{2}{0.64^2} \approx 4.8828125$.
La risposta è: 4.8828125

\subsubsection*{Quesito 5}
Calcolare la probabilità che il moscerino sopravviva al più 1.04 settimane e il sensore duri al più 1.04 settimane.

\textbf{Soluzione:}
Le due variabili aleatorie sono indipendenti, quindi possiamo calcolarne la funzione di ripartizione congiunta semplicemente moltiplicando le funzioni di ripartizione marginali. Abbiamo già $F_X$ e dobbiamo solo ricavare $F_Y$:
$$ F_Y(y) = \int_{-\infty}^y f_Y(\eta) d\eta = \int_0^y \alpha e^{-\alpha \eta} d\eta = 1 - e^{-\alpha y} $$
Ora valutiamo la congiunta nel punto richiesto:
$$ F_{X,Y}(1.04, 1.04) = F_X(1.04) \cdot F_Y(1.04) \approx 0.3814072 $$
La risposta corretta è: 0.3814072


% ==========================================
%               ESERCIZIO 4
% ==========================================
\vspace{0.5cm}
\begin{tcolorbox}[colback=lightergray!10, colframe=tocsec, title=\textbf{Esercizio 4}]
A un posto di controllo della Polizia, passano e vengono fermati in media 5 autoveicoli all’ora. Indichiamo con $X$ la variabile aleatoria che conta il numero di autoveicoli fermati tra le 8 e le 9 del mattino. (Suggerimento: $X$ può essere descritta con uno dei modelli visti a lezione, il parametro della distribuzione è la media indicata.)
\end{tcolorbox}

\vspace{0.3cm}
\textbf{Soluzione preliminare:} Possiamo descrivere questo fenomeno casuale usando la distribuzione di Poisson. Infatti quello che conosciamo è solamente il numero medio di veicoli fermati nell'intervallo temporale. Tra i modelli visti a lezione è quello più adatto. Quindi $X \sim Pois(\lambda)$, con $\lambda=5$.

\subsubsection*{Quesito 1}
Qual è la probabilità che la Polizia non fermi autoveicoli nell’ora indicata?

\textbf{Soluzione:}
Ci stiamo chiedendo quale sia $P(X=0)$ o equivalentemente $\varphi_X(0)$. Possiamo usare la definizione della densità discreta di una variabile aleatoria di Poisson:
$$ \varphi_X(k) = \frac{\lambda^k}{k!} e^{-\lambda} $$
che in questo caso con $k=0$ ci dà $e^{-5}$. Potremmo fare lo stesso calcolo usando la funzione \texttt{dpois} di R, calcolata in 0. (Il risultato che ci aspettiamo è abbastanza piccolo, dal momento che $\lambda=5>0$).
La risposta è: 0.0067379

\subsubsection*{Quesito 2}
Qual è la probabilità che nell’ora indicata ne fermi esattamente 5?

\textbf{Soluzione:}
In questo caso siamo interessati al valore della densità discreta di $X$ valutata in 5, ma la funzione è la stessa vista sopra. Come prima, possiamo usare \texttt{dpois} per calcolare il valore che ci interessa:
$$ \varphi_X(5) = \frac{5^5}{5!} e^{-5} \approx 0.1754674 $$
La risposta è: 0.1754674

\subsubsection*{Quesito 3}
Ora, sia $Y$ una v.a. di Gauss di parametri $\mu, \sigma^2$ rispettivamente la media (o valore atteso) e la varianza della variabile $X$. Qual è la probabilità che $Y$ valga almeno 8 e non più di 13?

\textbf{Soluzione:}
La media e la varianza di un processo di Poisson $X \sim Pois(\lambda)$ coincidono entrambe con $\lambda$, per cui l'approssimazione è $Y \sim N(\lambda,\lambda) = N(5, 5)$.
Qui ci interessa la probabilità che $Y$ cada nell’intervallo $[8,13]$. Ricordando che $Y$, essendo Gaussiana, ha una densità definita su tutto $\mathbb{R}$, vale l'uguaglianza dei bordi $P(Y \in [8,13]) = P(X \in (8,13])$.
Per trovare questa probabilità possiamo calcolare l’integrale della densità Gaussiana, con i parametri dati, fra i due estremi oppure usare la funzione di ripartizione \texttt{pnorm} in 13 e 8 e farne la differenza.
\textit{Attenzione ai parametri della funzione: \texttt{pnorm(q, mean, sd)}, per cui ricordiamoci la radice, i.e. \texttt{sd = sqrt(lambda)}!}
$$ P(8 \leq Y \leq 13) = \Phi\left(\frac{13-5}{\sqrt{5}}\right) - \Phi\left(\frac{8-5}{\sqrt{5}}\right) \approx 0.0896829 $$
La risposta è: 0.0896829


% ==========================================
%               ESERCIZIO 5
% ==========================================
\vspace{0.5cm}
\begin{tcolorbox}[colback=lightergray!10, colframe=tocsec, title=\textbf{Esercizio 5}]
È stato osservato che il tempo trascorso tra il passaggio di due veicoli successivi sotto una videocamera del traffico ha una distribuzione esponenziale di media 11 minuti. Chiamiamo $T$ la corrispondente variabile aleatoria. 
\textit{Memo: la distribuzione esponenziale di media $\mu$ ha densità $f(x) := \frac{1}{\mu} e^{-\frac{x}{\mu}}$, per $x>0$.}
\end{tcolorbox}

\vspace{0.3cm}
\textbf{Soluzione preliminare:} Data una variabile aleatoria esponenziale $T \sim Exp(\lambda)$, sappiamo che la sua media è $\frac{1}{\lambda}$.
Detto questo, è immediato ricavare il parametro di intensità o tasso di decadimento che nel nostro caso è $\lambda = \frac{1}{11} \approx 0.0909091$.

\newpage

\subsubsection*{Quesito 1}
Qual è la probabilità che il tempo trascorso tra il passaggio di due veicoli successivi sia minore di 9.2 minuti?

\textbf{Soluzione:}
La probabilità richiesta si ottiene integrando la densità o calcolando la CDF in quel punto:
$$ P(T < 9.2) = \int_0^{9.2} 0.0909091 e^{-0.0909091 t} dt = 1 - e^{-0.0909091 \cdot 9.2} \approx 0.5667168 $$
Possiamo anche calcolarla computazionalmente in R con la funzione \texttt{pexp}.
La risposta corretta è: 0.5667168

\subsubsection*{Quesito 2}
Qual è l’intervallo di tempo $t$ (in minuti) tale per cui siamo certi al 94\% che il tempo trascorso tra il passaggio di due veicoli sia maggiore di $t$ minuti?

\textbf{Soluzione:}
In questo caso ci viene data la probabilità e dobbiamo determinare il quantile inverso $t$, cioè siamo interessati a 1 meno la funzione quantile calcolata in 0.94. Questo equivale a invertire la funzione di sopravvivenza:
$$ P(T > t) = \int_t^{+\infty} 0.0909091 e^{-0.0909091 t} dt = e^{-0.0909091 t} = 0.94 $$
da cui, prendendo i logaritmi naturali:
$$ t = \frac{-\log(0.94)}{0.0909091} \approx 0.6806294 $$
In alternativa possiamo usare la funzione quantile \texttt{qexp}, ricordandoci di usare il parametro \texttt{lower.tail = FALSE}.
La risposta corretta è: 0.6806294

\subsubsection*{Quesito 3}
Sapendo che sono già trascorsi 9.7 minuti dal passaggio dell’ultimo veicolo, qual è la probabilità che si debba attendere al più altri 4 minuti per il passaggio del veicolo successivo?

\textbf{Soluzione:}
La distribuzione esponenziale, come la geometrica, gode della proprietà di assenza di memoria (o \textit{memoryless property}), quindi il tempo già atteso non influenza in alcun modo l'attesa residua:
$$ P(T < 9.7+4 \mid T > 9.7) = P(T < 4) = 1 - e^{-4 \cdot (1/11)} \approx 0.3048561 $$
(Nota bene: abbiamo enunciato l’assenza di memoria in una forma leggermente diversa, con la prima disuguaglianza di segno opposto, ma questo non fa alcuna differenza rispetto alla formulazione duale del principio).
La risposta corretta è: 0.3048561

\newpage

\subsubsection*{Quesito 4}
Si consideri ora la variabile casuale $U = \sqrt{T}$. Qual è la funzione di densità di $U$?

\textbf{Soluzione:}
Si può usare il teorema sulla trasformazione di variabili aleatorie continue. In questo caso la funzione $g(t) = \sqrt{t}$ è continua e strettamente crescente sulla semiretta $t>0$, con inversa $t = g^{-1}(u) = u^2$. Inoltre derivando l'inversa:
$$ \frac{d}{du} g^{-1}(u) = 2u $$
Alternativamente, se deriviamo la funzione diretta otteniamo $\frac{d}{dt} g(t) = \frac{1}{2} t^{-1/2}$. 
Applichiamo la formula di trasformazione delle densità:
$$ f_U(u) = f_T(u^2) \left| \frac{d}{du}(u^2) \right| = (\lambda e^{-\lambda u^2}) \cdot 2u = 2\lambda u e^{-\lambda u^2} \quad \text{per } u>0 $$
In alternativa si può passare dalla funzione di ripartizione, per $u>0$:
$$ F_U(u) = P(U \leq u) = P(T \leq u^2) = F_T(u^2) = 1 - e^{-\lambda u^2} $$
sfruttando il fatto che $u>0$. A questo punto basta derivare in $u$ per avere la densità trovata sopra. 
Possiamo scrivere questa funzione in codice:
\begin{verbatim}
f_U <- function(u) {
  ifelse(u > 0, (2 * u * lambda) * exp(- lambda * (u ^ 2)), 0)
}
\end{verbatim}

% ==========================================
%               ESERCIZIO 6
% ==========================================
\vspace{0.5cm}
\begin{tcolorbox}[colback=lightergray!10, colframe=tocsec, title=\textbf{Esercizio 6}]
Nella pavimentazione di una stanza vengono adoperati listelli di legno di lunghezza media 16.5 cm. Dai macchinari per la produzione, si sa che le lunghezze di tali listelli seguono una legge normale con deviazione standard uguale a $\sigma=0.13$ cm.
Per ricoprire il buco lasciato dai lavori per inserire un tubo, bisogna creare una striscia larga quanto un listello e lunga 494.0387469 cm, con un opportuno numero $n$ di listelli.
\end{tcolorbox}

\subsubsection*{Quesito 1}
Qual è la probabilità che servano più di $n=30$ listelli per creare una striscia lunga 494.0387469 cm?

\textbf{Soluzione:}
Innanzitutto, sia $X_i$ la v.a. che indica la lunghezza dell’$i$-esimo listello. Allora, $X_i \sim N(\mu,\sigma)$, con $\mu=16.5$ cm e $\sigma=0.13$ cm.
Sia poi $S_n$ la v. a. che indica la lunghezza della striscia formata da $n$ listelli. Allora, 
$$ S_n = \sum_{i=1}^n X_i $$
Poiché le $X_i$ sono indipendenti tra di loro, si ha che $S_n \sim N(n \cdot \mu, \sqrt{n} \cdot \sigma)$ dal momento che la distribuzione Gaussiana è riproducibile (la varianza è $n\sigma^2$, quindi la deviazione standard è $\sqrt{n}\sigma$).
Ora, servono più di 30 listelli per una striscia lunga 494.0387469 cm se la somma dei primi 30 listelli non è sufficiente a coprire la distanza, ovvero se $S_{30} < 494.0387469$. Dobbiamo calcolare quindi la $P(S_{30} < 494.0387469)$.
Ricordiamo che se $X \sim N(\mu,\sigma)$, allora $P(X \leq x) = P\left(Z \leq \frac{x-\mu}{\sigma}\right)$, dove $Z \sim N(0,1)$.
Quindi nel nostro caso abbiamo:
$$ P(S_{30} < 494.0387469) = P\left(Z < \frac{494.0387469 - 30 \cdot 16.5}{0.13 \cdot \sqrt{30}}\right) $$
Possiamo allora usare la funzione di R \texttt{pnorm(x)}. Alternativamente, avremmo potuto guardare le tavole della normale standard.
Notiamo che la funzione \texttt{pnorm} è la cdf di una normale standard, tuttavia se non ci fossimo ricondotti ad una normale standard avremmo potuto comunque usare \texttt{pnorm(x, mean = 30*16.5, sd = sqrt(30)*0.13)}, passando quindi in input i valori della media e della deviazione standard della nostra v.a. gaussiana di partenza $S_{30}$.

Risposta: 0.088508

\subsubsection*{Quesito 2}
Qual è la probabilità che servano meno di $n=30$ listelli per creare una striscia lunga 494.0387469 cm?

\textbf{Soluzione:}
Servono meno di 30 listelli per una striscia lunga 494.0387469 cm se i primi 30 listelli da soli eccedono già la lunghezza richiesta, ovvero se $S_{30} > 494.0387469$. Dobbiamo calcolare quindi la $P(S_{30} > 494.0387469)$.
Basta quindi calcolarsi l'evento complementare \texttt{1-pnorm(x)}, sfruttando il fatto che la Gaussiana è una distribuzione continua.
$$ P(S_{30} > 494.0387469) = 1 - P(S_{30} \leq 494.0387469) \approx 1 - 0.088508 = 0.911492 $$

Risposta: 0.911492

\subsubsection*{Quesito 3}
Quanto dovrebbe essere lunga la striscia affinché 30 listelli siano sufficienti per ricoprirla con una probabilità maggiore o uguale al 64.519\%?

\textbf{Soluzione:}
Dobbiamo trovare quel valore $L$ della lunghezza della striscia per cui valga:
$$ P(S_{30} \geq L) \geq 0.64519 $$
Sia $L^* = \frac{L - n \cdot \mu}{\sqrt{n} \cdot \sigma}$, con $n=30$, $\mu=16.5$ e $\sigma=0.13$, allora:
$$ P(S_n \geq L) = P(Z \geq L^*) = 1 - \Phi(L^*) \geq p $$
dove $p=0.64519$, da cui deduciamo che $\Phi(L^*) \leq 1-p$. Invertendo tramite la funzione quantile (CDF inversa):
$$ L^* = \Phi^{-1}(1-p) \implies \frac{L - n\mu}{\sqrt{n}\sigma} = \Phi^{-1}(1-p) $$
$$ L = \sqrt{n}\sigma \cdot \Phi^{-1}(1-p) + n\mu $$
$\Phi^{-1}(1-p)$ può essere calcolato con \texttt{qnorm(1-p)}.
Sostituendo i valori: $L = \sqrt{30}(0.13) \cdot \Phi^{-1}(1 - 0.64519) + 495 \approx 494.7348604$.

Risposta: 494.7348604


% ==========================================
%               ESERCIZIO 7
% ==========================================
\vspace{0.5cm}
\begin{tcolorbox}[colback=lightergray!10, colframe=tocsec, title=\textbf{Esercizio 7}]
Per una variabile aleatoria bivariata discreta si conosce che le modalità della prima componente ($X$) sono 
\texttt{-9, 2, 7} 
mentre le modalità della seconda componente ($Y$) sono 
\texttt{-7, -5, 0, 9}.

La funzione di probabilità marginale di $X$ è:
\begin{center}
\begin{tabular}{ccc}
\toprule
-9 & 2 & 7 \\
0.35 & 0.35 & 0.30 \\
\bottomrule
\end{tabular}
\end{center}

mentre le distribuzioni condizionate $p_{Y|X}(y|x)$ sono riportate nelle righe della seguente tabella:
\begin{center}
\begin{tabular}{c|cccc}
\toprule
$X \setminus Y$ & -7 & -5 & 0 & 9 \\
\midrule
-9 & 0.24 & 0.20 & 0.28 & 0.28 \\
2 & 0.24 & 0.20 & 0.28 & 0.28 \\
7 & 0.21 & 0.25 & 0.26 & 0.28 \\
\bottomrule
\end{tabular}
\end{center}
dove sulle righe abbiamo la variabile $X$ e sulle colonne la variabile $Y$.
\end{tcolorbox}

\subsection*{Soluzione}
Per le prime due domande si usano la funzione di probabilità marginale di $X$ e le definizioni di speranza matematica (valore atteso) e varianza.

\subsubsection*{Quesito 1}
1. Calcolare $\mathbb{E}(X)$.

\textbf{Soluzione:}
$$ \mathbb{E}(X) = \sum_{x} x \cdot p_X(x) = (-9)(0.35) + (2)(0.35) + (7)(0.30) = -3.15 + 0.70 + 2.10 = -0.35 $$
Risposta: -0.35

\subsubsection*{Quesito 2}
2. Calcolare la $Var(X)$.

\textbf{Soluzione:}
$$ \mathbb{E}(X^2) = \sum_{x} x^2 \cdot p_X(x) = (81)(0.35) + (4)(0.35) + (49)(0.30) = 28.35 + 1.40 + 14.70 = 44.45 $$
$$ Var(X) = \mathbb{E}(X^2) - \mathbb{E}(X)^2 = 44.45 - (-0.35)^2 = 44.45 - 0.1225 = 44.3275 $$
Risposta: 44.3275

\newpage

\subsubsection*{Quesito 3}
3. Calcolare $Var(\mathbb{E}(X|Y))$.

\textbf{Soluzione:}
Abbiamo bisogno della funzione di probabilità condizionata $p_{X|Y}(x|y)$ per calcolare i valori attesi $\mathbb{E}(X|Y)$ e la troviamo usando la definizione di funzione di probabilità condizionata, p.136 note B:
$$ p_{X|Y}(X=x|Y=y) = \frac{p_{X,Y}(x,y)}{p_Y(y)} = \frac{p_{Y|X}(y|x)p_X(x)}{p_Y(y)} \quad \text{per } y \in R_Y (p_Y(y)>0) $$
Le probabilità congiunte (numeratore) $p_{X,Y}(x,y) = p_{Y|X}(y|x)p_X(x)$ sono riportate nella seguente tabella:

\begin{center}
\begin{tabular}{c|cccc}
\toprule
$X \setminus Y$ & -7 & -5 & 0 & 9 \\
\midrule
-9 & 0.0840 & 0.0700 & 0.0980 & 0.0980 \\
2 & 0.0840 & 0.0700 & 0.0980 & 0.0980 \\
7 & 0.0630 & 0.0750 & 0.0780 & 0.0840 \\
\bottomrule
\end{tabular}
\end{center}

da cui è anche facile derivare la marginale di $Y$, $p_Y(y)$, prendendo la somma sulle colonne – \texttt{colSums()} – della matrice:
\begin{itemize}
    \item $p_Y(-7) = 0.0840 + 0.0840 + 0.0630 = 0.231$
    \item $p_Y(-5) = 0.0700 + 0.0700 + 0.0750 = 0.215$
    \item $p_Y(0) = 0.0980 + 0.0980 + 0.0780 = 0.274$
    \item $p_Y(9) = 0.0980 + 0.0980 + 0.0840 = 0.280$
\end{itemize}

I valori attesi condizionati $\mathbb{E}(X|Y=y) = \sum_x x \cdot \frac{p_{X,Y}(x,y)}{p_Y(y)}$ sono:
\begin{itemize}
    \item $\mathbb{E}(X|Y=-7) = -0.6363636$
    \item $\mathbb{E}(X|Y=-5) = 0.1627907$
    \item $\mathbb{E}(X|Y=0) = -0.5109489$
    \item $\mathbb{E}(X|Y=9) = -0.35$
\end{itemize}

Infine, 
$$ Var(\mathbb{E}(X|Y)) = \sum_y (\mathbb{E}(X|Y=y) - \mathbb{E}(X))^2 p_Y(y) $$
dove abbiamo usato il fatto che $\mathbb{E}(\mathbb{E}(X|Y)) = \mathbb{E}(X)$. Sostituendo i valori e le probabilità marginali di $Y$ si ottiene la varianza desiderata.

Risposta: 0.082576

\subsubsection*{Quesito 4}
4. Calcolare $\mathbb{E}(Var(X|Y))$.

\textbf{Soluzione:}
Dalla formula di decomposizione della varianza, Varianza e varianza condizionale, pag.140 note B abbiamo che:
$$ Var(X) = \mathbb{E}(Var(X|Y)) + Var(\mathbb{E}(X|Y)) $$
Poiché $Var(X)$ e $Var(\mathbb{E}(X|Y))$ sono già note, possiamo ottenere il valore richiesto prendendone semplicemente la differenza.
$$ \mathbb{E}(Var(X|Y)) = Var(X) - Var(\mathbb{E}(X|Y)) = 44.3275 - 0.082576 = 44.244924 $$

Risposta: 44.244924


% ==========================================
%               ESERCIZIO 8
% ==========================================
\vspace{0.5cm}
\begin{tcolorbox}[colback=lightergray!10, colframe=tocsec, title=\textbf{Esercizio 8}]
Sono date due variabili casuali discrete: $X$ prende valori in $\{-8, -7.5, -5\}$, mentre $Y(\omega) \in \{-7.5, -6.5, 0.5, 6\}$.
Si conoscono: la funzione di probabilità marginale $p_X(x)=P(X=x)$
\begin{center}
\begin{tabular}{ccc}
\toprule
-8 & -7.5 & -5 \\
0.31 & 0.21 & 0.48 \\
\bottomrule
\end{tabular}
\end{center}
e le probabilità condizionate $p_{Y|X}(y|x)$ (righe della seguente tabella)
\begin{center}
\begin{tabular}{c|cccc}
\toprule
$X \setminus Y$ & -7.5 & -6.5 & 0.5 & 6 \\
\midrule
-8 & 0.34 & 0.21 & 0.19 & 0.26 \\
-7.5 & 0.25 & 0.15 & 0.29 & 0.31 \\
-5 & 0.30 & 0.22 & 0.23 & 0.25 \\
\bottomrule
\end{tabular}
\end{center}
\end{tcolorbox}

\subsection*{Soluzione}

\subsubsection*{Quesito 1}
1. Qual è la probabilità di $(-7.5, 6)$?

\textbf{Soluzione:}
Definizione di probabilità condizionata:
$$ p_{X|Y}(X=x|Y=y) = P(\{X=x\}|\{Y=y\}) = \frac{p_{X,Y}(x,y)}{p_Y(y)} \quad \text{per } y \in R_Y(p_Y(y)>0) $$
da cui, analogamente per l'altra direzione:
$$ p_{X,Y}(x,y) = p_{Y|X}(Y=y|X=x) \cdot p_X(x) $$
Sostituendo le coordinate $(-7.5, 6)$:
$$ P(X=-7.5, Y=6) = p_{Y|X}(6 \mid -7.5) \cdot p_X(-7.5) = 0.31 \cdot 0.21 = 0.0651 $$

Risposta: 0.0651

\subsubsection*{Quesito 2}
2. Determinare $Cov(X,Y)$.

\textbf{Soluzione:}
E’ possibile calcolare la covarianza in diversi modi.
Seguendo la definizione di covarianza (prima centrare le variabili attorno alla rispettiva media e poi valore atteso del prodotto), $Cov(X,Y) = \mathbb{E}[(X-\mathbb{E}(X))(Y-\mathbb{E}(Y))]$.

$\mathbb{E}(X)$ si ottiene immediatamente, poiché la marginale di $X$ è nota:
$$ \mathbb{E}(X) = (-8)(0.31) + (-7.5)(0.21) + (-5)(0.48) = -6.455 $$
Per calcolare $\mathbb{E}(Y)$ possiamo usare la regola dell’aspettazione totale $\mathbb{E}(Y) = \mathbb{E}(\mathbb{E}(Y|X))$ o ricavare la marginale di $Y$ dalla funzione di probabilità congiunta ricavata nel punto precedente. $\mathbb{E}(Y) = -1.86925$.
Dopodiché centriamo le variabili $X$ e $Y$ attorno alle rispettive medie, chiamiamo $\tilde{X}$ e $\tilde{Y}$ le rispettive v.a. centrate, e calcoliamo il valore atteso del prodotto $\mathbb{E}(\tilde{X}\tilde{Y})$.

Usando la formula che segue la definizione (valore atteso del prodotto meno prodotto dei valori attesi):
Ottenere $\mathbb{E}(XY)$ usando la Proposizione 4, pag. 67 note A:
$$ \mathbb{E}(XY) = \sum_i \sum_j x_i y_j p_{X,Y}(x_i, y_j) $$
Ottenere $\mathbb{E}(X), \mathbb{E}(Y)$ come sopra e poi:
$$ Cov(X,Y) = \mathbb{E}(XY) - \mathbb{E}(X)\mathbb{E}(Y) \approx -0.1743338 $$

Risposta: -0.1743338

\subsubsection*{Quesito 3}
3. $X$ e $Y$ sono non correlate? Rispondere TRUE o FALSE.

\textbf{Soluzione:}
Due variabili casuali $X$ e $Y$ sono dette non correlate (o incorrelate) se la loro covarianza è nulla. Essendo nel nostro caso la covarianza diversa da zero, le variabili risultano correlate.
Si osservi che se non sapessimo niente delle distribuzioni di $X$ e $Y$ e ci venisse detto solo che $Cov(X,Y)=0$ non potremmo rispondere alla domanda \textit{sono $X$ e $Y$ indipendenti?}
Mentre se ci viene detto che $Cov(X,Y) \neq 0$ abbiamo che dall’implicazione 
$$ X,Y \text{ indipendenti} \implies Cov(X,Y)=0 $$
o equivalentemente 
$$ Cov(X,Y) \neq 0 \implies X,Y \text{ dipendenti} $$
possiamo concludere che $X$ e $Y$ sono dipendenti.

Risposta: FALSE

\newpage

\subsubsection*{Quesito 4}
4. Determinare $\rho(X,Y)$.

\textbf{Soluzione:}
La correlazione di $X$ e $Y$ è data dalla covarianza standardizzata con le deviazioni standard marginali:
$$ \rho(X,Y) = \frac{Cov(X,Y)}{\sqrt{Var(X) \cdot Var(Y)}} $$
Calcolando le varianze $\mathbb{E}(X^2) - \mathbb{E}(X)^2$ e $\mathbb{E}(Y^2) - \mathbb{E}(Y)^2$ e sostituendo, otteniamo il coefficiente.

Risposta: -0.0219907


% ==========================================
%               ESERCIZIO 9
% ==========================================
\vspace{0.5cm}
\begin{tcolorbox}[colback=lightergray!10, colframe=tocsec, title=\textbf{Esercizio 9}]
\textbf{Quesito 1.} Un viaggiatore che ha perso la sua mappa si trova a un bivio da cui partono tre strade:
\begin{itemize}
    \item la prima strada riporta al punto in cui si trova ora, dopo un’ora.
    \item anche la seconda strada riporta al punto in cui il viaggiatore si trova ora, dopo 2 ore.
    \item la terza strada porta alla città successiva, dopo 3 ore di cammino.
\end{itemize}
Ogni volta che il viaggiatore si trova al bivio sceglie una strada con la stessa probabilità (il viaggiatore non ha memoria).
\end{tcolorbox}

\subsubsection*{Soluzione Quesito 1}
Sia $T$ il tempo necessario per raggiungere la città e $X_i =$ ``il viaggiatore prende la strada $i$'', dunque $X_i \in \{1,2,3\}$.
Determinare il tempo medio che il viaggiatore impiega per raggiungere la città.
Usare il fatto che $\mathbb{E}(T|X_i) = \mathbb{E}(T) + k_i$ ore e la speranza matematica della speranza matematica condizionale.

Bisognava osservare che $\mathbb{E}(T|X_i) = \mathbb{E}(T) + k_i$ vale solo per la strada 1 e 2, poiché significa che aggiungo ore al tempo medio e ricomincio il processo da capo, mentre se il viaggiatore prende l’ultima strada, quella che porta effettivamente alla città e interrompe l'esperimento, $\mathbb{E}(T|X_3) = 3$.
Dopodiché per la legge dell'aspettazione totale:
$$ \begin{aligned} \mathbb{E}(T) &= \mathbb{E}(\mathbb{E}(T|X_i)) \\ &= \mathbb{E}(T|X_1)P(X_1) + \mathbb{E}(T|X_2)P(X_2) + \mathbb{E}(T|X_3)P(X_3) \\ &= \frac{1}{3} (\mathbb{E}(T) + 1 + \mathbb{E}(T) + 2 + 3) \\ &= \frac{1}{3} (2\mathbb{E}(T) + 6) \end{aligned} $$
Risolviamo l'equazione per la variabile $\mathbb{E}(T)$:
$$ \mathbb{E}(T) = \frac{2}{3}\mathbb{E}(T) + 2 \implies \frac{1}{3}\mathbb{E}(T) = 2 \implies \mathbb{E}(T) = 6 $$

Risposta: 6

\vspace{0.5cm}
\begin{tcolorbox}[colback=lightergray!10, colframe=tocsec]
\textbf{Quesito 2.} Sia $Y$ una v.a. discreta che prende i valori $\{1,2,3\}$ con le seguenti probabilità $0.29, 0.24, 0.47$.
Sono date inoltre le probabilità condizionate $p_{X|Y}(x|y)$ (righe della seguente tabella):
\begin{center}
\begin{tabular}{c|cccc}
\toprule
$Y \setminus X$ & -9 & 6 & 8 & 8.5 \\
\midrule
1 & 0.22 & 0.23 & 0.31 & 0.24 \\
2 & 0.19 & 0.24 & 0.31 & 0.26 \\
3 & 0.26 & 0.25 & 0.28 & 0.21 \\
\bottomrule
\end{tabular}
\end{center}
\end{tcolorbox}

\subsubsection*{Soluzione 2.1: Determinare $\mathbb{E}(X)$}
$$ \mathbb{E}(X) = \mathbb{E}(\mathbb{E}(X|Y)) = \sum_{y=1}^3 \left[ \sum_{x \in R_X} x \cdot p_{X|Y}(x|y) \right] p_Y(y) $$
Per cui prima determiniamo $\mathbb{E}(X|Y)$ per ogni riga:
\begin{itemize}
    \item $\mathbb{E}(X|Y=1) = -9(0.22) + 6(0.23) + 8(0.31) + 8.5(0.24) = 3.92$
    \item $\mathbb{E}(X|Y=2) = -9(0.19) + 6(0.24) + 8(0.31) + 8.5(0.26) = 4.42$
    \item $\mathbb{E}(X|Y=3) = -9(0.26) + 6(0.25) + 8(0.28) + 8.5(0.21) = 3.185$
\end{itemize}
Sostituendo e pesando con $p_Y(y)$: $\mathbb{E}(X) = 3.92(0.29) + 4.42(0.24) + 3.185(0.47) = 3.69455$.

Risposta: 3.69455

\subsubsection*{Soluzione 2.2: Determinare $Var(X)$}
Vediamo in Varianza e varianza condizionale, pag.140 note B che è possibile ottenere la varianza di $X$ come:
$$ Var(X) = \mathbb{E}(Var(X|Y)) + Var(\mathbb{E}(X|Y)) $$
Questa è anche nota come ``legge della varianza totale''. Iniziamo da $Var(\mathbb{E}(X|Y))$, i valori attesi $\mathbb{E}(X|Y)$ sono noti dal punto precedente, per calcolare la loro varianza abbiamo bisogno della media $\mathbb{E}(\mathbb{E}(X|Y)) = \mathbb{E}(X)$, anch’essa nota, dopodiché:
$$ Var(\mathbb{E}(X|Y)) = \sum_y (\mathbb{E}(X|Y=y) - \mathbb{E}(X))^2 p_Y(y) = 0.263078 $$
Per quanto riguarda $\mathbb{E}(Var(X|Y))$ abbiamo bisogno di un po’ più di lavoro, ricavando per ogni riga $Var(X|Y=y) = \mathbb{E}(X^2|Y=y) - \mathbb{E}(X|Y=y)^2$:
\begin{itemize}
    \item $Var(X|Y=1) = \sum_x (x - \mathbb{E}(X|Y=1))^2 p_{X|Y}(x|y=1) = 47.9136$
    \item $Var(X|Y=2) = \sum_x (x - \mathbb{E}(X|Y=2))^2 p_{X|Y}(x|y=2) = 43.1186$
    \item $Var(X|Y=3) = \sum_x (x - \mathbb{E}(X|Y=3))^2 p_{X|Y}(x|y=3) = 53.008275$
\end{itemize}
dopodiché 
$$ \mathbb{E}(Var(X|Y)) = \sum_y Var(X|Y=y)p_Y(y) = 49.1572972 $$
Sommando i due risultati parziali otteniamo la risposta desiderata $49.1572972 + 0.263078 = 49.4203753$.

Risposta: 49.4203753


% ==========================================
%               ESERCIZIO 10
% ==========================================
\vspace{0.5cm}
\begin{tcolorbox}[colback=lightergray!10, colframe=tocsec, title=\textbf{Esercizio 10}]
Sia $X \sim P(\lambda)$, sapendo che $P(X=2) = P(X=3)$.
\end{tcolorbox}

\subsubsection*{Quesito 1}
1. Calcolare il valore di $\lambda$.

\textbf{Soluzione:}
Scriviamo la funzione per generici $k$ e $k+1$:
$$ P(X=k) = \frac{\lambda^k e^{-\lambda}}{k!} \quad \text{e} \quad P(X=k+1) = \frac{\lambda^{k+1} e^{-\lambda}}{(k+1)!} $$
Ora, $P(X=k)=P(X=k+1)$ con $k=2$ equivale a:
$$ \frac{\lambda^k e^{-\lambda}}{k!} = \frac{\lambda^{k+1} e^{-\lambda}}{(k+1)!} $$
Semplificando $e^{-\lambda}$ e raccogliendo i termini:
$$ \frac{1}{k!} = \frac{\lambda}{(k+1) \cdot k!} \implies 1 = \frac{\lambda}{k+1} $$
da cui $\lambda = k+1$. Sostituendo $k=2$, troviamo che $\lambda = 3$.

Risposta: 3

\vspace{0.3cm}
\textit{Sia ora $Y$ un’altra variabile aleatoria di Poisson, con media 5 e indipendente da $X$. E’ noto che la somma di v.a. indipendenti di Poisson, $Z=X+Y$ è anch’essa distribuita con legge di Poisson con parametro $\lambda$ dato dalla somma dei parametri di $X$ e $Y$.
Dopo aver derivato la funzione di probabilità condizionata $X|Z=n$ (scrivetela con $n$, $\lambda_X$ e $\lambda_Y$ generici, tornerà utile nella domanda 3) rispondere alla seguente domanda:}

\subsubsection*{Quesito 2}
2. Qual è la probabilità $P(X=5|Z=11)$.

\textbf{Soluzione:}
La media di $Y \sim Pois(\lambda_Y)$ è $\lambda_Y$ da cui capiamo che $Y \sim Pois(\lambda_Y=5)$.
Come suggerito dal testo e per la proprietà di riproducibilità, sappiamo anche che $Z \sim Pois(8)$.
$$ \begin{aligned} P(X=k|Z=n) &= \frac{P(X=k, Z=n)}{P(Z=n)} \\ &= \frac{P(X=k, X+Y=n)}{P(Z=n)} \\ &= \frac{P(X=k, Y=n-k)}{P(Z=n)} \\ &= \frac{P(X=k)P(Y=n-k)}{P(Z=n)} \end{aligned} $$
dove abbiamo usato la definizione di probabilità condizionata (nella prima eguaglianza), la definizione di $Z=X+Y$ nella seconda, l’ovvia osservazione che l’intersezione di eventi $(X=k, Y=n-k)$ si può avere solo se $Y=n-k$ e infine l’indipendenza di $X$ e $Y$.
Sviluppando algebricamente in Binomiale (come mostrato nel quesito 3) si ottiene che la variabile si comporta come $Bin(n=11, p=\frac{3}{8})$.
$$ P(X=5|Z=11) = \binom{11}{5} \left(\frac{3}{8}\right)^5 \left(\frac{5}{8}\right)^6 \approx 0.2042107 $$

Risposta: 0.2042107

\subsubsection*{Quesito 3}
3. Determinare il valore atteso condizionato $\mathbb{E}(X|Z=k[2])$.
\textit{Suggerimento: La $p_{X|Z}(x|z=n)$ risulterà essere una funzione di densità nota, se i passaggi sono giusti!}

\textbf{Soluzione:}
Qui torna utile sviluppare l’espressione derivata sopra. Per semplicità scriviamo $\lambda_X$ e $\lambda_Y$ per i parametri di $X$ e $Y$ rispettivamente.
$$ \begin{aligned} P(X=k|Z=n) &= \frac{P(X=k)P(Y=n-k)}{P(Z=n)} \\ &= \frac{ \frac{\lambda_X^k e^{-\lambda_X}}{k!} \frac{\lambda_Y^{n-k} e^{-\lambda_Y}}{(n-k)!} }{ \frac{(\lambda_X+\lambda_Y)^n e^{-(\lambda_X+\lambda_Y)}}{n!} } \\ &= \frac{\lambda_X^k \lambda_Y^{n-k}}{(\lambda_X+\lambda_Y)^{k+n-k}} \cdot \frac{n!}{(n-k)!k!} \\ &= \binom{n}{k} \left(\frac{\lambda_X}{\lambda_X+\lambda_Y}\right)^k \left(\frac{\lambda_Y}{\lambda_X+\lambda_Y}\right)^{n-k} \end{aligned} $$
poichè $\frac{1}{(n-k)!k!} \cdot \frac{n!}{1} = \frac{n!}{(n-k)!k!} = \binom{n}{k}$, e $e^{-\lambda_X} e^{-\lambda_Y} = e^{-(\lambda_X+\lambda_Y)}$.

Osserviamo che la funzione di probabilità condizionata è una legge binomiale con parametri $n$ e $p = \frac{\lambda_X}{\lambda_X+\lambda_Y}$ per cui sappiamo immediatamente che il valore atteso di $X|Z=n$ è $np = n\frac{\lambda_X}{\lambda_X+\lambda_Y}$.
Sostituendo con i valori del problema ($n=11$, $\lambda_X=3$, $\lambda_Y=5$) si ottiene:
$$ \mathbb{E}(X|Z=11) = 11 \cdot \frac{3}{8} = 4.125 $$

Risposta: 4.125

\section*{Probabilità e Statistica - Esercizi 8}
\addcontentsline{toc}{section}{Probabilità e Statistica - Esercizi 8}
\markboth{Esercizi 8}{Esercizi 8}

% ==========================================
%               ESERCIZIO 1
% ==========================================
\begin{tcolorbox}[colback=lightergray!10, colframe=tocsec, title=\textbf{Esercizio 1}]
Per una variabile aleatoria bivariata discreta si conosce che le modalità della prima componente ($X$) sono 
\texttt{-8, -7, -3} 
mentre le modalità della seconda componente ($Y$) sono 
\texttt{-9, -3, 0, 1}

La funzione di probabilità marginale di $X$ è:
\begin{center}
\begin{tabular}{ccc}
\toprule
-8 & -7 & -3 \\
0.36 & 0.34 & 0.30 \\
\bottomrule
\end{tabular}
\end{center}

mentre le distribuzioni condizionate $p_{Y|X}(y|x)$ sono riportate nelle righe della seguente tabella:
\begin{center}
\begin{tabular}{c|cccc}
\toprule
$X \setminus Y$ & -9 & -3 & 0 & 1 \\
\midrule
-8 & 0.26 & 0.30 & 0.26 & 0.18 \\
-7 & 0.27 & 0.28 & 0.16 & 0.29 \\
-3 & 0.28 & 0.23 & 0.23 & 0.26 \\
\bottomrule
\end{tabular}
\end{center}
dove sulle righe abbiamo la variabile $X$ e sulle colonne la variabile $Y$.
\end{tcolorbox}

\subsubsection*{Quesito 1}
1. Calcolare $\mathbb{E}(X)$

\textbf{Soluzione:}
Applicando la definizione di valore atteso per la distribuzione marginale:
$$ \mathbb{E}(X) = \sum_x x \cdot p_X(x) = (-8)(0.36) + (-7)(0.34) + (-3)(0.30) = -6.16 $$
La risposta è: -6.16

\subsubsection*{Quesito 2}
2. Calcolare la $Var(X)$

\textbf{Soluzione:}
Calcoliamo il momento non centrato di ordine 2:
$$ \mathbb{E}(X^2) = \sum_x x^2 \cdot p_X(x) = (64)(0.36) + (49)(0.34) + (9)(0.30) = 23.04 + 16.66 + 2.7 = 42.4 $$
Quindi la varianza sarà:
$$ Var(X) = \mathbb{E}(X^2) - \mathbb{E}(X)^2 = 42.4 - (-6.16)^2 = 42.4 - 37.9456 = 4.4544 $$
La risposta è: 4.4544

\newpage

\subsubsection*{Quesito 3}
3. Calcolare $Var(\mathbb{E}(X|Y))$

\textbf{Soluzione:}
Abbiamo bisogno della funzione di probabilità condizionata $p_{X|Y}(x|y)$ per calcolare i valori attesi $\mathbb{E}(X|Y)$ e la troviamo usando la definizione di funzione di probabilità condizionata, p.136 note B
$$ p_{X|Y}(X=x|Y=y) = \frac{p_{X,Y}(x,y)}{p_Y(y)} \quad \text{per } y \in R_Y(p_Y(y)>0) $$
$$ = \frac{p_{Y|X}(y|x)p_X(x)}{p_Y(y)} $$
Le probabilità congiunte (numeratore) sono riportate nella seguente tabella:
\begin{center}
\begin{tabular}{c|cccc}
\toprule
$X \setminus Y$ & -9 & -3 & 0 & 1 \\
\midrule
-8 & 0.0936 & 0.1080 & 0.0936 & 0.0648 \\
-7 & 0.0918 & 0.0952 & 0.0544 & 0.0986 \\
-3 & 0.0840 & 0.0690 & 0.0690 & 0.0780 \\
\bottomrule
\end{tabular}
\end{center}
da cui è anche facile derivare la marginale di $Y$, $p_Y(y)$, prendendo la somma sulle colonne – \texttt{colSums()} – della matrice. I valori marginali $p_Y(y)$ sono rispettivamente $0.2694, 0.2722, 0.2170, 0.2414$.

I valori attesi condizionati risultano:
\begin{itemize}
    \item $\mathbb{E}(X|Y=-9) = -6.1002227$
    \item $\mathbb{E}(X|Y=-3) = -6.3828068$
    \item $\mathbb{E}(X|Y=0) = -6.159447$
    \item $\mathbb{E}(X|Y=1) = -5.9759735$
\end{itemize}
Infine, 
$$ Var(\mathbb{E}(X|Y)) = \sum_y (\mathbb{E}(X|Y=y) - \mathbb{E}(X))^2 p_Y(y) $$
dove abbiamo usato il fatto che $\mathbb{E}(\mathbb{E}(X|Y)) = \mathbb{E}(X)$.
La risposta è: 0.0226507

\subsubsection*{Quesito 4}
4. Calcolare $\mathbb{E}(Var(X|Y))$.

\textbf{Soluzione:}
Dalla formula di decomposizione della varianza abbiamo che 
$$ Var(X) = \mathbb{E}(Var(X|Y)) + Var(\mathbb{E}(X|Y)) $$
Poiché $Var(X)$ e $Var(\mathbb{E}(X|Y))$ sono già note, possiamo ottenere il valore richiesto prendendone la differenza.
$$ \mathbb{E}(Var(X|Y)) = 4.4544 - 0.0226507 = 4.4317493 $$
La risposta è: 4.4317493


% ==========================================
%               ESERCIZIO 2
% ==========================================
\vspace{0.5cm}
\begin{tcolorbox}[colback=lightergray!10, colframe=tocsec, title=\textbf{Esercizio 2}]
Sono date due variabili casuali discrete: $X$ prende valori in $\{-5, 0.5, 6.5\}$, mentre $Y(\omega) \in \{-4.5, -1, 6, 9\}$.
Si conoscono: la funzione di probabilità marginale $p_X(x)=P(X=x)$
\begin{center}
\begin{tabular}{ccc}
\toprule
-5 & 0.5 & 6.5 \\
0.26 & 0.31 & 0.43 \\
\bottomrule
\end{tabular}
\end{center}
e le probabilità condizionate $p_{Y|X}(y|x)$ (righe della seguente tabella)
\begin{center}
\begin{tabular}{c|cccc}
\toprule
$X \setminus Y$ & -4.5 & -1 & 6 & 9 \\
\midrule
-5 & 0.32 & 0.23 & 0.19 & 0.26 \\
0.5 & 0.30 & 0.20 & 0.25 & 0.25 \\
6.5 & 0.23 & 0.27 & 0.25 & 0.25 \\
\bottomrule
\end{tabular}
\end{center}
\end{tcolorbox}

\subsubsection*{Quesito 1}
1. Qual è la probabilità dell’evento congiunto $(X=-5, Y=-4.5)$?

\textbf{Soluzione:}
Definizione di probabilità condizionata 
$$ p_{X|Y}(X=x|Y=y) = P(\{X=x\}|\{Y=y\}) = \frac{p_{X,Y}(x,y)}{p_Y(y)} \quad \text{per } y \in R_Y(p_Y(y)>0) $$
da cui
$$ p_{X,Y}(x,y) = p_{Y|X}(Y=y|X=x) \cdot p_X(x) $$
Sostituendo i valori della tabella: $0.32 \cdot 0.26 = 0.0832$.
La risposta è: 0.0832

\subsubsection*{Quesito 2}
2. Determinare $Cov(X,Y)$

\textbf{Soluzione:}
E’ possibile calcolare la covarianza in diversi modi.

\textbf{1.} Seguendo la definizione di covarianza (prima centrare le variabili attorno alla rispettiva media e poi valore atteso del prodotto), 
$$ Cov(X,Y) = \mathbb{E}[(X-\mathbb{E}(X))(Y-\mathbb{E}(Y))]. $$
$\mathbb{E}(X)$ si ottiene immediatamente, poiché la marginale di $X$ è nota, $\mathbb{E}(X)=1.65$. Per calcolare $\mathbb{E}(Y)$ possiamo usare la regola dell’aspettazione totale $\mathbb{E}(Y) = \mathbb{E}(\mathbb{E}(Y|X))$ o ricavare la marginale di $Y$ dalla funzione di probabilità congiunta ricavata nel punto precedente. $\mathbb{E}(Y)=2.20395$.
Dopodiché centriamo le variabili $X$ e $Y$ attorno alle rispettive medie, chiamiamo $\tilde{X}$ e $\tilde{Y}$ le rispettive v.a. centrate, e calcoliamo il valore atteso del prodotto $\mathbb{E}(\tilde{X}\tilde{Y})$.

\textbf{2.} Usando la formula valore atteso del prodotto meno prodotto dei valori attesi:
Ottenere $\mathbb{E}(XY)$ 
$$ \mathbb{E}(XY) = \sum_i \sum_j x_i y_j p_{X,Y}(x_i, y_j). $$
Ottenere $\mathbb{E}(X)$, $\mathbb{E}(Y)$ come sopra e poi 
$$ Cov(X,Y) = \mathbb{E}(XY) - \mathbb{E}(X)\mathbb{E}(Y). $$
La risposta è: 1.1852575

\subsubsection*{Quesito 3}
3. $X$ e $Y$ sono non correlate? Rispondere TRUE o FALSE.

\textbf{Soluzione:}
Due variabili casuali $X$ e $Y$ sono dette non correlate (o incorrelate) se la loro covarianza è nulla. Nel nostro caso $Cov(X,Y) \neq 0$.

La risposta è: FALSE

Si osservi che se non sapessimo niente delle distribuzioni di $X$ e $Y$ e ci venisse detto solo che $Cov(X,Y)=0$ non potremmo rispondere alla domanda \textit{sono $X$ e $Y$ indipendenti?}
Mentre se ci viene detto che $Cov(X,Y) \neq 0$ abbiamo che dall’implicazione 
$$ X,Y \text{ indipendenti} \implies Cov(X,Y)=0 $$
o equivalentemente 
$$ Cov(X,Y) \neq 0 \implies X,Y \text{ dipendenti} $$
possiamo concludere che $X$ e $Y$ sono dipendenti.

\subsubsection*{Quesito 4}
4. Determinare $\rho(X,Y)$.

\textbf{Soluzione:}
La correlazione di $X$ e $Y$ è data da
$$ \rho(X,Y) = \frac{Cov(X,Y)}{\sqrt{Var(X) \cdot Var(Y)}} $$
La risposta è: 0.046228


% ==========================================
%               ESERCIZIO 3
% ==========================================
\vspace{0.5cm}
\begin{tcolorbox}[colback=lightergray!10, colframe=tocsec, title=\textbf{Esercizio 3}]
\textbf{Testo 1.} Un viaggiatore che ha perso la sua mappa si trova a un bivio da cui partono tre strade:
\begin{itemize}
    \item la prima strada riporta al punto in cui si trova ora, dopo un’ora.
    \item anche la seconda strada riporta al punto in cui il viaggiatore si trova ora, dopo 3 ore.
    \item la terza strada porta alla città successiva, dopo 2 ore di cammino.
\end{itemize}
Ogni volta che il viaggiatore si trova al bivio sceglie una strada con la stessa probabilità (il viaggiatore non ha memoria)
Sia $T$ il tempo necessario per raggiungere la città e $X_i =$ ``il viaggiatore prende la strada $i$'', dunque $X_i \in \{1,2,3\}$.
\end{tcolorbox}

\newpage

\subsubsection*{Soluzione Testo 1}
Determinare il tempo medio che il viaggiatore impiega per raggiungere la città.
Usare il fatto che $\mathbb{E}(T|X_i) = \mathbb{E}(T) + k_i$ ore è la \textit{speranza matematica della speranza matematica condizionale}.

Bisognava osservare che $\mathbb{E}(T|X_i) = \mathbb{E}(T) + k_i$ vale solo per la strada 1 e 2, poiché significa che aggiungo ore al tempo medio, mentre se il viaggiatore prende l’ultima strada, quella che porta effettivamente alla città, $\mathbb{E}(T|X_3) = 2$.
Dopodiché per la legge delle probabilità totali:
$$ \begin{aligned} \mathbb{E}(T) &= \mathbb{E}(\mathbb{E}(T|X_i)) \\ &= \mathbb{E}(T|X_1)P(X_1) + \mathbb{E}(T|X_2)P(X_2) + \mathbb{E}(T|X_3)P(X_3) \\ &= \frac{1}{3} (\mathbb{E}(T) + 1 + \mathbb{E}(T) + 3 + 2) \\ &= \frac{1}{3} (2\mathbb{E}(T) + 6) \end{aligned} $$
Da cui si ottiene facilmente:
$$ \mathbb{E}(T) - \frac{2}{3}\mathbb{E}(T) = 2 \implies \frac{1}{3}\mathbb{E}(T) = 2 \implies \mathbb{E}(T) = 6 $$
La risposta è: 6

\vspace{0.5cm}
\begin{tcolorbox}[colback=lightergray!10, colframe=tocsec]
\textbf{Testo 2.} Sia $Y$ una v.a. discreta che prende i valori $\{1,2,3\}$ con le seguenti probabilità $0.32, 0.44, 0.24$.
Sono date inoltre le probabilità condizionate $p_{X|Y}(x|y)$ (righe della seguente tabella):
\begin{center}
\begin{tabular}{c|cccc}
\toprule
$Y \setminus X$ & -8 & -6 & 3 & 8.5 \\
\midrule
1 & 0.41 & 0.22 & 0.23 & 0.14 \\
2 & 0.25 & 0.27 & 0.28 & 0.20 \\
3 & 0.23 & 0.24 & 0.31 & 0.22 \\
\bottomrule
\end{tabular}
\end{center}
\end{tcolorbox}

\subsubsection*{Soluzione 2.1}
2. Determinare $\mathbb{E}(X)$.

\textbf{Soluzione:}
$$ \mathbb{E}(X) = \mathbb{E}(\mathbb{E}(X|Y)) = \sum_{y=1}^3 \left[ \sum_{x \in R_X} x \cdot p_{X|Y}(x|y) \right] p_Y(y) $$
Per cui prima determiniamo $\mathbb{E}(X|Y)$:
\begin{itemize}
    \item $\mathbb{E}(X|Y=1) = -2.72$
    \item $\mathbb{E}(X|Y=2) = -1.08$
    \item $\mathbb{E}(X|Y=3) = -0.48$
\end{itemize}
Pesando i valori condizionati con le probabilità marginali di $Y$ ($0.32, 0.44, 0.24$):
La risposta è: -1.4608

\subsubsection*{Soluzione 2.2}
3. Determinare $Var(X)$.

\textbf{Soluzione:}
Vediamo in Varianza e varianza condizionale, pag.140 note B che è possibile ottenere la varianza di $X$ come 
$$ Var(X) = \mathbb{E}(Var(X|Y)) + Var(\mathbb{E}(X|Y)). $$
Questa è anche nota come ``legge della varianza totale''. Iniziamo da $Var(\mathbb{E}(X|Y))$, i valori attesi $\mathbb{E}(X|Y)$ sono noti dal punto precedente, per calcolare la loro varianza abbiamo bisogno della media $\mathbb{E}(\mathbb{E}(X|Y)) = \mathbb{E}(X)$, anch’essa nota, dopodiché: 
$$ Var(\mathbb{E}(X|Y)) = \sum_y (\mathbb{E}(X|Y=y) - \mathbb{E}(X))^2 p_Y(y) = 0.8020634. $$
Per quanto riguarda $\mathbb{E}(Var(X|Y))$ abbiamo bisogno di un po’ più di lavoro: 
\begin{itemize}
    \item $Var(X|Y=1) = \sum_x (x - \mathbb{E}(X|Y=1))^2 p_{X|Y}(x|y=1) = 38.9466$
    \item $Var(X|Y=2) = \sum_x (x - \mathbb{E}(X|Y=2))^2 p_{X|Y}(x|y=2) = 41.5236$
    \item $Var(X|Y=3) = \sum_x (x - \mathbb{E}(X|Y=3))^2 p_{X|Y}(x|y=3) = 41.8146$
\end{itemize}
dopodiché 
$$ \mathbb{E}(Var(X|Y)) = \sum_y Var(X|Y=y)p_Y(y) = 40.7688. $$
Sommando i due risultati parziali otteniamo la risposta desiderata.
La risposta è: 41.5708634


% ==========================================
%               ESERCIZIO 4
% ==========================================
\vspace{0.5cm}
\begin{tcolorbox}[colback=lightergray!10, colframe=tocsec, title=\textbf{Esercizio 4}]
Una piccola azienda confeziona e vende sacchetti di caramelle gommose. Ogni sacchetto contiene caramelle di due diversi colori: rosse e gialle.
In fabbrica i sacchetti sono riempiti da due diverse macchine – che operano in maniera indipendente. Una versa nel sacchetto le caramelle rosse, l’altra quelle gialle. Le macchine non sono molto precise, la prima macchina immette nel sacchetto in media 34 caramelle rosse, con una deviazione standard di $\sigma_1=2$; dopodiché il sacchetto passa sotto la seconda macchina, che ha una media di 29 caramelle gialle con deviazione standard $\sigma_2=6$.
Una caramella pesa un grammo e il prezzo dei sacchetti dipende dal loro peso, in particolare il prezzo è 0.15 Euro volte il peso del sacchetto.
Chiamiamo $X$ e $Y$ le variabili casuali che rappresentano il numero (aleatorio) di caramelle distribuite dalle due macchine. $\mathbb{E}(X)$, $\mathbb{E}(Y)$, $Var(X)=\sigma_1^2$ e $Var(Y)=\sigma_2^2$ sono dati dei problemi.
\end{tcolorbox}

\subsubsection*{Quesito 1}
1. Quanto costa in media un sacchetto di caramelle?

\textbf{Soluzione:}
Sfruttiamo la linearità dell'operazione valore atteso. Considerando il costo unico di 0.15 per entrambe le caramelle:
$$ 0.15 \cdot \mathbb{E}(X+Y) = 0.15 \cdot (\mathbb{E}(X)+\mathbb{E}(Y)) = 0.15 \cdot (34 + 29) = 0.15 \cdot 63 = 9.45 $$
La risposta è: 9.45

\vspace{0.3cm}
\textit{Ora, le caramelle rosse richiedono una preparazione più lunga e costosa, per cui i produttori decidiono di modificare il prezzo delle caramelle: una caramella rossa costa 0.25, mentre una gialla costa 0.15.}

\subsubsection*{Quesito 2}
2. Quanto costa in media un sacchetto di caramelle, con i nuovi prezzi?

\textbf{Soluzione:}
Siano $a$ e $b$ i due costi.
$$ \mathbb{E}(aX+bY) = a\mathbb{E}(X) + b\mathbb{E}(Y) = 0.25(34) + 0.15(29) = 8.5 + 4.35 = 12.85 $$
La risposta è: 12.85

\subsubsection*{Quesito 3}
3. Qual è la varianza del prezzo di un sacchetto di caramelle (con i nuovi prezzi)?

\textbf{Soluzione:}
$$ Var(aX+bY) = a^2 Var(X) + b^2 Var(Y) + 2ab \cdot Cov(X,Y) $$
$X$ e $Y$ sono indipendenti, ergo $Cov(X,Y)=0$.
$$ Var(0.25X + 0.15Y) = (0.25)^2 (4) + (0.15)^2 (36) = 0.0625 \cdot 4 + 0.0225 \cdot 36 = 0.25 + 0.81 = 1.06 $$
La risposta è: 1.06

\subsubsection*{Quesito 4}
4. Qual è la probabilità $p$ che il costo di un sacchetto di caramelle abbia un costo maggiore o uguale a 11.4? Stabilire un limite superiore al valore di $p$ usando la disuguaglianza di Markov.

\textbf{Soluzione:}
Dalla Disuguaglianza di Markov abbiamo che per una v.a. non negativa $W = aX+bY$:
$$ P(aX+bY \geq \alpha) \leq \frac{\mathbb{E}(aX+bY)}{\alpha} $$
da cui sostituendo la soglia $\alpha = 11.4$:
$$ P(Costo \geq 11.4) \leq \frac{12.85}{11.4} \approx 1.127193 $$
La risposta è: 1.127193

Ora, in alcuni casi (quando $\alpha < \mathbb{E}(aX+bY)$) il valore restituito dalla disuguaglianza di Markov è maggiore di 1. 
Questo è ok, ma se volessimo essere proprio precisi potremmo restituire 1, ossia il limite più stringente e fisicamente sensato per una probabilità. Una versione migliorata sarebbe quindi 
$$ P(aX+bY \geq \alpha) \leq \min\left\{\frac{\mathbb{E}(aX+bY)}{\alpha}, 1\right\} $$

\newpage


% ==========================================
%               ESERCIZIO 5
% ==========================================
\vspace{0.5cm}
\begin{tcolorbox}[colback=lightergray!10, colframe=tocsec, title=\textbf{Esercizio 5}]
Siano $X$ e $Y$ due variabili casuali indipendenti distribuite secondo una Poisson di parametri $\lambda_1=11$ e $\lambda_2=3$ rispettivamente.
\end{tcolorbox}

\subsubsection*{Quesito 1}
1. Dopo aver determinato $\mathbb{E}(X+Y)$ usare la disuguaglianza di Markov per ottenere un limite superiore della probabilità che la somma $X+Y$ sia maggiore o uguale a 23.

\textbf{Soluzione:}
Essendo indipendenti, la loro somma è ancora una Poisson con parametro pari alla somma dei parametri:
$$ \mathbb{E}(X+Y) = \mathbb{E}(X) + \mathbb{E}(Y) = \lambda_1 + \lambda_2 = 14. $$
Per la disuguaglianza di Markov:
$$ P(X+Y \geq a) \leq \frac{\mathbb{E}(X+Y)}{a} = \frac{14}{23} \approx 0.6086957 $$
La risposta è: 0.6086957

\subsubsection*{Quesito 2}
2. Usando la disuguaglianza di Markov ottenere un limite superiore della probabilità condizionata $P(X \geq 8 \mid X+Y=23)$.

\textbf{Soluzione:}
Abbiamo visto in esercizi precedenti che la distribuzione di $X$ data la somma $X+Y=n$ è binomiale di parametri $N=n$ e $p = \frac{\lambda_1}{\lambda_1+\lambda_2}$.
Il valore atteso in questo caso condizionato è $\mathbb{E}(X \mid X+Y=n) = Np = n \frac{\lambda_1}{\lambda_1+\lambda_2} = 23 \frac{11}{14}$.
Applicando Markov alla distribuzione condizionata:
$$ P(X \geq k \mid X+Y=23) \leq \frac{\mathbb{E}(X \mid X+Y=n)}{k} = \frac{23 \cdot \frac{11}{14}}{8} \approx 2.2589286 $$
La risposta è: 2.2589286

\subsubsection*{Quesito 3}
3. Sia ora $X$ una v.a. normale standard e sia $Y=aX$ dove $a=1.2$. Usare la disuguaglianza di Chebychev per ottenere un limite inferiore alla probabilità $P(|Y|<z)$ per $z=2.4$.

\textbf{Soluzione:}
Da $X \sim N(0,1)$ segue che $Y \sim N(0, a^2)$ (check it!) per le proprietà delle trasformazioni lineari di normali.
Sfruttando l'evento complementare:
$$ P(|Y| < z) = 1 - P(|Y| \geq z) $$
La disuguaglianza di Chebychev afferma che $P(|Y| \geq z) \leq \frac{\sigma_Y^2}{z^2} = \frac{a^2}{z^2}$.
Per cui se da 1 tolgo qualcosa di più grande o uguale non ho più l’uguaglianza bensì la disuguaglianza (si inverte il segno della relazione d'ordine originaria per l'evento complementare):
$$ P(|Y| < z) = 1 - P(|Y| \geq z) \geq 1 - \frac{a^2}{z^2} $$
e poiché abbiamo che $1 - \frac{a^2}{z^2} \leq P(|Y| < z)$, abbiamo ottenuto un limite inferiore.
Sostituendo $a=1.2$ e $z=2.4$:
$$ P(|Y| < 2.4) \geq 1 - \frac{1.44}{5.76} = 1 - 0.25 = 0.75 $$
La risposta è: 0.75

\subsubsection*{Quesito 4}
4. Qual è il valore esatto di $P(|Y|<2.4)$?

\textbf{Soluzione:}
Svolgendo il modulo come $P(-2.4 < Y < 2.4)$ e passando tramite la funzione di ripartizione continua, possiamo standardizzare i bordi dividendo per la deviazione standard $a=1.2$:
$$ \Phi\left(\frac{2.4}{1.2}\right) - \Phi\left(\frac{-2.4}{1.2}\right) = \Phi(2) - \Phi(-2) = 2\Phi(2) - 1 \approx 0.9544997 $$
La risposta è: 0.9544997


% ==========================================
%               ESERCIZIO 6
% ==========================================
\vspace{0.5cm}
\begin{tcolorbox}[colback=lightergray!10, colframe=tocsec, title=\textbf{Esercizio 6}]
Proprietà di campioni binomiali.
\end{tcolorbox}

\subsubsection*{Quesito 1}
1. Lanciando 12 volte una moneta equilibrata, l’evento ottenere $n/2$ teste è più (strettamente più) probabile del suo complementare? Rispondere TRUE o FALSE.

\textbf{Soluzione:}
Per $n=12$ e $p=0.5$, la probabilità di ottenere esattamente 6 teste è $\binom{12}{6}(0.5)^{12} \approx 0.225$. L'evento complementare (non ottenere 6 teste) ha probabilità $1 - 0.225 = 0.775$. Quindi non è strettamente più probabile.
La risposta è: FALSE

\subsubsection*{Quesito 2}
2. Qual è la probabilità di ottenere almeno 11 teste?

\textbf{Soluzione:}
Calcoliamo $\binom{12}{11}(0.5)^{12} + \binom{12}{12}(0.5)^{12} = \frac{12 + 1}{4096} = \frac{13}{4096} \approx 0.0031738$.
La risposta è: 0.0031738

\subsubsection*{Quesito 3}
3. In quanti possibili modi si possono distribuire 11 successi in una successione di 12 prove?

\textbf{Soluzione:}
Questo è esattamente il coefficiente binomiale $\binom{12}{11} = 12$.
La risposta è: 12

\subsubsection*{Quesito 4}
4. Supponendo che ora la probabilità di ottenere testa sia 0.74, qual è il numero medio di teste in 12 lanci?

\textbf{Soluzione:}
La speranza matematica di un processo Binomiale si calcola come:
$$ np = 12 \cdot 0.74 = 8.88 $$
La risposta è: 8.88

\newpage

% ==========================================
%               ESERCIZIO 7
% ==========================================
\vspace{0.5cm}
\begin{tcolorbox}[colback=lightergray!10, colframe=tocsec, title=\textbf{Esercizio 7}]
Abbiamo due variabili aleatorie binomiali indipendenti $X$ e $Y$, di parametri rispettivamente $(n=6, p=0.501)$ e $(n=5, p=0.501)$. Prendiamone anche la somma $S=X+Y$, che, come è noto, è distribuita come una binomiale di parametri $(n=11, p=0.501)$.
\end{tcolorbox}

\subsubsection*{Quesito 1}
Qual è la probabilità che $S$ sia compresa nell’intervallo $[0,7)$?

\textbf{Soluzione:}
Sappiamo che $S$ è binomiale di parametri $(n=6+5=11, p=0.501)$.
Pertanto possiamo usare le funzioni \texttt{pbinom} o \texttt{dbinom} per risolvere il problema, ad esempio calcolando la prima in $0-1$ e $7-1$ e sottraendo il primo valore dal secondo. Essendo un intervallo aperto a destra su una variabile discreta, $S \in [0,7)$ equivale a calcolare la funzione di ripartizione valutata in 6, i.e. $P(S \leq 6)$.
La risposta è: 0.7233256

\subsubsection*{Quesito 2}
Se sappiamo che $S=5$, qual è la probabilità che $X=5$?

\textbf{Soluzione:}
Stiamo cercando quale sia la densità discreta di $X|S$. Per farlo, partiamo dalla definizione
$$ \begin{aligned} \varphi_{X|S}(k|5) &= \frac{\varphi_X(k) \cdot \varphi_Y(5-k)}{\varphi_S(5)} \\ &= \frac{\binom{6}{k}p^k(1-p)^{6-k} \cdot \binom{5}{5-k}p^{5-k}(1-p)^{5-(5-k)}}{\binom{6+5}{5}p^5(1-p)^{6+5-5}} \\ &= \frac{\binom{6}{k}\binom{5}{5-k}}{\binom{6+5}{5}} \end{aligned} $$
in cui abbiamo usato l’indipendenza di $X$ e $Y$ e la forma della densità discreta determinata nel primo quesito.
Quella scritta nell’ultima riga è la densità discreta di una variabile aleatoria ipergeometrica. In altre parole, $X|S$ ha legge ipergeometrica di parametri $(6, 5, 5)$. Dobbiamo allora calcolare il valore della densità discreta di una tale ipergeometrica in $k=5$, cosa che possiamo fare con l’ausilio della funzione R \texttt{dhyper(5, 6, 5, 5)}.
La risposta è: 0.012987

\newpage

\subsubsection*{Quesito 3}
Se sappiamo che $S=6$, qual è la probabilità che $Y \in [1,2]$?

\textbf{Soluzione:}
In questo caso ci occorre la legge di $Y|S$, ma essa è, come abbiamo visto nella soluzione del Quesito 2 (scambiando i ruoli di $X$ e $Y$)
$$ \varphi_{Y|S}(k|6) = \frac{\binom{5}{k}\binom{6}{6-k}}{\binom{5+6}{6}}, $$
cioè un’ipergeometrica di parametri $(5, 6, 6)$. Dobbiamo sommare i valori della sua densità discreta per tutti i $k$ tra 1 e 2, oppure possiamo fare la differenza tra la funzione di ripartizione, che in R si calcola con la funzione \texttt{phyper} nel caso dell’ipergeometrica, valutata nei punti 2 e $1-1$.
La risposta è: 0.3896104

% ==========================================
%               ESERCIZIO 8
% ==========================================
\vspace{0.5cm}
\begin{tcolorbox}[colback=lightergray!10, colframe=tocsec, title=\textbf{Esercizio 8}]
In un cinema multisala vengono venduti solo biglietti in prevendita online e non all’entrata. Sapendo che il 20\% delle persone che acquistano un biglietto online per un film poi non si presenta al cinema, i gestori del sito di prevendita del cinema vendono 16 biglietti per il film proiettato nella sala 1, che ha 14 posti e 17 biglietti per quello nella sala 2, da 15 posti.
\end{tcolorbox}

\subsubsection*{Quesito 1}
Calcolare la probabilità che almeno una persona che ha comprato il biglietto online (per il primo film) non trovi posto nella sala 1.

\textbf{Soluzione:}
Sia $X_1$ la v.a che conta il numero di persone che si presentano al cinema per il film della sala 1.
Definendo “successo” l’evento “La persona che ha acquistato il biglietto si presenta al cinema”, la probabilità di successo è $p=1 - 0.20 = 0.8$. Quindi $X_1$ conterà il numero di successi (persone che si presentano) su 16 tentativi (biglietti venduti per la sala 1), cioè $X_1 \sim Bin(N_1=16, p=0.8)$.
Se si presentano al cinema per il primo film più di 14 persone, cioè più del numero di posti della sala 1, allora almeno una persona non troverà posto, quindi in definitiva dobbiamo calcolare $P(X_1 > 14)$:
$$ P(X_1 > 14) = \sum_{k=15}^{N_1} \binom{N_1}{k} p^k (1-p)^{N_1-k} = \binom{16}{15}(0.8)^{15}(0.2)^1 + \binom{16}{16}(0.8)^{16}(0.2)^0 $$
Con R possiamo farlo con \texttt{pbinom(q, n, p, lower.tail = FALSE)}, con \texttt{q=14}, \texttt{n=16} e \texttt{p=0.8}.

La risposta è: 0.1407375

\subsubsection*{Quesito 2}
Calcolare la probabilità che nessun acquirente del biglietto online per il secondo film abbia problemi a trovare posto nella sala 2.

\textbf{Soluzione:}
Ragionando come prima, chiamiamo $X_2$ la v.a. che conta il numero di persone che si presentano al cinema per il film della sala 2, allora $X_2 \sim Bin(N_2=17, p=0.8)$.
In questo caso dobbiamo calcolare la probabilità che il numero di persone che si presentano per il secondo film sia al più uguale al numero di posti della sala 2, i.e. $P(X_2 \leq 15)$.
Questa è calcolabile sfruttando l'evento complementare: $1 - P(X_2 > 15)$.

La risposta corretta è: 0.8817805

\subsubsection*{Quesito 3}
Calcolare il numero di biglietti che si potrebbero vendere per la sala 1, accettando che la probabilità di scontentare una persona (i.e. lasciarla fuori dalla sala) non sia superiore a 0.46.

\textbf{Soluzione:}
Sia $\alpha=0.46$. Fissata la probabilità $\alpha$, dobbiamo determinare il parametro $N_1$ tale per cui 
$$ \alpha \geq P(X_1 > 14) = \sum_{k=15}^{N_1} \binom{N_1}{k} p^k (1-p)^{N_1-k}. $$
Possiamo procedere per tentativi con un calcolo esplicito. Osserviamo che dobbiamo trovare il più grande intero $N_1$ per cui $P(X_1 > 14) \leq \alpha$. Allora, iniziamo provando con $N_1 = 14+1 = 15$ per capire qual è la probabilità di scontentare qualcuno vendendo solo un biglietto in più rispetto ai posti della sala 1. Se già questa probabilità è maggiore di $\alpha$, dobbiamo prendere $N_1=14$, cioè dobbiamo vendere un numero di biglietti pari al numero di posti in sala.
Altrimenti, aumentiamo di un’unità il valore di $N_1$ e calcoliamo la nostra probabilità; iteriamo il procedimento fino a quando non troviamo quel valore di $N_1$ per cui
$$ \sum_{k=15}^{N_1} \binom{N_1}{k} p^k (1-p)^{N_1-k} \leq \alpha \quad \text{e} \quad \sum_{k=15}^{N_1+1} \binom{N_1+1}{k} p^k (1-p)^{(N_1+1)-k} > \alpha. $$
Tale valore di $N_1$ sarà la nostra soluzione. Sostituendo si verifica che 17 biglietti mantengono il rischio accettabile, ma 18 no.

La risposta è: 17

\subsubsection*{Quesito 4}
In media quante persone si presentano per il film della sala 2?

\textbf{Soluzione:}
Dobbiamo calcolare il valore atteso di una v.a. binomiale di parametri $N_2=17$ e $p=0.8$, cioè $\mathbb{E}[X_2] = N_2 \cdot p = 17 \cdot 0.8 = 13.6$.

La risposta è: 13.6


% ==========================================
%               ESERCIZIO 9
% ==========================================
\vspace{0.5cm}
\begin{tcolorbox}[colback=lightergray!10, colframe=tocsec, title=\textbf{Esercizio 9}]
Un giocatore di pallacanestro ha una probabilità del 41\% di fare canestro con un tiro da 2 punti.
\end{tcolorbox}

\subsubsection*{Quesito 1}
Determinare la probabilità che il giocatore non faccia più di 54 punti in 59 tiri tutti da 2 punti, utilizzando l’approssimazione normale (Teorema Limite Centrale) e la correzione di continuità.

\textbf{Soluzione:}
Sia $X_{59}$ la v.a che conta il numero di canestri su 59 tiri da 2 punti, quindi $X_{59} \sim Bin(n=59, p=0.41)$. Dato che ogni canestro segnato fa guadagnare 2 punti, dovremo calcolarci la probabilità che il giocatore faccia al più $\frac{54}{2} = 27$ canestri.
Sappiamo che $\mathbb{E}[X_{59}] = n \cdot p = 59 \cdot 0.41 = 24.19$ e $Var(X_{59}) = np(1-p) = 24.19 \cdot 0.59 = 14.2721$.
Allora, usando l’approssimazione con TLC e la correzione per continuità estendiamo il bordo di $0.5$ per coprire interamente l'istogramma discreto del valore 27:

\newpage

$$ P(X_{59} \leq 27) = P\left(\frac{X_{59} - 24.19}{\sqrt{14.2721}} \leq \frac{(27+0.5) - 24.19}{\sqrt{14.2721}}\right) \approx \Phi\left(\frac{27.5 - 24.19}{\sqrt{14.2721}}\right) = \Phi\left(\frac{3.31}{3.7778}\right) $$
La risposta è: 0.8095288

\subsubsection*{Quesito 2}
Qual è l’errore assoluto tra la probabilità ottenuta con l’approssimazione calcolata al quesito 1 e il valore esatto?

\textbf{Soluzione:}
Dobbiamo confrontare la soluzione del quesito 1 con il valore della funzione di ripartizione della variabile aleatoria binomiale $X_{59}$ calcolata esplicitamente in 27, $F_{X_{59}}(27)$. Possiamo, come al solito, usare per questo la funzione in R \texttt{pbinom}.
Per concludere, sottraiamo al valore esatto ottenuto quello approssimato calcolato al quesito 1, e prendiamone il valore assoluto.
$$ |P(X_{59} \leq 27)_{\text{esatto}} - P(X_{59} \leq 27)_{\text{TLC}}| $$
La risposta corretta è: $4.4128511 \times 10^{-4}$

\subsubsection*{Quesito 3}
Determinare il numero minimo di tiri che il giocatore deve effettuare affinché la probabilità di segnare almeno 84 punti sia non inferiore a 0.89, utilizzando l’approssimazione normale e la correzione di continuità.
\textit{Suggerimento: \texttt{qnorm(p, mean=0, sd=1)} calcola il quantile corrispondente alla probabilità p secondo la distribuzione normale standard.}

\textbf{Soluzione:}
Sia $X_n$ la v.a. che indica il numero di canestri su $n$ tiri, quindi $X_n \sim Bin(n, p=0.41)$.
Dobbiamo quindi determinare $n$ in modo che si verifichi $P(X_n \geq 42) \geq 0.89$, dato che per ottenere almeno 84 punti bisogna fare almeno 42 canestri. Utilizzando l’approssimazione normale e la correzione di continuità (da $42$ scendiamo a $41.5$), abbiamo:
$$ \begin{aligned} 0.89 &\leq P(X_n \geq 42) \\ &= P(X_n > 42 - 0.5) \\ &= P\left(\frac{X_n - n \cdot 0.41}{\sqrt{n \cdot 0.41(1-0.41)}} > \frac{(42 - 0.5) - n \cdot 0.41}{\sqrt{n \cdot 0.41(1-0.41)}}\right) \\ &\approx 1 - \Phi\left(\frac{41.5 - n \cdot 0.41}{\sqrt{n \cdot 0.41 \cdot 0.59}}\right) \end{aligned} $$
da cui segue riordinando la disuguaglianza:
$$ \Phi\left(\frac{41.5 - n \cdot 0.41}{\sqrt{n \cdot 0.2419}}\right) \leq 0.11 \iff \frac{41.5 - n \cdot 0.41}{\sqrt{n \cdot 0.2419}} \leq \Phi^{-1}(0.11). $$
Come suggerito, possiamo quindi usare \texttt{qnorm} per calcolarci $\Phi^{-1}(0.11)$, dopodiché risolvendo la disuguaglianza di secondo grado associata si ottiene $n \geq 117.1442415$, ossia arrotondando per eccesso per avere almeno la probabilità richiesta $n \geq 118$.

La risposta è: 118


% ==========================================
%               ESERCIZIO 10
% ==========================================
\vspace{0.5cm}
\begin{tcolorbox}[colback=lightergray!10, colframe=tocsec, title=\textbf{Esercizio 10}]
Delle caramelle artigianali hanno un peso distribuito come una normale con media $\mu=5.579$ g e deviazione standard $\sigma=0.604$ g. Al controllo qualità le caramelle con un peso (strettamente) superiore a 6.3606904 g o (strettamente) inferiore a 4.717553 g vengono scartate.
\end{tcolorbox}

\subsubsection*{Quesito 1}
Qual è la probabilità che una caramella sia scartata al controllo qualità?

\textbf{Soluzione:}
Sia $X \sim N(\mu=5.579, \sigma^2=0.364816)$ la v.a. che rappresenta il peso di una caramella artigianale.
La probabilità che una caramella sia scartata è uguale alla somma delle code di reiezione:
$$ \begin{aligned} P(X > 6.3606904) + P(X < 4.717553) &= (1 - P(X \leq 6.3606904)) + P(X < 4.717553) \\ &= 0.1747 \end{aligned} $$
Dobbiamo ora calcolare $P(X \leq x)$. Per farlo possiamo usare la standardizzazione:
$$ P(X \leq x) = \Phi\left(\frac{x - \mu}{\sigma}\right). $$
Possiamo quindi usare \texttt{pnorm((x-mu)/sigma)} o le tavole. In alternativa, possiamo usare direttamente \texttt{pnorm(x, mean = mu, sd = sigma)}.

La risposta è: 0.1747

\subsubsection*{Quesito 2}
Qual è la probabilità che se vengono controllate 73 caramelle, queste pesino complessivamente più di 412.737213 g ?

\textbf{Soluzione:}
Sia $S_n$ la v.a. che indica il peso di $n$ caramelle. Dato che $S_n$ è somma di $n$ v.a. gaussiane indipendenti di media $\mu$ e varianza $\sigma^2$, avremo $S_n \sim N(n \cdot \mu, n \cdot \sigma^2)$ (la Gaussiana è riproducibile rispetto alla somma).
Quindi, per $n=73$ e soglia $W=412.737213$, dobbiamo calcolare:
$$ P(S_n > W) = P\left(\frac{S_n - n\mu}{\sqrt{n}\sigma} > \frac{W - n\mu}{\sqrt{n}\sigma}\right) = P\left(Z > \frac{W - n\mu}{\sqrt{n}\sigma}\right) = 1 - \Phi\left(\frac{W - n\mu}{\sqrt{n}\sigma}\right), $$
dove $Z \sim N(0,1)$.
Possiamo quindi usare \texttt{1-pnorm(x)} o le tavole.
Alternativamente \texttt{1-pnorm(x, mean = n*mu, sd = sqrt(n)*sigma)}.

La risposta è: 0.1445723

\subsubsection*{Quesito 3}
Qual è la probabilità che su 73 caramelle al più 11 siano da scartare?

\textbf{Soluzione:}
Sia $Y$ la v.a. che conta le caramelle da scartare sulle $n=73$ da controllare. Dobbiamo allora calcolare $P(Y \leq 11)$.
Osserviamo che dobbiamo contare il numero di successi (caramelle scartate al controllo) su $n=73$ tentativi (caramelle controllate). Per cui $Y \sim Bin(n, p)$ dove $p=0.1747$ è la probabilità che una caramella singola sia scartata, calcolata al quesito 1.
Possiamo quindi usare la funzione in R \texttt{pbinom(11, 73, p)}.

La risposta corretta è: 0.3606461


% ==========================================
%               ESERCIZIO 11
% ==========================================
\vspace{0.5cm}
\begin{tcolorbox}[colback=lightergray!10, colframe=tocsec, title=\textbf{Esercizio 11}]
Abbiamo quattro variabili aleatorie: $X$, $Y$, $W$, $Z$.
Le prime due, $X$ e $Y$, sono indipendenti tra loro e hanno entrambe distribuzione poissoniana, di parametri $\lambda_X=2$ e $\lambda_Y=10$ rispettivamente.
La variabile $Z$ è una normale (o Gaussiana) standard, mentre $W=1.1 \cdot Z$.
\end{tcolorbox}

\subsubsection*{Quesito 1}
Dare la migliore possibile stima dall’alto della probabilità che la somma $X+Y$ sia maggiore o uguale a 18, usando la disuguaglianza di Markov.

\textbf{Soluzione:}
La disuguaglianza di Markov, per una variabile aleatoria non negativa $T$, è:
$$ P(T \geq a) \leq \frac{\mathbb{E}[T]}{a}. $$
Per poter usare la disuguaglianza di Markov non abbiamo bisogno di conoscere (o ricordare) la funzione di ripartizione o la densità discreta. Abbiamo però bisogno di calcolare la speranza della variabile aleatoria $X+Y$: la speranza della somma è la somma delle speranze, quindi, siccome $X$ e $Y$ sono di Poisson e quindi riproducibili,
$$ \mathbb{E}[X+Y] = \mathbb{E}[X] + \mathbb{E}[Y] = \lambda_X + \lambda_Y = 12. $$
Ora possiamo scrivere la disuguaglianza nel nostro caso (osserviamo che entrambe le variabili aleatorie sono non negative e dunque tale è anche la loro somma):
$$ P(X+Y \geq a) \leq \frac{\mathbb{E}[X+Y]}{a} \implies P(X+Y \geq 18) \leq \frac{12}{18}. $$
Questa è la miglior stima possibile purché sia minore o uguale di 1.

La risposta corretta è: 0.6666667

\subsubsection*{Quesito 2}
Dare la migliore possibile stima dall’alto della probabilità (condizionata) $P(X \geq 14.03 \mid X+Y=18)$, usando la disuguaglianza di Markov.

\textbf{Soluzione:}
Sappiamo che la distribuzione di $X$ sapendo che $X+Y=n$ è una binomiale di parametri $n$ e $p = \frac{\lambda_X}{\lambda_X+\lambda_Y}$.
Il valore atteso condizionato è allora:
$$ \mathbb{E}[X \mid X+Y=n] = n \cdot p = n\frac{\lambda_X}{\lambda_X+\lambda_Y}. $$
Possiamo ora andare a considerare la disuguaglianza di Markov per lo spazio di probabilità condizionato:
$$ P(X \geq 14.03 \mid X+Y=18) \leq \frac{\mathbb{E}[X \mid X+Y=18]}{14.03} = \frac{18 \cdot 2}{14.03 \cdot (2+10)}. $$
Questa, però, è la miglior stima possibile solamente se è minore o uguale di 1 e se consideriamo la natura discreta della variabile.
Infatti siccome la stima data dalla disuguaglianza di Markov decresce al crescere del denominatore, ci conviene prendere il denominatore più grande possibile, a parità di probabilità a primo membro. In questo caso la cosa è rilevante, perché $X$ assume solo valori in $\mathbb{N}$, quindi l'evento $X \geq y$ equivale strettamente all'evento $X \geq \lceil y \rceil$.
Dunque possiamo stimare più finemente scalando la soglia all'intero immediatamente superiore:
$$ P(X \geq 14.03 \mid X+Y=18) = P(X \geq 15 \mid X+Y=18) \leq \frac{\mathbb{E}[X \mid X+Y=18]}{15} = \frac{18 \cdot 2}{15 \cdot (2+10)}. $$
La risposta è: 0.2

\subsubsection*{Quesito 3}
Dare una stima dal basso della probabilità $P(|W| < 1.6)$, usando la disuguaglianza di Chebychev.

\textbf{Soluzione:}
Siccome $Z \sim N(0,1)$, ha speranza 0 e varianza 1, per cui $W = 1.1 \cdot Z$ ha speranza 0 e varianza $1.1^2 = 1.21$ (si tratta di una trasformazione lineare).
Non solo, in realtà non ci occorre per questo quesito (pur essendo utile nel prossimo), ma conosciamo la distribuzione di $W$: $W \sim N(0, 1.1^2)$.
Osserviamo che possiamo riscrivere la quantità che ci interessa nel modo seguente:
$$ P(|W| < 1.6) = 1 - P(|W| \geq 1.6). $$
Scriverla in questa forma ci permette di mettere in evidenza il termine che possiamo stimare con la disuguaglianza di Chebychev, che quantifica quanto la variabile si allontana dal proprio valore atteso nullo:
$$ P(|W| \geq 1.6) \leq \frac{1.21}{1.6^2} = \frac{1.21}{2.56}, $$
dove abbiamo usato il fatto che la varianza di $W$ è 1.21.
Questa è una stima dall’alto, ma siccome la stiamo sottraendo da 1, l’uguaglianza $P(|W| < 1.6) = 1 - P(|W| \geq 1.6)$ diventa la disuguaglianza rovesciata:
$$ P(|W| < 1.6) = 1 - P(|W| \geq 1.6) \geq 1 - \frac{1.21}{2.56}. $$
Abbiamo dunque ottenuto la stima dal basso cercata.

La risposta è: 0.5273438

\subsubsection*{Quesito 4}
Qual è l’errore assoluto (rispetto al valore vero di $P(|W| < 1.6)$) fatto nella stima precedente?

\textbf{Soluzione:}
Qui, a differenza del quesito precedente, vogliamo sfruttare il fatto che sappiamo la forma esatta della distribuzione di $W$. Possiamo infatti riscrivere la quantità cercata nel modo seguente:
$$ P(|W| < 1.6) = P(-1.6 < W < 1.6) = F_W(1.6) - F_W(-1.6). $$
Osserviamo che la media di $W$ è 0 (abbiamo solo riscalato $Z$, non l’abbiamo traslata). Questo ci dice che la densità è perfettamente simmetrica rispetto all'origine. Possiamo fare la stessa considerazione dopo aver standardizzato la variabile e sfruttando le proprietà della CDF normale $\Phi(-x) = 1 - \Phi(x)$:
$$ P(|W| < 1.6) = \Phi\left(\frac{1.6}{1.1}\right) - \Phi\left(\frac{-1.6}{1.1}\right) = 2\Phi\left(\frac{1.6}{1.1}\right) - 1, $$
quantità che possiamo calcolare in R come \texttt{2*pnorm(x)-1} con \texttt{x=1.6/1.1} o con l’ausilio delle tavole.
Per finire, non ci resta che sottrarre a questo valore quello approssimato ottenuto al quesito 3 e prenderne il valore assoluto.

La risposta corretta è: 0.3268611


% ==========================================
%               ESERCIZIO 12
% ==========================================
\vspace{0.5cm}
\begin{tcolorbox}[colback=lightergray!10, colframe=tocsec, title=\textbf{Esercizio 12}]
Sia $(X_1,\dots,X_n)$ un campione casuale da una distribuzione $N(\mu=22, \sigma^2=5)$ (i.e. le $X_i$ sono indipendenti ed identicamente distribuite).
Si considerino i seguenti stimatori della media:
$$ T_1 = \frac{\sum_{i=1}^n X_i}{n} = \bar{X}_n \quad , \quad T_2 = 2\bar{X}_n - X_1 $$
Sia inoltre $n=30$.
\end{tcolorbox}

\subsubsection*{Quesito 1}
Calcolare la distorsione degli stimatori $T_1, T_2$.

\textbf{Soluzione:}
Per calcolare la distorsione di uno stimatore $T$, $\mathbb{E}(T-\mu)$, dobbiamo innanzitutto calcolarne il valore atteso.
Per la media campionaria è noto che:
$$ \mathbb{E}(T_1) = \mathbb{E}(\bar{X}_n) = \mathbb{E}(X_i) = \mu $$
Applicando la linearità per il secondo stimatore:
$$ \mathbb{E}(T_2) = 2\mathbb{E}(\bar{X}_n) - \mathbb{E}(X_1) = 2\mu - \mu = \mu $$
Per cui entrambi gli stimatori sono centrati sul parametro reale, e la loro distorsione risulta nulla: $\mu - \mu = 0$. Entrambi gli stimatori sono non distorti.

La risposta è: 0, 0

\subsubsection*{Quesito 2}
Calcolare la varianza degli stimatori $T_1, T_2$.

\textbf{Soluzione:}
La varianza di uno stimatore $T$ è, per definizione, $Var(T) = \mathbb{E}(T - (\mathbb{E}(T))^2)$.
Ricordiamo inoltre che la varianza per trasformazioni lineari di variabili casuali soddisfa:
$$ Var(aX + bY + c) = a^2 Var(X) + b^2 Var(Y) + 2ab \cdot Cov(X,Y) $$
Per il primo stimatore, la varianza della media campionaria si riduce per l'indipendenza:
$$ Var(T_1) = \frac{Var(X_i)}{n} = \frac{\sigma^2}{n} $$
Per il secondo stimatore la questione è più complessa a causa della covarianza non nulla tra la media e il suo primo elemento:
$$ \begin{aligned} Var(T_2) &= 4Var(\bar{X}_n) + Var(X_1) - 4Cov(\bar{X}_n, X_1) \\ &= 4\frac{Var(X_i)}{n} + Var(X_1) - 4Cov\left(\frac{X_1}{n} + \frac{X_2}{n} \dots + \frac{X_n}{n}, X_1\right) \\ &= \frac{4\sigma^2}{n} + \sigma^2 - 4\left(Cov\left(\frac{X_1}{n}, X_1\right) + \dots + Cov\left(\frac{X_n}{n}, X_1\right)\right) \\ &= \frac{4\sigma^2}{n} + \sigma^2 - 4\left(\frac{1}{n}Var(X_1) + 0 + \dots + 0\right) \\ &= \frac{4\sigma^2}{n} + \sigma^2 - 4\frac{\sigma^2}{n} \\ &= \sigma^2 \end{aligned} $$
Sostituendo $\sigma^2 = 5$ e $n=30$, troviamo i valori richiesti.

La risposta corretta è: 0.166666666666667, 5

\subsubsection*{Quesito 3}
Calcolare l’errore quadratico medio degli stimatori $T_1, T_2$.

\textbf{Soluzione:}
A questo punto abbiamo già tutto quello che ci serve, poiché l’errore quadratico medio (mean square error, MSE) di uno stimatore $T$ per il parametro $\mu$ è definito dalla scomposizione bias-varianza: 
$$ MSE(T, \mu) = Var(T) + (\mathbb{E}(T - \mu))^2. $$
Poiché in precedenza abbiamo accertato che entrambi sono non distorti (hanno bias nullo), l'MSE coincide matematicamente con la sola varianza:
$$ MSE(T_1, \mu) = Var(T_1), \quad MSE(T_2, \mu) = Var(T_2) $$
La risposta è: 0.166666666666667, 5


% ==========================================
%               ESERCIZIO 13
% ==========================================
\vspace{0.5cm}
\begin{tcolorbox}[colback=lightergray!10, colframe=tocsec, title=\textbf{Esercizio 13}]
Durante uno studio medico si misura il parametro toracico di $n=6$ persone, i valori sono riportati nella seguente tabella:
\begin{center}
84.60, 79.10, 84.50, 79.60, 74.00, 74.10
\end{center}
\end{tcolorbox}
\subsubsection*{Quesito 1}
Calcolare una stima puntuale del parametro toracico medio $\mu$.

\textbf{Soluzione:}
Si vogliono stimare media e varianza, $\mu, \sigma^2$, del campione casuale $\{X_1,\dots,X_6\}$.
Usiamo lo stimatore classico della media campionaria:
$$ \bar{X}_n = \frac{1}{n} \sum_{i=1}^n X_i = \frac{84.60 + 79.10 + 84.50 + 79.60 + 74.00 + 74.10}{6} = \frac{475.9}{6} $$
La risposta è: 79.3166667

\subsubsection*{Quesito 2}
Calcolare una stima puntuale della deviazione standard, utilizzando il seguente stimatore della varianza campionaria:
$$ S^2(X_1,\dots,X_6) = \frac{1}{n} \sum_{i=1}^6 (X_i - \bar{X})^2 $$

\textbf{Soluzione:}
La deviazione standard è la radice quadrata della varianza campionaria fornita nella formula.
Svolgendo i calcoli degli scarti quadratici della singola misurazione dalla media $\bar{X}$ trovata e mediandoli per $n=6$, otteniamo la varianza. Per cui usiamo infine $\sqrt{S^2(X_1,\dots,X_6)}$ per ricondurci all'unità di misura originale.

La risposta è: 4.2892955

\subsubsection*{Quesito 3}
Lo stimatore $S^2(X_1,\dots,X_6)$ del punto precedente è non-distorto?
Se non-distorto, allora rispondere 1, altrimenti inserire il fattore (costante moltiplicativa) per cui si deve moltiplicare $S^2(X_1,\dots,X_6)$ per ottenere una stima migliore della varianza.

\textbf{Soluzione:}
Come avete visto a lezione, lo stimatore diviso per $n$ tende sistematicamente a sottostimare la reale ampiezza della variabilità di popolazione:
$$ \mathbb{E}\left(\frac{1}{n} \sum_{i=1}^n (Y_i - \bar{Y})^2\right) = \frac{n-1}{n} \sigma^2 $$
per cui no, lo stimatore $S^2(X_1,\dots,X_6)$ è distorto (sottostimando in media di un fattore proporzionale). Il fattore correttivo richiesto per recuperare lo stimatore corretto di Bessel risulta essere l'inverso del bias, cioè $\frac{n}{n-1}$.
Essendo $n=6$, il fattore è $\frac{6}{5} = 1.2$.

La risposta è: 1.2

\newpage

% ==========================================
%               ESERCIZIO 14
% ==========================================
\vspace{0.5cm}
\begin{tcolorbox}[colback=lightergray!10, colframe=tocsec, title=\textbf{Esercizio 14}]
Sia $\{X_1, X_2, X_3, X_4\}$ un campione casuale da distribuzione esponenziale con parametro $\lambda=0.87$.
Si considerino i seguenti stimatori:
$$ T_1 = X_3 \quad , \quad T_2 = \frac{X_1 + 2X_2}{3} \quad , \quad T_3 = \frac{X_1 + X_2 + X_3}{6} + \frac{X_3 + X_4}{4} $$
$X_i \sim E(\lambda)$, $f_{X_i}(x) = \lambda\exp\{-\lambda x\}$, $\mathbb{E}(X) = \frac{1}{\lambda}$.
\end{tcolorbox}

\subsubsection*{Quesito 1}
Quali tra i precedenti stimatori sono non-distorti?
\textit{Inserire il numero corrispondente allo stimatore, e.g. 1 per $T_1$, 7 per "tutti e tre" ecc.}

\textbf{Soluzione:}
Dobbiamo calcolare $\mathbb{E}(T_i)$ per $i=1,2,3$.
$$ \mathbb{E}(T_1) = \mathbb{E}(X_3) = \frac{1}{\lambda} $$
$$ \mathbb{E}(T_2) = \frac{1}{3} \mathbb{E}(X_1 + 2X_2) = \frac{1}{3} (\mathbb{E}(X_1) + 2\mathbb{E}(X_2)) = \frac{1}{3}\left(\frac{3}{\lambda}\right) = \frac{1}{\lambda} $$
$$ \mathbb{E}(T_3) = \frac{1}{6} \mathbb{E}(X_1 + X_2 + X_3) + \frac{1}{4} \mathbb{E}(X_3 + X_4) = \frac{3}{6} \frac{1}{\lambda} + \frac{2}{4} \frac{1}{\lambda} = \frac{1}{2}\frac{1}{\lambda} + \frac{1}{2}\frac{1}{\lambda} = \frac{1}{\lambda} $$
Ergo tutti gli stimatori mantengono intatta la speranza matematica del parametro, perciò tutti e tre sono non-distorti.

La risposta è: 7

\subsubsection*{Quesito 2}
Calcolare la varianza dello stimatore $T_1$.

\textbf{Soluzione:}
Calcoliamo la varianza di ognuno degli stimatori, ricordando che $Cov(X_i, X_j) = 0$ per $i \neq j$ (per indipendenza, cfr. definizione campione casuale). Sappiamo che la varianza per l'esponenziale è $Var(X) = \frac{1}{\lambda^2} = \lambda^{-2}$.
$$ Var(T_1) = Var(X_3) = \lambda^{-2} $$
Sostituendo $\lambda=0.87$, ricaviamo $\frac{1}{0.87^2} = \frac{1}{0.7569} \approx 1.3211785$.
(Aggiungiamo per completezza le altre per il quesito successivo):
$$ Var(T_2) = Var\left(\frac{X_1 + 2X_2}{3}\right) = \frac{1}{9} (Var(X_1) + 4Var(X_2)) = \frac{5}{9} \lambda^{-2} $$
$$ Var(T_3) = Var\left(\frac{2X_1 + 2X_2 + 5X_3 + 3X_4}{12}\right) = \frac{1}{144} Var(2X_1 + 2X_2 + 5X_3 + 3X_4) = \frac{42}{144} \lambda^{-2} $$

La risposta è: 1.3211785

\newpage

\subsubsection*{Quesito 3}
Quale degli stimatori è preferibile (in termini di errore quadratico medio)?

\textbf{Soluzione:}
Ricordiamo che l’errore quadratico medio di uno stimatore $T$ di $\theta$, $EQM(T) = \mathbb{E}[(T-\theta)^2]$, può essere decomposto nella somma della varianza dello stimatore più il quadrato della sua distorsione. Essendo gli stimatori non-distorti, ci basta guardare alla varianza degli stimatori e chiederci quale abbia minore varianza per decretare la migliore efficienza.
$$ \min\left\{1, \frac{5}{9}, \frac{42}{144}\right\} = \frac{42}{144} $$
dunque lo stimatore che minimizza la dispersione è $T_3$, che risulta essere lo stimatore preferibile.

La risposta è: 3

\subsubsection*{Quesito 4}
Dopo aver determinato lo stimatore di massima verosimiglianza $\hat{\lambda}$ per $\lambda$, inserirne la stima per il seguente campione casuale: 
\begin{center}
1.10, 1.80, 0.40, 0.60, 1.20, 1.40, 0.50, 8.10, 0.90, 0.10
\end{center}

\textbf{Soluzione:}
Scriviamo la funzione di verosimiglianza che, per l’indipendenza delle $X_i$, è data dal prodotto delle densità $f_{X_i}$:
$$ L(\lambda; x_1,\dots,x_n) = \prod_{i=1}^n \lambda \exp\{-\lambda x_i\} = \lambda^n \exp\left(-\lambda \sum_{j=1}^n x_j\right) $$
e la log-verosimiglianza:
$$ l(\lambda; x_1,\dots,x_n) = n \ln(\lambda) - \lambda \sum_{j=1}^n x_j. $$
Derivando rispetto al parametro $\lambda$ otteniamo la cosiddetta score function:
$$ S(\lambda) = \frac{d l(\lambda)}{d\lambda} = \frac{n}{\lambda} - \sum_{j=1}^n x_j $$
e la score equation $S(\lambda) = 0$ ci dà lo stimatore richiesto:
$$ \hat{\lambda}_n = \frac{n}{\sum_{j=1}^n x_j} = \frac{1}{\bar{X}_n}. $$
Sostituendo la somma dei 10 valori del campione (la cui somma è $16.1$) si ottiene la stima per $\lambda = \frac{10}{16.1}$.

La risposta è: 0.621118

\newpage

% ==========================================
%               ESERCIZIO 15
% ==========================================
\vspace{0.5cm}
\begin{tcolorbox}[colback=lightergray!10, colframe=tocsec, title=\textbf{Esercizio 15}]
Sia $X$ una v.a. esponenziale di parametro $\lambda$, i.e. $f_X(x) = \lambda e^{-\lambda x}$. Purtroppo non possiamo campionare direttamente $X$, però possiamo avere un campione casuale da $Y=8X$, $(Y_1,\dots,Y_n)$ ove $Y_i$ sono i.i.d. (indipendenti e identicamente distribuite).
\end{tcolorbox}

\subsubsection*{Quesito 1}
Dopo aver determinato la funzione di densità di $Y$ e lo stimatore di massima verosimiglianza per $\lambda$, inserirne il valore per il seguente campione:
\begin{center}
    0.94, 0.79, 0.30, 8.38, 0.15, 3.33, 3.88, 6.53, 0.75
\end{center}
\textbf{Soluzione:}
Sfruttiamo l'invarianza degli stimatori di Massima Verosimiglianza e la trasformazione della densità per un fattore di scala. La media campionaria è soggetta alla trasformazione $\bar{Y} = 8\bar{X}$, quindi $\bar{X} = \bar{Y}/8$. Essendo lo stimatore MLE dell'esponenziale il reciproco della media originaria:
$$ \hat{\lambda} = \frac{1}{\bar{X}} = \frac{8}{\bar{Y}} $$
Valutando numericamente dal campione fornito si ricava il valore approssimato.

La risposta è: 2.8742515

\subsubsection*{Quesito 2}
Si consideri ora $Z=X-10$, dopo aver determinato la funzione di densità di $Z$ e lo stimatore di massima verosimiglianza per $\lambda$ date $Z_1,\dots,Z_n$ i.i.d. dalla distribuzione di $Z$, inserirne il valore per il seguente campione:
\begin{center}
0.14, 0.04, 0.01, 0.00, 0.03, 0.11, 0.03, 0.11, 0.17
\end{center}

\textbf{Soluzione:}
Similmente, si applica uno shift lineare. Se $Z = X-10 \implies X = Z+10$. Di conseguenza la media campionaria originale è stimata da $\bar{X} = \bar{Z} + 10$.
Lo stimatore MLE di base rimane l'inverso del parametro centrale ricostruito:
$$ \hat{\lambda} = \frac{1}{\bar{Z} + 10} $$
Inserendo la media del nuovo campione calcolato, si deduce la stima empirica finale.

La risposta è: 0.0992939


\newpage
\section*{Probabilità e Statistica - Esercizi 9}
\addcontentsline{toc}{section}{Probabilità e Statistica - Esercizi 9}
\markboth{Esercizi 9}{Esercizi 9}

% ==========================================
%               ESERCIZIO 1
% ==========================================
\vspace{0.5cm}
\begin{tcolorbox}[colback=lightergray!10, colframe=tocsec, title=\textbf{Esercizio 1}]
Sia $(X_1,\dots,X_n)$ un campione casuale (i.e. le $X_i$ sono indipendenti ed identicamente distribuite) da una distribuzione binomiale di parametri $m=14$, numero di prove, e $\theta$, probabilità di successo.
Sia inoltre dato il seguente campione estratto:
\begin{center}
(13, 11, 9, 9, 7, 10, 13, 10, 10, 11, 9)
\end{center}
\end{tcolorbox}

\subsubsection*{Quesito 1}
Trovare lo stimatore di massima verosimiglianza di $\theta$. E riportare la stima di massima verosimiglianza.

\textbf{Soluzione:}
Per trovare lo stimatore di massima verosimiglianza, partiamo con lo scrivere la funzione di verosimiglianza sfruttando l'indipendenza del campione:
$$ \begin{aligned} L(x_1, \dots, x_n; \theta) &= P_{X_1, \dots, X_n}(x_1, \dots, x_n) \\ &= \prod_{i=1}^n P_{X_i}(x_i) \\ &= \prod_{i=1}^n \binom{m}{x_i} \theta^{x_i} (1-\theta)^{m-x_i} \\ &= \left[ \prod_{i=1}^n \binom{m}{x_i} \right] \theta^{\sum_i x_i} (1-\theta)^{nm - \sum_i x_i} \end{aligned} $$
ove nell’ultima uguaglianza abbiamo usato la commutatività del prodotto. Si osservi inoltre che $\left[ \prod_{i=1}^n \binom{m}{x_i} \right] = C$ è una costante e non dipende dal parametro $\theta$.

Prendiamo ora la log-verosimiglianza (il logaritmo naturale di $L$):
$$ \ell(x_1, \dots, x_n; \theta) = \log C + \sum_i x_i \log\theta + \left(nm - \sum_i x_i\right)\log(1-\theta) $$
e deriviamo rispetto a $\theta$, ponendo la derivata nota come \textit{score function} uguale a zero per trovare il punto di massimo:
$$ 0 = \frac{\partial}{\partial\theta}\ell(x_1, \dots, x_n; \theta) = \frac{\sum_i x_i}{\theta} - \frac{nm - \sum_i x_i}{1-\theta} $$
Risolvendo l'equazione rispetto a $\theta$:
$$ 0 = (1-\theta)\sum_i x_i - \theta\left(nm - \sum_i x_i\right) \implies \hat{\theta}_{ML} = \frac{\sum_i x_i}{nm} = \frac{\bar{X}_n}{m} $$
Sostituendo la somma dei dati estratti ($\sum x_i = 112$), $n=11$ e $m=14$, otteniamo il valore della stima.

La risposta è: 0.7272727

\subsubsection*{Quesito 2}
Calcolare l’intervallo di confidenza per $\theta$ (alla Wald) a livello di confidenza del 91\%, per il campione estratto.

\textbf{Soluzione:}
Sappiamo che per le singole v.a. $X_i \sim Bin(m,\theta)$, $\mathbb{E}(X_i) = m\theta$ e $Var(X_i) = m\theta(1-\theta)$.
Abbiamo determinato lo stimatore $\hat{\theta} = \frac{1}{nm} \sum_{i=1}^n X_i$, per le cui proprietà di correttezza sappiamo che:
$$ \mathbb{E}(\hat{\theta}) = \theta \quad \text{e} \quad Var(\hat{\theta}) = \frac{\theta(1-\theta)}{nm} $$
Inoltre, possiamo assumere la convergenza asintotica per il Teorema Centrale del Limite (infatti potremmo anche vedere le nostre $n$ variabili binomiali come $nm$ variabili Bernoulliane indipendenti, ed essendo $nm \gg 1$ l'approssimazione regge):
$$ Z_n = \frac{\hat{\theta} - \theta}{\sqrt{\frac{\theta(1-\theta)}{nm}}} \xrightarrow{d} N(0,1) $$
Ci chiediamo quindi per quali quantili si verifica:
$$ 1-\alpha \approx P\left( z_{\frac{\alpha}{2}} \leq \frac{\hat{\theta} - \theta}{\sqrt{\frac{\theta(1-\theta)}{nm}}} \leq z_{1-\frac{\alpha}{2}} \right) $$
Nell'approccio ``alla Wald'', si sostituisce il parametro $\theta$ incognito a denominatore con il suo stimatore puntuale $\hat{\theta}$:
$$ \begin{aligned} 1-\alpha &\approx P\left( z_{\frac{\alpha}{2}} \leq \frac{\hat{\theta} - \theta}{\sqrt{\frac{\hat{\theta}(1-\hat{\theta})}{nm}}} \leq z_{1-\frac{\alpha}{2}} \right) \\ &= P\left( z_{\frac{\alpha}{2}} \sqrt{\frac{\hat{\theta}(1-\hat{\theta})}{nm}} \leq \hat{\theta} - \theta \leq z_{1-\frac{\alpha}{2}} \sqrt{\frac{\hat{\theta}(1-\hat{\theta})}{nm}} \right) \\ &= P\left( -\hat{\theta} + z_{\frac{\alpha}{2}} \sqrt{\dots} \leq -\theta \leq -\hat{\theta} + z_{1-\frac{\alpha}{2}} \sqrt{\dots} \right) \\ &= P\left( \hat{\theta} - z_{1-\frac{\alpha}{2}} \sqrt{\frac{\hat{\theta}(1-\hat{\theta})}{nm}} \leq \theta \leq \hat{\theta} - z_{\frac{\alpha}{2}} \sqrt{\frac{\hat{\theta}(1-\hat{\theta})}{nm}} \right) \end{aligned} $$
Da cui, ricordando la simmetria della normale per cui $z_{\frac{\alpha}{2}} = -z_{1-\frac{\alpha}{2}}$, ricaviamo l'intervallo stimato:
$$ \hat{\theta} \pm z_{1-\frac{\alpha}{2}} \sqrt{\frac{\hat{\theta}(1-\hat{\theta})}{nm}} $$
(Suggerimento: usare \texttt{qnorm()} per trovare i quantili della normale corrispondenti al livello di confidenza richiesto).

La risposta è: (0.666427829248847, 0.788117625296607)

\newpage

% ==========================================
%               ESERCIZIO 2
% ==========================================
\vspace{0.5cm}
\begin{tcolorbox}[colback=lightergray!10, colframe=tocsec, title=\textbf{Esercizio 2}]
Sia $(X_1,\dots,X_n)$ un campione casuale (i.e. le $X_i$ sono indipendenti ed identicamente distribuite) da una distribuzione esponenziale di parametro $\theta$.
Sia inoltre dato il seguente campione estratto:
\begin{center}
(0.57, 0.13, 0.07, 0.02, 0.19, 1.66, 0.3, 0.35, 1.37, 1.08, 2.11, 1.77)
\end{center}
\end{tcolorbox}

\subsubsection*{Quesito 1}
Trovare lo stimatore di massima verosimiglianza di $\theta$ e calcolarne la stima per il campione estratto e riportare il valore della stima.

\textbf{Soluzione:}
$X_i \sim Exp(\theta)$ per $i=1,2,\dots,n$.
Scriviamo la funzione di verosimiglianza per il campione $(x_1, x_2, \dots, x_n)$:
$$ L(x_1, \dots, x_n; \theta) = \prod_{i=1}^n f_{X_i}(x_i; \theta) = \prod_{i=1}^n \theta e^{-\theta x_i} = \theta^n e^{-\theta \sum_{i=1}^n x_i} $$
Prendiamo la log-verosimiglianza:
$$ \ell(x_1, \dots, x_n; \theta) = n\log\theta - \theta\sum_{i=1}^n x_i $$
Infine deriviamo e poniamo la derivata uguale a zero per massimizzare la funzione:
$$ \frac{\partial \ell}{\partial \theta} = \frac{n}{\theta} - \sum_{i=1}^n x_i = 0 \implies \hat{\theta}_{ML} = \frac{n}{\sum_{i=1}^n x_i} $$
Ossia, lo stimatore di massima verosimiglianza per il parametro dell’esponenziale è dato dal reciproco della media campionaria $\hat{\theta}_{ML} = (\bar{X}_n)^{-1}$.

La risposta corretta è: sol1 (n.d.r. testo originale)

\subsubsection*{Quesito 2}
Stimare la media (valore atteso) della popolazione, utilizzando il campione estratto e lo stimatore del quesito 1.

\textbf{Soluzione:}
Sappiamo che il valore atteso di una v.a. esponenziale di parametro $\theta$ è $\mathbb{E}(X) = \frac{1}{\theta}$. Per il principio di invarianza degli stimatori ML, ci basta prendere il reciproco della stima precedente, ossia calcolare direttamente la media campionaria:
$$ \widehat{\mathbb{E}(X)} = \frac{1}{\hat{\theta}_{ML}} = \bar{X}_n $$
Calcolando la media dei 12 valori osservati si ottiene il risultato.

La risposta corretta è: 0.8016667

\subsubsection*{Quesito 3}
Usando il teorema centrale del limite, calcolare l’intervallo di confidenza (approssimato) per $\theta$ a livello di confidenza del 91\%, dato il campione estratto.

\textbf{Soluzione:}
Sia $X_i \sim Exp(\theta)$, con $\mathbb{E}(X_i) = \frac{1}{\theta}$ e $Var(X_i) = \frac{1}{\theta^2}$.
Usando il Teorema Centrale del Limite abbiamo che:
$$ Z_n = \frac{\bar{X}_n - \frac{1}{\theta}}{\frac{1}{\theta \sqrt{n}}} \xrightarrow{d} N(0,1) $$
in distribuzione. Possiamo quindi utilizzare la distribuzione normale standard per stimare gli estremi dell’intervallo di confidenza. In particolare:
$$ 1-\alpha \approx P\left( z_{\frac{\alpha}{2}} \leq \frac{\bar{X}_n - \frac{1}{\theta}}{\frac{1}{\theta \sqrt{n}}} \leq z_{1-\frac{\alpha}{2}} \right) $$
Manipoliamo l'espressione per isolare il parametro teorico $\theta$ al centro (usando $\bar{X}_n = \frac{1}{\hat{\theta}}$):
$$ \begin{aligned} 1-\alpha &\approx P\left( z_{\frac{\alpha}{2}} \frac{1}{\sqrt{n}} \leq \frac{\frac{1}{\hat{\theta}} - \frac{1}{\theta}}{\frac{1}{\theta}} \leq z_{1-\frac{\alpha}{2}} \frac{1}{\sqrt{n}} \right) \\ &= P\left( z_{\frac{\alpha}{2}} \frac{1}{\sqrt{n}} \leq \frac{\theta}{\hat{\theta}} - 1 \leq z_{1-\frac{\alpha}{2}} \frac{1}{\sqrt{n}} \right) \\ &= P\left( 1 + z_{\frac{\alpha}{2}} \frac{1}{\sqrt{n}} \leq \frac{\theta}{\hat{\theta}} \leq 1 + z_{1-\frac{\alpha}{2}} \frac{1}{\sqrt{n}} \right) \\ &= P\left( \hat{\theta}\left(1 + \frac{z_{\alpha/2}}{\sqrt{n}}\right) \leq \theta \leq \hat{\theta}\left(1 + \frac{z_{1-\alpha/2}}{\sqrt{n}}\right) \right) \end{aligned} $$
Da cui, ricordando che $z_{\frac{\alpha}{2}} = -z_{1-\frac{\alpha}{2}}$, ricaviamo analiticamente la formula per gli estremi dell'intervallo:
$$ \hat{\theta} \left(1 \pm \frac{z_{1-\frac{\alpha}{2}}}{\sqrt{n}}\right) $$

La risposta è: (0.636899173965222, 1.85790332083727)


\end{document}
